% concatenated from main.tex, section_introduction.tex, section_model.tex, section_solution.tex, section_mfg.tex, section_experiments.tex, section_conclusion.tex, section_solution_proofs.tex, section_mfg_proofs.tex
\documentclass[12pt]{article}

\usepackage[T1]{fontenc}
\usepackage{amsmath,amssymb,amsthm}
\usepackage{mathrsfs}
\usepackage{bm}
\usepackage[colorlinks,citecolor=blue,urlcolor=blue,linkcolor=blue]{hyperref}
\usepackage[numbers]{natbib}
\usepackage{graphicx}
\usepackage{booktabs}
\usepackage{float}      % if you want [H]
\usepackage{caption}
\usepackage{placeins}
\usepackage{authblk}

\makeatletter
\setlength{\@fptop}{0pt}
\makeatother

\setlength{\textwidth}{6in}
\setlength{\textheight}{8.5in}
\setlength{\topmargin}{-0.25in}
\setlength{\evensidemargin}{.25in}
\setlength{\oddsidemargin}{.25in}
\setlength{\parskip}{.65ex}

\numberwithin{equation}{section}

\newtheorem{theorem}{Theorem}[section]
\newtheorem{lemma}[theorem]{Lemma}
\newtheorem{proposition}[theorem]{Proposition}
\newtheorem{corollary}[theorem]{Corollary}
\newtheorem{remark}[theorem]{Remark}

\theoremstyle{definition}
\newtheorem{definition}[theorem]{Definition}
\newtheorem{assumption}[theorem]{Assumption}

\usepackage{algorithm}
\usepackage{algpseudocode}
\usepackage{xcolor}

\title{Strategic Index Reconstitution: Differential Games, Closed-Loop Equilibria and Mean-Field Dynamics}
\author[1]{Lukas-Benedikt Fiechtner}
\author[1,2]{Jose Blanchet}
\affil[1]{Institute for Computational and Mathematical Engineering, Stanford University}
\affil[2]{Department of Management Science and Engineering, Stanford University}
\date{\today}

\begin{document}

\vspace{-0.5cm}
\maketitle

\begin{abstract}
We study strategic trading around index reconstitution in a continuous-time,
multiasset game with transient cross-asset price impact and heterogeneous
beliefs about future index membership. Opportunistic traders position before
a public announcement, adjust to the revealed composition, and trade around
an indexer following a prescribed execution schedule. Under a
no-price-manipulation condition, we construct a subgame-perfect
Nash equilibrium on every finite horizon. The affine feedback policies are
computed from a non-standard matrix Riccati equation and linear ordinary differential
equations whose number and dimensions do not grow with the number of traders.
Aggregate inventories and price impact depend on beliefs only through the
population-average belief, while differences in beliefs affect individual
inventory positions. Under mean-field scaling, we construct a
mean-field equilibrium and obtain quantitative convergence and
approximate-Nash bounds. Numerical illustrations show how competition and impact
decay determine the balance between adverse price displacement from
anticipatory trading and savings in the indexer's execution costs when
opportunists trade against its orders during implementation.
\end{abstract}

\begin{description}
\item[Mathematics Subject Classification (2020):] 91A15, 91A16, 91G15, 93E20
\item[JEL Classification:] C73, C61, G11, G14
\item[Keywords:] index rebalancing, stochastic differential games, closed-loop Nash equilibrium, mean-field games, transient market impact, heterogeneous beliefs
\end{description}

\bigskip

\section{Introduction}\label{sec:introduction}

When an index changes its constituents, index funds must adjust their holdings
to track the new portfolio. The changes become public on the announcement date and index funds typically execute the required trades in the closing auction on
the later implementation date, immediately before the new index composition
takes effect. To profit from these required trades, opportunists may build
anticipatory positions before the announcement, revise them once the changes
are announced, and unwind them as index funds rebalance. Their problem is
dynamic: each trade moves prices with impact persisting for some time, and other
opportunists can respond to the resulting changes in prices and inventories.
Press accounts illustrate the financial stakes of anticipating index trades:
The Economist has called index rebalancing ``the biggest event in markets.''\footnote{The Economist,
\href{https://www.economist.com/finance-and-economics/2026/05/12/index-rebalancing-is-now-the-biggest-event-in-markets}
{\emph{Index Rebalancing Is Now the Biggest Event in Markets}}, May 12, 2026.}
Bloomberg reported that two Millennium
index-rebalancing teams
earned about \$3.7 billion in June 2026 and that the same teams had lost about
\$900 million in early 2025.\footnote{Nishant Kumar,
\href{https://www.bloomberg.com/news/articles/2026-07-06/two-millennium-trading-pods-made-about-3-7-billion-last-month}
{\emph{Two Millennium Trading Pods Made About \$3.7 Billion Last Month}},
Bloomberg, July 6, 2026. The report covers several June rebalancing events,
including quarter-end multi-asset rebalancing. Bei Hu, Katherine Burton, and
Nishant Kumar,
\href{https://www.bloomberg.com/news/articles/2025-03-08/millennium-loses-900-million-on-strategy-roiled-by-market-chaos}
{\emph{Millennium Loses \$900 Million on Strategy Roiled by Market Chaos}},
Bloomberg, March 7, 2025.}

The empirical record reinforces the importance of understanding how
opportunists trade around index rebalancing.
Studies of S\&P 500 additions in the late 1970s and early 1980s document
significant abnormal returns of roughly \(3\%\) on the first trading day
after announcement
\citep{Shleifer1986,HarrisGurel1986}, with subsequent reversals documented by
\citet{HarrisGurel1986}. However, average abnormal returns around announcement and
implementation have declined in recent years: comparing the 1990s with 2010--2020,
\citet{GreenwoodSammon2025} report a fall from \(4.94\%\) to \(1.02\%\)
around announcement and from \(3.55\%\) to \(0.26\%\) around
implementation.\footnote{For these event windows, abnormal returns are measured
net of the S\&P 500 return, from the preceding trading day's close to the
following trading day's close.}
On the other hand, abnormal returns over the 100 trading days \emph{before} announcement
increased from \(9.60\%\) to \(18.68\%\) over the same periods
\citep[Table~IV]{GreenwoodSammon2025}, suggesting increased anticipatory trading
even before the announcement.
\footnote{The increase may not only reflect trading in anticipation of
index inclusion, but also selection of past winners, or both \citep{GreenwoodSammon2025}.}

Despite the decline in abnormal returns around announcement and implementation,
S\&P 500 reconstitution volume relative to normal daily trading has risen
\citep{ChincoSammon2024}. For additions from outside the S\&P index family in
2010--2024, \citet{PegoraroSammonShim2026} report implementation-day turnover
near \(20\%\) of shares outstanding despite an average abnormal return of
\(-0.09\%\). Interestingly, implementation-day volume substantially exceeds the trades
implied by labeled index funds alone, consistent with additional trading by
institutions managing their own index portfolios and active managers closely
tracking benchmarks \citep{ChincoSammon2024}. Using equal-weight rebalances,
\citet{Nathan2026} finds substantial additional displayed buying and selling
interest in the closing auction.

Furthermore, for S\&P 500 additions from outside the index family,
\citet{PegoraroSammonShim2026} find short interest rising around implementation
and declining afterward, consistent with opportunists going short to supply
inventory to index funds in addition to selling the inventory they built up
before implementation.
\citet{PegoraroSammonShim2026} also relate the timing of inventory buildup to
predictability: less predictable Russell migrations exhibit more trading
between the rank and implementation dates, which they interpret as consistent
with earlier positioning when membership changes are easier to anticipate.

This empirical evidence motivates us to ask who holds the required inventory
before index funds arrive, how that inventory is transferred, and how
competition changes the cost of the transfer. We study how opportunists build
positions before announcement using their probabilities over possible index
reconstitutions, revise their positions when membership is revealed, and
transfer inventory as index funds rebalance.

\paragraph{Model.}
We briefly outline the model here and give the full specification later in
Section~\ref{sec:model}.
We consider \(n\geq2\) opportunists trading \(d\) assets over a finite
horizon \([0,T]\), with announcement at \(T_{\mathrm{ann}}\) and
implementation at \(T_{\mathrm{impl}}\). Before \(T_{\mathrm{ann}}\),
speculators consider \(m\) possible scenarios for the new index constitution.
Each scenario \(k\) determines the indexer's required trades, including
purchases of added assets and sales of deleted assets. The corresponding
trading rate \(g^k\) approximates execution in the closing auction by
spreading these purchases and sales uniformly over
\([T_{\mathrm{TWAP}},T_{\mathrm{impl}}]\), with
\(T_{\mathrm{ann}}<T_{\mathrm{TWAP}}<T_{\mathrm{impl}}\).
The realized scenario is announced publicly at \(T_{\mathrm{ann}}\) and determines the realized
trading rate \(g\) followed by the indexer.
The remaining interval from \(T_{\mathrm{impl}}\) to \(T\) allows opportunists to unwind
positions after implementation.
We allow traders to hold different beliefs about the possible index
constitutions: before the announcement, trader \(i\)'s subjective probability
measure \(\mathbb P_n^i\) assigns probabilities \(\pi_n^i\in\Delta_m\)
to the possible scenarios.
Each trader knows its own belief vector (which we also call prior) and the sample-average belief
\(\bar\pi_n=\frac1n\sum_{j=1}^n\pi_n^j\),
but need not know the other traders' beliefs. We assume that no further
information about index membership arrives before \(T_{\mathrm{ann}}\), so
beliefs remain fixed until the public announcement.

Trader \(i\)'s inventory \(X_n^i\) evolves as \(dX^i(t)=u^i(t)dt\) with trading rate \(u_n^i\).
The midprice is \(S=S^0+I\), where \(S^0\) is the fundamental price which is a square-integrable c\`adl\`ag
martingale independent of the realized index scenario under each \(\mathbb P_n^i\) and
\(I\) is the transient impact satisfisfying
\begin{equation}\label{eq:intro_dynamics}
dI(t)=\left[-RI(t)+\Gamma\left(\varpi_n\sum_{j=1}^n u_n^j(t)+g(t)\right)\right]dt.
\end{equation}
Here \(R\) governs the decay 
of existing impact, \(\Gamma\) maps new trades into direct and
cross-asset impact, and \(\varpi_n>0\) scales each trader's contribution to the impact.
Under the fixed-market-size scaling \(\varpi_n=1\), increasing \(n\)
adds opportunists of fixed size to the same market. Under the mean-field
scaling \(\varpi_n=\varpi/n\), total opportunist mass \(\varpi>0\)
stays fixed as \(n\) increases and individual traders become smaller.
Trader \(i\)'s execution price is
\[
\widehat S^i(t)=S(t)+\Lambda\left(u_n^i(t)
+\chi_u\varpi_n\sum_{j\ne i}u_n^j(t)+\chi_g g(t)\right),
\]
where \(\Lambda\succ0\) and \(\chi_u,\chi_g\geq0\).
Trader \(i\) minimizes expected trading costs, deducting the change in
inventory value at fundamental prices and adding quadratic running and
terminal inventory penalties:
\[
\begin{aligned}
\mathbb E_n^i\Bigg[
&\int_0^T \big(\widehat S^i(t)\big)^\top u_n^i(t)\,dt
-\left((S^0(T))^\top X_n^i(T)-(S^0(0))^\top X_n^i(0)\right)\\
&+\frac{1}{2}\int_0^T (X_n^i(t))^\top Q X_n^i(t)\,dt
+\frac{1}{2}(X_n^i(T))^\top Q_T X_n^i(T)
\Bigg],
\end{aligned}
\]
where \(\mathbb E_n^i\) denotes expectation under \(\mathbb P_n^i\)
and \(Q,Q_T\succeq0\) penalize running and terminal inventory, respectively.

Traders choose their admissible trading rates using \emph{closed-loop} strategies,
that is, feedback functions of the
histories of its own inventory, average inventory, the midprice, and the
public announcement up to the current time (and also of the trader's own belief, the population-average belief). Under a unilateral deviation,
trader \(i\) changes its feedback function while the other traders keep
theirs fixed. The deviation can change the realized inventory and price
histories entering the other traders' feedback functions. Their
\emph{realized trading rates} can therefore change despite their fixed
feedback functions. In the \emph{open-loop} equilibria commonly studied in
execution games, the other traders' trading-rate processes instead remain
fixed under a deviation by trader \(i\)
\citep{strehle_2017,NeumanVoss2023}.
We seek a \emph{subgame-perfect} Nash equilibrium: a closed-loop Nash equilibrium
in which, from every starting time and admissible state, including off the
equilibrium path, no trader can lower its remaining expected cost by
unilaterally switching to another admissible feedback function.

\paragraph{Main Results and Contributions.}
\begin{enumerate}
\item \textbf{Construction and Global Existence of a Feedback Equilibrium.}
In Theorem~\ref{thm:stitched_hjb_verification}, we prove a verification
result: sufficiently smooth solutions of a system of
coupled Hamilton--Jacobi--Bellman (HJB) equations satisfying
terminal and announcement conditions yield a subgame-perfect Nash
equilibrium through their admissible minimizing feedbacks that respect the information 
structure restrictions we impose.
The result covers all admissible history-dependent unilateral deviations.

We construct solutions of the HJB equations that are quadratic in the
reduced state
\(
z_n^i=\bigl(X_n^i-\bar X_n,\,\bar X_n,\,I\bigr)\in\mathbb R^{3d}
\),
where \(\bar X_n=\frac1n\sum_{j=1}^n X_n^j\) is average inventory.
The inventory deviation \(X_n^i-\bar X_n\) is private to trader \(i\),
while \(\bar X_n\) and \(I\) are common to all traders.
For a generic reduced state \(z^i\), we use the quadratic ansatz
\begin{equation}\label{eq:intro_quadratic_value_ansatz}
\begin{aligned}
V_n^{i,0}(t,z^i)
&=
\frac{1}{2}(z^i)^\top H_n(t)z^i
+(r_n^{i,0}(t))^\top z^i+c_n^{i,0}(t),\\
V_n^{i,k}(t,z^i)
&=
\frac{1}{2}(z^i)^\top H_n(t)z^i
+(r_n^k(t))^\top z^i+c_n^k(t),
\qquad k=1,\ldots,m,
\end{aligned}
\end{equation}
before announcement and in scenario \(k\) afterward, respectively.
Substitution into the HJB equations yields a nonstandard
\(3d\times3d\) matrix Riccati ordinary differential equation (ODE) for
the common quadratic coefficient \(H_n\). The remaining coefficients are
obtained by solving linear ODEs sequentially. For fixed \(d\) and \(m\),
only a fixed number of \(3d\)-dimensional linear ODEs is needed to determine
the feedback. The complexity of computing the Nash equilibrium, measured
by the number and dimensions of the ODEs, therefore does not increase with
\(n\). Theorem~\ref{thm:full_horizon_closed_loop_equilibrium} shows that,
whenever the coefficient equations have a solution on the full horizon,
the quadratic value functions and associated affine feedbacks satisfy
the verification theorem.

Standard ODE theory guarantees only local solvability of the Riccati
equation. The cost contains indefinite cross terms between trading rates
and outstanding impact, so there is no immediate bound ruling out
finite-time blow-up before the backward solution reaches time zero.
Yet, under the economically reasonable conditions ruling out price manipulation
stated in Assumption~\ref{ass:standing_model},
Theorem~\ref{thm:global_riccati_solvability} establishes upper and lower bounds
\begin{equation}\label{eq:intro_riccati_bounds}
-\underline H(t)\preceq H_n(t)\preceq\overline H(t),
\qquad 0\leq t\leq T,
\end{equation}
uniformly over \(n\geq2\). The bounds prevent finite-time blow-up and
yield a unique Riccati solution on every finite horizon and hence a
subgame-perfect Nash equilibrium. No short-horizon restriction or additional
smallness condition on transient impact is required.
To our knowledge, this is the first globally solvable continuous-time,
finite-player, closed-loop execution game with transient price impact.
Most related finite-player execution games instead study open-loop equilibria
\citep{strehle_2017,schied_zhang_2017,schied_strehle_zhang_2017,
schied_zhang_2019,luo_schied_2019,NeumanVoss2023,
campbell2026optimalexecutionntraders} for which existence of a Nash equilibrium is
usually easier to prove.

\item \textbf{Beliefs and Public Revelation.}
Our solution shows how the sample-average belief and individual belief
deviations affect inventories, price impact, and trading rates.
In Remark~\ref{rem:full_procedure_decomposition}, we decompose the equilibrium
feedbacks and state dynamics into common and individual components.
Before announcement, the common state \((\bar X_n,I)\) and the inventory
deviations \(X_n^i-\bar X_n\) satisfy decoupled linear systems, with forcing
terms linear in the sample-average belief \(\bar\pi_n\) and player \(i\)'s belief deviation
\(\pi_n^i-\bar\pi_n\), respectively.
After announcement, the realized indexer schedule affects only the common
system, while inventory deviations evolve identically across scenarios
from their announcement-time values, preserving the effect of pre-announcement
disagreement.

\item \textbf{Mean-Field Approximation.}
The finite-player feedbacks use current average inventory, which may need
not always be observable in practice. Additionally, in practice the number of players
may be high, but not known exactly. For these reasons, we
construct mean-field policies using aggregate paths determined by the
mean-field equilibrium, without requiring real-time observation of
aggregate inventory or the exact number of competitors.

Under the mean-field scaling, we define a game whose player types are
belief vectors \(\pi\in\Delta_m\) over membership scenarios. The mean field
averages trading across a distribution of belief types, which may be
continuous. Under the conditions ruling out price
manipulation in Assumption~\ref{ass:standing_model}, we construct an explicit
equilibrium and prove that the finite-dimensional linear system determining average
inventory and impact is nonsingular, yielding uniqueness within the affine
construction. A trader implements its policy using its own belief and
inventory, the population-average belief, and equilibrium paths for average
inventory and impact for each announcement scenario
(Theorems~\ref{thm:mf_solution} and \ref{thm:mf_implementability}).

Let \(\bar\pi\) denote the average belief in the limiting population and set
\begin{equation}\label{eq:intro_epsilon_n}
\varepsilon_n=\frac1n+\|\bar\pi_n-\bar\pi\|.
\end{equation}
For zero initial inventories and impact and \(\bar\pi_n\to\bar\pi\),
the constructed finite-player Nash equilibria converge to the mean-field
equilibrium as \(n\to\infty\). The errors in trading rates, inventories,
and impact are \(O(\varepsilon_n)\), uniformly over time, scenario, and
trader, comparing each trader with a mean-field trader with the same belief vector
(Theorem~\ref{thm:finite_to_mf_convergence}).

When all traders use the mean-field policies in the finite game, the
resulting profile is an approximate Nash equilibrium from the prescribed
initial state of zero initial inventory and impact: no trader can lower its expected cost by more than
\(C\varepsilon_n^2\) through an admissible
unilateral deviation, with \(C\) independent of \(n\) and the finite
collection of beliefs (Theorem~\ref{thm:mf_approximate_nash}).
For i.i.d.\ belief types, the path and Nash errors are \(O_p(n^{-1/2})\)
and \(O_p(n^{-1})\), respectively. When a deterministic population matches
the limiting average belief at rate \(1/n\), the errors are \(O(n^{-1})\)
and \(O(n^{-2})\), respectively.

\item \textbf{Numerical Illustrations.}
We use a stylized five-asset index rebalance to investigate qualitative
properties of the equilibrium. In the atomic player setting \(\varpi=1\)
we obatin the following insights:

\begin{itemize}
\item \textbf{Aggregate Inventory:}
Before announcement, opportunists build aggregate inventory according to
the average belief: long in assets the indexer is expected to buy and short
in those it is expected to sell. After announcement, aggregate positions grow in
correctly anticipated additions and deletions and reverse where the
revealed indexer order has the opposite sign.
During implementation, opportunists trade against the indexer, taking
aggregate inventory in added and deleted assets through zero before
unwinding afterward. In added stocks, they thus sell more than their
accumulated inventory and subsequently cover the resulting aggregate short
positions, consistent with the rise and subsequent decline in short
interest documented by \citet{PegoraroSammonShim2026}.
Adding opportunists of fixed size generally increases aggregate inventory
magnitudes at announcement and the start of implementation, with the largest
increases at small \(n\).

\item \textbf{Belief Dispersion and Trading Volume:}
While aggregate inventory follows expected indexer trading under the
population-average belief, individual positions also reflect a trader's
expectations relative to those of other opportunists. A trader can therefore
enter the announcement short in a stock it expects to be added if its
expected indexer purchases are sufficiently below the population average.
With the average belief fixed,
greater dispersion of traders' beliefs increases gross trading through
offsetting purchases and sales among opportunists, leaving net aggregate
inventory and the impact trajectory unchanged. Additionally, at a given belief dispersion, the additional
gross volume per opportunist is larger with more players.

\item \textbf{Price Impact:}
Price impact follows the direction of expected indexer trading before
announcement and adjusts to the realized order afterward.
Adding opportunists of fixed size generally raises impact magnitudes at
announcement but reduces them at implementation.

\item \textbf{Indexer Costs and Opportunist Wealth:}
Adding opportunists of fixed size lowers the indexer's implementation
shortfall with diminishing marginal savings. Both average and total
terminal trading wealth of opportunists decline.

\item \textbf{Parameter Sensitivity:}
We investigate how equilibrium inventories, price impact, indexer costs,
and opportunist wealth respond to changes in trading windows, execution
costs, impact strength, and resilience, the rate of impact decay.
Among these parameters,
only resilience produces a non-monotonic response in expected indexer cost
savings relative to the no-opportunist benchmark at the same parameter
values: savings peak at intermediate resilience.
With slow decay, adverse impact from opportunists' earlier
trades persists into implementation and can outweigh the benefits of their
opposing trades. Intermediate resilience allows much of this earlier impact
to decay while opposing trades still substantially reduce impact costs
during implementation. With fast decay, the indexer's own impact dissipates
quickly, reducing the benefit of offsetting its trades.
\end{itemize}
Finally, under the mean-field scaling \(\varpi_n=\varpi/n\),
the illustrations also show convergence to the mean-field inventory and
impact paths.
\end{enumerate}

\paragraph{Related Literature.}
The closest economic model is
\citet{PegoraroSammonShim2026}.  They study a five-period, two-stock model with
three trading dates (with an extension to arbitrarily many trading
dates). Index membership is public before trading begins (i.e. there is no pre-announcement trading), and an
indexer chooses its schedule against multiple identical Cournot
speculators, either without commitment or as a Stackelberg leader.  Purchases
of the incoming stock and sales of the outgoing stock use the same schedule,
and impact is asset-separable, with permanent and instantaneous components.
Their model answers a question that ours leaves exogenous: when should the
index fund trade? We fix that schedule without optimizing the indexer's tradeoff
between tracking error and execution cost, and study how several traders respond
before and after a public announcement when beliefs differ, impact propagates
across assets, and impact decays over time.  \citet{Tasitsiomi2025} studies a
complementary one-stock model with one manager and one anticipatory trader
after an index addition has been announced. Both papers study inventory
acquisition \emph{after} announcement and liquidity provision at implementation
using \emph{open-loop} trading schedules. Our \emph{closed-loop} game
with transient price impact includes inventory buildup \emph{before} announcement
under differing beliefs about index membership and position adjustments
after membership is revealed. We construct equilibrium policies that react
to the observed state and prove equilibrium existence on every finite
horizon.

The optimal-execution literature provides the foundations for our
transient-impact dynamics and no-manipulation assumption.
The resilient limit-order-book model of \citet{OBIZHAEVA20131} provides the
basic dynamics for the impact state. The no-manipulation restrictions for one- and multiasset
transient impact developed by
\citet{gatheral_2012,alfonsi2016,Schneider02012019,jaber2026} motivate the
nonnegative self-impact-cost condition in Assumption~\ref{ass:standing_model},
which underpins the global Riccati bounds discussed above.

Market-impact games model strategic interaction among traders whose orders
affect both their own and their rivals' execution prices.
In discrete time, \citet[Chapter~9]{schoeneborn2008trade} studies a
two-player transient-impact game in both open-loop and closed-loop
formulations. The open-loop two-player model of \citet{schied_zhang_2019}
is extended to arbitrary numbers of risk-averse traders by
\citet{luo_schied_2019} and to multiple assets by \citet{CordoniLillo2024}.
\citet{schied_strehle_zhang_2017} study the high-frequency limit of the
Schied--Zhang game with exponentially decaying impact. \citet{NutzProsperi2025}
extend the high-frequency analysis to \(n\) traders and recover a
continuous-time open-loop equilibrium with costs on the initial and terminal
block trades.

In continuous-time open-loop market-impact games, rivals' trading processes
remain fixed under a unilateral deviation. For transient impact, \citet{strehle_2017}
proves existence and uniqueness for any finite number of traders under a
general positive-type impact kernel, a positive quadratic trading cost, and
hard liquidation.
\citet{campbell2026optimalexecutionntraders} give explicit equilibria for
\(N\) traders in a one-asset Obizhaeva--Wang model, both with finite terminal
penalties and under hard liquidation.
In a companion paper, \citet{CampbellNutzRandomization} show that private randomization cannot
restore existence and prove that there is at most one equilibrium in a broad
class of transient-impact games.
\citet{CarteaJaimungalSanchezBetancourt2026} characterize a one-asset game
between a broker and an informed trader by coupled forward--backward stochastic
differential equations, with a horizon bound in their general solvability
result for genuinely resilient impact. The July 2026 revision of
\citet{AbiJaberNeumanVoss2023} treats a broad class of linear-quadratic
games over trading paths, including a transient-impact liquidation problem,
through open-loop variational methods. With permanent impact,
\citet{schied_zhang_2017} and \citet{DrapeauLuoSchiedXiong2021} obtain
finite-horizon open-loop equilibria under terminal constraints or monotonicity
conditions.

The literature on closed-loop market-impact games is much smaller than its
open-loop counterpart. Here, rivals retain their feedback
functions under a deviation, so their realized trading rates can respond to
the changed state. In continuous time, \citet{carlin2007} derive a feedback
HJB/Riccati system with permanent impact and instantaneous execution costs.
\citet{CarmonaYang2008PredatoryTrading} formulate a two-player feedback
game with permanent impact and instantaneous execution costs and analyze its
degenerate HJB system numerically. \citet{micheli2023} give a
continuous-time, finite-player closed-loop equilibrium with a common signal,
but their one-asset infinite-horizon model has no transient-impact state and
requires price impact to be sufficiently small. With transient price impact,
\citet{Ohnishi02102020,Ohnishi03072024} obtain Markov-perfect equilibria with
two players in discrete time. The missing case is a continuous-time,
finite-player transient-impact game with
subgame-perfect strategies and full-horizon solvability.  No prior result
supplies this conjunction even for one asset. Our theorem allows matrix-valued
cross-impact across \(d\) assets and differing beliefs.

In closed-loop trading models with prices determined by market clearing rather
than by prescribed transient-impact dynamics,
\citet{ChoiLarsenSeppi2019,ChoiLarsenSeppi2021Benchmarks}
study private trading targets and intraday benchmarks.
\citet{ChenChoiLarsenSeppi2022} introduce learning about latent demand and
construct a finite-horizon subgame-perfect market equilibrium.
\citet{ChenChoiLarsenSeppi2023} prove global existence for arbitrary finite
horizons in a one-stock consumption and asset-pricing model with endogenous
inventory-dependent price impact.

Mean-field convergence results connect finite-player market-impact games
to limiting games with a continuum of traders.
\citet{NeumanVoss2023} construct finite-player and
mean-field open-loop equilibria under transient impact and a common predictive
signal, and prove convergence and approximate-Nash bounds.  The
linear--quadratic game over trading paths of
\citet{AbiJaberNeumanVoss2023} also yields open-loop finite-player convergence
and approximate equilibria. Our analysis instead establishes convergence of
explicitly constructed finite-player \emph{closed-loop} equilibria and approximate-Nash
bounds for the mean-field policies. Our estimates also separate the finite-player
term \(1/n\) from the discrepancy between the finite-player and limiting
average beliefs. These model-specific results complement general closed-loop
mean-field convergence theory \citep{Lacker2020ClosedLoop,LackerLeFlem2023}
and quantitative convergence results for games with interaction through
controls \citep{LauriereTangpi2022}.

Mean-field models of trading and portfolio choice allow heterogeneity in
beliefs, preferences, and trading horizons. \citet{Casgrain2020} allow traders to have different
objectives and subjective probability measures.
In portfolio games without price impact, \citet{Lacker2019,Lacker2020}
allow heterogeneous risk tolerances and relative-performance concerns.
\citet{LeclereRosenbaum2026} study a continuum
of heterogeneous horizon types. Our types are instead beliefs about a common event
that is publicly revealed.

Our paper also connects to the literature on major--minor mean-field games.
Under our mean-field scaling, the indexer plays the role of a major agent
because its exogenously prescribed trading flow remains unscaled, whereas
the opportunists are minor players whose individual market impact vanishes
as the population grows.
\citet{Huang04032019} study anticipatory trading in a major--minor game in
which both the institutional trader and HFTs trade continuously.
\citet{ChengWangXu2024} introduce a trading-speed asymmetry: the large trader
trades at discrete dates and HFTs trade continuously, with both prescribed
and optimized large-trader schedules considered.
\citet{chen2024periodictradingactivitiesfinancial} study an open-loop
major--minor liquidation game in which the major trader balances execution
costs against inventory deviations from a prescribed periodic trading
benchmark.
These papers construct mean-field equilibria and establish finite-population
approximate-Nash properties. We explicitly solve the finite-player feedback
game, define and solve the corresponding mean-field game, and then prove
that the constructed finite-player equilibria converge to the mean-field
equilibrium.

Finally, the roles of predation and liquidity provision around a predictable
order have a long history.  Strategic traders may amplify a forced order or
supply inventory against it
\citep{brunnermeier_2005,carlin2007,schoeneborn2009liquidation,strehle_2017}.
\citet{BessembinderCarrionTuttleVenkataraman2016} show, theoretically and
empirically, how strategic trading can improve execution around predictable
ETF roll orders when impact is sufficiently resilient.
\citet{Herzing2026} studies a two-player continuous-time open-loop game with
uncertainty about opponents' required trading volumes and no belief updating.
Further evidence on execution costs,
predictability, and crowding around index events appears in
\citep{Petajisto2011,Li2021,MicheliNeuman2022,JiangWuYao2025}.  Our
contribution is to place these mechanisms inside one finite-player feedback
game that spans uncertainty, public revelation, implementation, and subsequent
unwinding, and to prove that the game is solvable on the entire horizon.

The remainder of the paper is organized as follows.
Section~\ref{sec:model} defines the scenario-dependent indexer flow, the
traders' objectives, and the admissible feedback policies.
Section~\ref{sec:solution} constructs the finite-player equilibrium and proves
global solvability.
Section~\ref{sec:mean_field_limit} develops the mean-field equilibrium,
finite-player convergence, and approximate-Nash bounds.
Section~\ref{sec:experiments} presents numerical illustrations and Section~\ref{sec:conclusion} concludes.
The proofs are collected in the appendices.

\paragraph{Notation.}
The symbol \(\odot\) denotes componentwise (Hadamard) multiplication.
We write \(\Delta_m:=\{p\in[0,1]^m:\sum_{k=1}^m p^k=1\}\) for the
probability simplex.
For any dimension \(r\), \(\mathbb S^r\) denotes the real symmetric
\(r\times r\) matrices, with \(\mathbb S_+^r\) and \(\mathbb S_{++}^r\)
their positive semidefinite and positive definite subsets.
For symmetric matrices, \(\mathsf M\preceq\mathsf N\)
(respectively, \(\mathsf M\prec\mathsf N\)) means that
\(\mathsf N-\mathsf M\) is positive semidefinite (respectively, positive
definite).
We use \(\|\cdot\|\) for the Euclidean norm of vectors and any fixed matrix
norm of matrices.

\section{Model}\label{sec:model}

In this section we describe our model for trading around an index reconstitution involving an
\emph{indexer} and $n$ \emph{opportunistic} traders. The indexer must rebalance
its portfolio to match the new index composition and, once that composition is
announced, follows a deterministic trading rule. The opportunists anticipate
the resulting order flow and trade strategically around the event. Thus, the
strategic interaction is entirely among the opportunists.

\subsection{Index Reconstitution Mechanics}\label{sec:index-mechanics}

The rebalance unfolds over a deterministic horizon $[0,T]$ with relevant dates $0<T_{\mathrm{ann}}<T_{\mathrm{impl}}<T$.
At the \emph{announcement time} $T_{\mathrm{ann}}$, the new index constitution becomes
public information. At the \emph{implementation time} $T_{\mathrm{impl}}$, the indexer
must have completed the rebalance and hold the new index portfolio.

We work on a finite eligible universe of $d$
tradable stocks. Here, $d$ denotes the size of the eligible universe, which
contains the current index constituents and a small set of plausible entrants. For
example, an S\&P 500 application might use $d\approx 510$.
The new index reconstitution is characterized by two quantities, both of which
are made public at \(T_{\mathrm{ann}}\).

The first announced quantity is a vector of \emph{post-announcement membership
indicators}. Since this vector is uncertain before the announcement, we model
it as a random variable \(\Xi\) and denote a generic realization of \(\Xi\) by
\(\xi\in\{0,1\}^d\). Here, $\xi_a=1$ means that asset $a$ is included in the
index after the implementation time and $\xi_a=0$ means that it is not. We consider
\(m\) possible reconstitution scenarios \(\{\xi^1,\dots,\xi^m\}\subset\{0,1\}^d\) until
the true index constitution is revealed. Before the announcement, trader \(i\in\{1,\dots,n\}\) evaluates membership
uncertainty under a \emph{subjective probability} \(\mathbb P_n^i\), with
\[
\mathbb P_n^i(\Xi=\xi^k)=\pi_n^{i,k},
\qquad k=1,\dots,m.
\]
We write \(\pi_n^i=(\pi_n^{i,1},\dots,\pi_n^{i,m})\in\Delta_m\) for trader \(i\)'s
prior vector, which may differ across traders, and define the \emph{finite-population
average prior} by
\begin{equation}\label{eq:finite_population_average_prior}
\bar\pi_n:=\frac{1}{n}\sum_{j=1}^n\pi_n^j.
\end{equation}
We assume that \(\bar\pi_n\) is common knowledge. Trader \(i\) knows its own
prior \(\pi_n^i\), but not the individual priors of the other traders.
Traders receive no additional information about
\(\Xi\) before \(T_{\mathrm{ann}}\) and stick to their individual beliefs about
scenario probabilities. At
\(T_{\mathrm{ann}}\), the realized membership vector is publicly revealed, so
all traders then know which scenario has occurred.
We encode this information arrival through the public announcement process
\begin{equation}\label{eq:announcement_process}
A(t):=\Xi\mathbf 1_{\{t\geq T_{\mathrm{ann}}\}},
\qquad 0\leq t\leq T.
\end{equation}

The second announced quantity is a vector of \emph{floating share counts}
$f\in\mathbb{R}_+^d$, where $f_a$ is the number of floating shares of stock
$a$ used in float-adjusted market-cap index calculations.\footnote{Floating
shares are shares available for public trading. Unlike total shares
outstanding, they exclude restricted shares and strategic holdings, such as
those of insiders or controlling shareholders.}
In our model, we treat \(f\) as fixed and known to all players, so \(\Xi\) is the only
uncertain component of the announced index constitution. Accordingly, the
notation for the scenario-dependent indexer quantities below displays only
their dependence on the generic realization \(\xi\) while dependence on fixed model
inputs \(f\) is suppressed.

\begin{remark}[Predictability of Index Membership]\label{rem:membership_predictability}
Russell index membership is largely determined by published eligibility and
market-capitalization ranking rules. Observed capitalizations therefore
provide a direct basis for updating membership forecasts before the ranking
date \citep{PegoraroSammonShim2026}. Although unexpected price movements can
still change outcomes near a ranking cutoff, membership is determined once
all relevant ranking-date inputs are known. As a modeling assumption, one
could let market capitalizations evolve according to stochastic differential
equations with stock-specific volatilities, starting from their current
values. Membership would then be determined from their rankings at the
ranking date. This would admit a
dynamic-programming formulation with additional capitalization state
variables. However, the ranking-based membership rule would generally
destroy the linear--quadratic structure developed later in
Section~\ref{sec:solution}.

By contrast, S\&P 500 selection also involves committee discretion. Public
information can improve forecasts, but observing capitalizations and
eligibility need not resolve the selection decision before the
announcement.\footnote{For example, Datadog's July 2025 addition to the
S\&P 500 surprised some investors who had expected Robinhood or AppLovin
instead. See Silin Chen,
\href{https://www.thestreet.com/investing/analyst-resets-datadog-stock-price-target-after-surprise-addition-to-s-p-500}
{\emph{Analyst Resets Datadog Stock Price Target After Surprise Addition to
S\&P 500}}, TheStreet, July 4, 2025.}
Committee discretion also provides a natural motivation for heterogeneous
membership beliefs: traders observing the same public information may differ
in their assessments of the committee's eventual choice. Our trader-specific
scenario priors represent these differences in judgment.
\end{remark}

Once the realized membership vector \(\xi\) and the vector of floating shares \(f\) have been announced,
the required trade and resulting flow of an indexer replicating the
index can be computed. To do so, let
$h^0\in\mathbb{R}^d$ denote the indexer's pre-announcement holdings
\emph{measured in shares}, so that $h_a^0=0$ for names not held
before the event. After the announcement, the indexer's target stock holdings
are proportional to $f\odot \xi$, i.e.,
\begin{equation}\label{eq:target_holdings_proportional}
h(\xi)=\kappa(\xi)\,(f\odot \xi)
\end{equation}
for some positive proportionality factor $\kappa(\xi)$. In our model this factor is chosen using a
fixed reference-price approximation: Fix
\(s^{\mathrm{ref}}\in\mathbb{R}_{++}^d\), which we usually choose as the initial
midprice \(S(0)\), with \(S(t)\) defined later in \eqref{eq:midprice}, and impose a
self-financing rebalance at this reference price.
Writing \(\Delta h(\xi):=h(\xi)-h^0\) for the indexer's total
rebalance quantity, the self-financing condition is
\begin{equation}\label{eq:cash_neutral_condition}
0 = (s^{\mathrm{ref}})^\top \Delta h(\xi)
  = (s^{\mathrm{ref}})^\top\big(\kappa(\xi)(f\odot \xi)-h^0\big).
\end{equation}
Solving for $\kappa(\xi)$ gives
\begin{equation}
\kappa(\xi)
:=\frac{(s^{\mathrm{ref}})^\top h^0}
{(s^{\mathrm{ref}})^\top(f\odot \xi)}.
\label{eq:kappa_ref}
\end{equation}
Substituting this expression into \(\Delta h(\xi)\) gives the indexer's required
rebalance quantity:
\begin{equation}\label{eq:order_quantity}
\Delta h(\xi)
=\frac{(s^{\mathrm{ref}})^\top h^0}
{(s^{\mathrm{ref}})^\top(f\odot \xi)}(f\odot \xi)-h^0.
\end{equation}
As mentioned before,
we do not model the indexer as a player choosing how to execute this order.
Instead, we assume that it follows a deterministic TWAP that executes the
quantity \(\Delta h(\xi)\) uniformly over
\([T_{\mathrm{TWAP}},T_{\mathrm{impl}}]\), where
\(T_{\mathrm{ann}}<T_{\mathrm{TWAP}}\). Its trading rate is
\begin{equation}\label{eq:twap_rate}
g(t,\xi)
:=
\frac{\Delta h(\xi)}{\Delta_{\mathrm{TWAP}}}
\mathbf 1_{\{T_{\mathrm{TWAP}}\le t\le T_{\mathrm{impl}}\}},
\end{equation}
where \(\Delta_{\mathrm{TWAP}}:=T_{\mathrm{impl}}-T_{\mathrm{TWAP}}\).
For scenario \(k\), define the target holdings and required order by
\begin{equation}\label{eq:scenario_indexer_quantities}
h^k:=h(\xi^k),
\qquad
\Delta h^k:=\Delta h(\xi^k),
\qquad k=1,\dots,m,
\end{equation}
and the corresponding indexer flow by
\begin{equation}\label{eq:scenario_indexer_flows}
g^k(t):=g(t,\xi^k),
\qquad k=1,\dots,m.
\end{equation}
The realized indexer-flow process is
\begin{equation}\label{eq:realized_indexer_flow}
g(t)
:=
\sum_{k=1}^m
\mathbf 1_{\{\Xi=\xi^k\}}\,g^k(t).
\end{equation}

\begin{remark}[TWAP Approximation of the Closing Auction]%
\label{rem:twap_approximation}
We interpret \(T_{\mathrm{impl}}\) as the closing auction on the trading day
before the new index constitution becomes effective. Index funds seeking to
minimize tracking error concentrate their rebalancing trades near the market
close, and many public-index-tracking ETFs execute entirely at the closing-auction price
\citep{Li2021,JiangWuYao2025}. Closing-auction execution aligns the indexer's
transaction prices with the closing prices used to calculate index returns,
avoiding tracking deviations caused by execution at different prices. We
approximate this concentrated auction trade by spreading the required quantity
uniformly over the short interval \([T_{\mathrm{TWAP}},T_{\mathrm{impl}}]\). If
\(\Delta_{\mathrm{TWAP}}\leq1\) and time is measured in trading days, this
interval begins no earlier than the previous close. The indexer therefore
still holds the pre-reconstitution portfolio at the previous close and
completes the rebalance by \(T_{\mathrm{impl}}\).
We discuss a discrete closing-auction extension later in
Remark~\ref{rem:discrete_closing_auction}.
\end{remark}

\begin{remark}[Fixed Reference-Price Approximation]
In practice, transaction prices during the implementation window
\([T_{\mathrm{TWAP}},T_{\mathrm{impl}}]\) fluctuate over time. They are also
affected by the indexer's own trading and by the trades of other market
participants, including the opportunists described in the next subsection.
Using the realized transaction prices rather than the fixed reference price
\(s^{\mathrm{ref}}\) to determine \(\kappa\) would cause the indexer's total
order to depend on prices generated within the model. The indexer's flow would
then have to be determined jointly with the strategic equilibrium. Using
\(s^{\mathrm{ref}}\) is therefore a tractability approximation. It replaces
the realized, endogenous transaction prices with a fixed valuation point, so
that \(\kappa(\xi)\), and hence the indexer flow \(g(t,\xi)\), is known once
the index constitution is announced. This preserves the linear--quadratic
structure used in Section~\ref{sec:solution}.
\end{remark}

\begin{remark}[Simultaneous Rebalancing Across Indices]
In practice, a stock added to one index may be deleted from another at the same
time which is called a migration. The indexer tracking the first index then buys the stock, while the
indexer tracking the second sells it, so their orders partly offset. The same
mechanism applies when a stock is deleted from the index studied here and added
to another index. Our specification abstracts from such simultaneous
rebalancing. Accordingly, \(\Delta h(\xi)\) denotes only the trade required by
the single indexer modeled above and excludes orders from indexers tracking
other indices. Simultaneous rebalancing can be incorporated without changing
the strategic analysis by defining \(g(t,\xi)\) as the net deterministic flow
obtained by summing the orders of all affected indexers.
\end{remark}

\subsection{Trading Dynamics and Objectives}\label{sec:opportunists}

We now specify the opportunists' inventory and price dynamics and their
trading objectives.

\paragraph{Price and Inventory Dynamics.} Let \(S^0=(S^0(t))_{0\le t\le T}\) be an exogenous
\(\mathbb R^d\)-valued process, which we call the \emph{fundamental price}.
Under each trader's subjective probability \(\mathbb P_n^i\), we assume that
\(S^0\) is a square-integrable càdlàg martingale with respect to its natural
filtration and is independent of \(\Xi\), and hence of the announcement process
\(A\).

Trader \(i\in\{1,\dots,n\}\) chooses a trading rate
\(u_n^i(t)\in\mathbb R^d\). Its inventory \(X_n^i(t)\in\mathbb R^d\) evolves as
\begin{equation}\label{eq:inventory_dynamics}
    dX_n^i(t)=u_n^i(t)\,dt,
    \qquad
    X_n^i(0)=x_0^i
\end{equation}
where we usually choose \(x_0^i=0\) so that the opportunists start flat.
We denote the average opportunist inventory by
\begin{equation}\label{eq:average_opportunist_inventory}
\bar X_n(t):=\frac{1}{n}\sum_{j=1}^n X_n^j(t)
\end{equation}
and the aggregate opportunist trading rate and the aggregate
rate of all opportunists other than trader \(i\) by
\begin{equation}\label{eq:aggregate_opportunist_rates}
U_n(t):=\sum_{j=1}^n u_n^j(t),
\qquad
U_n^{-i}(t):=\sum_{j\ne i}u_n^j(t),
\end{equation}
respectively. The \emph{midprice} \(S(t)\in\mathbb{R}^d\) is given by
\begin{equation}\label{eq:midprice}
S(t)=S^0(t)+I(t).
\end{equation}
Here, \(I(t)\in\mathbb R^d\) is the \emph{transient impact} state with dynamics
\begin{equation}\label{eq:transient_impact}
dI(t)
=
\left[
-R I(t)
+\Gamma\left(\varpi_n U_n(t)+g(t)\right)
\right]dt,
\qquad
I(0)=\iota_0
\end{equation}
where \(R, \Gamma\in\mathbb{R}^{d\times d}\) and \(\varpi_n>0\).
We usually choose \(\iota_0=0\), so that the initial impact is zero.
The matrix \(R\) governs the decay of transient impact, while \(\Gamma\) maps
the net trading rate \(\varpi_nU_n+g\) into its contribution to the rate of
change of \(I\). Thus, \(\Gamma_{ab}\) measures the change in the impact rate
of asset \(a\) per unit of net trading rate in asset \(b\). The population
weight \(\varpi_n\) determines the contribution of each opportunist's rate
to transient impact. Typical choices are \(\varpi_n=\varpi\) for a finite-player
game with fixed player size and \(\varpi_n=\varpi/n\) for a model with mean-field scaling,
with \(\varpi>0\) independent of \(n\). We discuss population scaling further in
Subsection~\ref{subsec:parameter_choices_and_restrictions}.

For trader \(i\), trading generates the \emph{instantaneous impact}
\(\Lambda\bigl(u_n^i(t)+\chi_u\varpi_n U_n^{-i}(t)+\chi_g g(t)\bigr)\) where \(\chi_u,\chi_g\geq0\), so
the \emph{execution price} for player \(i\) is
\begin{equation}\label{eq:execution_price}
\widehat S^i(t)
=
S(t)+\Lambda\left(
u_n^i(t)+\chi_u\varpi_n U_n^{-i}(t)+\chi_g g(t)
\right).
\end{equation}
The parameters \(\chi_u\) and \(\chi_g\)
measure the execution-price loadings on the contemporaneous flows of the
other opportunists and the indexer, respectively.
The choice \(\varpi_n=\chi_u=\chi_g=1\) gives all opportunist and indexer
flows equal weight in the execution price, so opposing buy and sell rates
offset. The execution price then is
\[
\widehat S^i(t)
=
S(t)+\Lambda\bigl(U_n(t)+g(t)\bigr),
\]
which is the same for every trader. In the scalar case
\(\Lambda=\lambda>0\) and for a single trader, this specification can be interpreted as walking
through a flat order book of constant height \(\lambda^{-1}\), measured in
shares per unit price. A net contemporaneous flow \(q\) then moves the
execution price by \(\lambda q\). The scaling and interpretation of instantaneous impact are discussed further
later in Subsection~\ref{subsec:parameter_choices_and_restrictions}.

\begin{remark}[Symmetry and Diagonality of the Execution-Cost Matrix]
The symmetry assumption on \(\Lambda\) can be relaxed: the analysis in
Sections~\ref{sec:solution} and \ref{sec:mean_field_limit} extends to
nonsymmetric \(\Lambda\in\mathbb R^{d\times d}\) whose symmetric part
\(\Lambda_{\mathrm{s}}:=\frac{1}{2}(\Lambda+\Lambda^\top)\) is positive
definite.
However, the symmetric specification is also empirically motivated. In particular, as
summarized in Remark~2.1 of \citet{jaber2026}, the findings of
\citet{capponi_cont_2020}, \citet{cont_cucuringu_zhang_2023}, and
\citet{lecoz_mastromatteo_challet_benzaquen_2024} motivate the more restrictive
choice of a diagonal \(\Lambda\).
\end{remark}


\paragraph{Objective.} For a trading-rate profile
\(u_n=(u_n^1,\ldots,u_n^n)\), write
\(u_n^{-i}:=(u_n^j)_{j\ne i}\) and define trader \(i\)'s cost by
\begin{equation}\label{eq:overview_player_objective}
\begin{aligned}
J_n^i(u_n^i;u_n^{-i})
:=
\mathbb E_n^i\Bigg[
&\int_0^T \big(\widehat S^i(t)\big)^\top u_n^i(t)\,dt
-\left((S^0(T))^\top X_n^i(T)-(S^0(0))^\top x_0^i\right)\\
&+\int_0^T \frac{1}{2}(X_n^i(t))^\top Q X_n^i(t)\,dt
+\frac{1}{2}(X_n^i(T))^\top Q_T X_n^i(T)
\Bigg],
\end{aligned}
\end{equation}
where \(Q,Q_T\in\mathbb S_+^d\), and \(\mathbb E_n^i\) denotes expectation
under \(\mathbb P_n^i\). The first integral is the cumulative cash paid for
trading. From this amount we subtract the change in the fundamental value of
the trader's inventory. The final two terms penalize running and terminal
inventory, respectively. We typically choose the terminal penalty sufficiently
large that it effectively forces terminal inventories to be very close to
zero.

\begin{remark}[Fundamental-Price Terminal Valuation]
We value terminal inventory at \(S^0(T)\), rather than at the impacted midprice
\(S(T)\). This avoids crediting a trader with self-generated terminal price
impact that cannot be realized without selling and thereby reversing it. The
convention also underlies the a priori Riccati bounds in
Theorem~\ref{thm:global_riccati_solvability}, which guarantee full-horizon
existence of the closed-loop Nash equilibrium. Our proof of these bounds does
not extend to midprice marking, as explained later in
Remark~\ref{rem:impacted_terminal_mark_solvability}. In the related open-loop
games of \citet{NeumanVoss2023} and
\citet{campbell2026optimalexecutionntraders},
fundamental-price
marking preserves convexity of costs
making first-order conditions sufficient and supporting uniqueness. Additionally, since
\(Q_T\) is typically large, terminal inventories are close to zero and the
choice has little quantitative effect.
\end{remark}

\subsubsection{Parameter Choices and Restrictions}%
\label{subsec:parameter_choices_and_restrictions}

In this subsection, we describe some assumptions we make on the market impact parameters and
also comment on the population weight \(\varpi_n\) in our model.


\paragraph{No Dynamic Arbitrage.} We choose the impact parameters \(\Gamma\) and \(R\) to rule out price
manipulation in the sense of \citet{huberman_stanzl}. Such manipulation means
profiting from price movements generated solely by one's own trades while
ending with the initial portfolio. For this test, all other opportunists and
the indexer are inactive, and the initial transient impact is zero.
Formally, a deterministic trading rate
\(v\in L^2([0,T],\mathbb R^d)\) is a \emph{round trip} if
\(\int_0^T v(t)\,dt=0\). With \(\Lambda\in\mathbb S_{++}^d\) denoting the
execution-cost matrix in \eqref{eq:execution_price}, its total self-induced
cost is
\begin{align*}
\mathcal C_{\mathrm{tot}}(v)
:=
\int_0^T v(t)^\top\Lambda v(t)\,dt+\varpi_n\mathcal C(v),\quad\mathcal C(v)
:=
\int_0^T v(t)^\top
\left(\int_0^t e^{-R(t-s)}\Gamma v(s)\,ds\right)dt.
\end{align*}
Absence of price manipulation requires \(\mathcal C_{\mathrm{tot}}(v)\geq 0\)
for every round trip. We impose the stronger condition
\(\mathcal C(v)\geq 0\) for every
\(v\in L^2([0,T],\mathbb R^d)\). Sufficient conditions for this requirement are
\begin{equation}\label{eq:transient_kernel_parameter_conditions}
R\in\mathbb S_+^d,
\qquad
\Gamma\in\mathbb S_+^d,
\qquad
R\Gamma=\Gamma R.
\end{equation}
Indeed, these conditions give for every
\(v\in L^2([0,T],\mathbb R^d)\) that
\begin{equation}\label{eq:transient_kernel_psd}
\mathcal C(v)
=
\frac{1}{2}\int_0^T\int_0^T
v(t)^\top\Gamma^{1/2}e^{-R\lvert t-s\rvert}
\Gamma^{1/2}v(s)\,ds\,dt
\geq 0.
\end{equation}

See \citet{jaber2026}, \citet{alfonsi2016}, and
\citet{gatheral_2012} for conditions that preclude price
manipulation in more general transient-impact models.

\paragraph{Interpretation of the Population Scaling.}
The weight $\varpi_n$ represents each opportunist's mass. The model variables
$X_n^i$, $u_n^i$, and $J_n^i$ are quantities per unit of trader mass. Their
physical counterparts are
\(X_{n,\mathrm{phys}}^i:=\varpi_nX_n^i\),
\(u_{n,\mathrm{phys}}^i:=\varpi_nu_n^i\), and
\(J_{n,\mathrm{phys}}^i:=\varpi_nJ_n^i\).
Consequently, $dX_{n,\mathrm{phys}}^i(t)
=u_{n,\mathrm{phys}}^i(t)\,dt$ and
$\varpi_nU_n=\sum_{j=1}^n u_{n,\mathrm{phys}}^j$. The transient-impact
dynamics can therefore be written in physical units as
\[
dI(t)
=
\left[
-RI(t)
+
\Gamma\left(
\sum_{j=1}^n u_{n,\mathrm{phys}}^j(t)+g(t)
\right)
\right]dt.
\]
Thus, for every $\varpi_n>0$, $R$ and $\Gamma$ act on aggregate physical order
flow and retain the same interpretation.

The interpretation of the private instantaneous and inventory costs is
different. In physical units, the execution price becomes
\[
\widehat S^i(t)
=
S(t)
+
\Lambda\left(
\frac{1}{\varpi_n}u_{n,\mathrm{phys}}^i(t)
+
\chi_u\sum_{j\ne i}u_{n,\mathrm{phys}}^j(t)
+
\chi_g g(t)
\right).
\]
Moreover, multiplying the objective by $\varpi_n$ and substituting
$X_n^i=X_{n,\mathrm{phys}}^i/\varpi_n$ and
$u_n^i=u_{n,\mathrm{phys}}^i/\varpi_n$ gives
\begin{align*}
J_{n,\mathrm{phys}}^i
=
\mathbb E_n^i\Bigg[
&
\int_0^T
\bigl(\widehat S^i(t)\bigr)^\top
u_{n,\mathrm{phys}}^i(t)\,dt
-
\left(
(S^0(T))^\top X_{n,\mathrm{phys}}^i(T)
-(S^0(0))^\top X_{n,\mathrm{phys}}^i(0)
\right)
\\
&
+
\frac{1}{2\varpi_n}\int_0^T
\bigl(X_{n,\mathrm{phys}}^i(t)\bigr)^\top
QX_{n,\mathrm{phys}}^i(t)\,dt
\\
&
+
\frac{1}{2\varpi_n}
\bigl(X_{n,\mathrm{phys}}^i(T)\bigr)^\top
Q_TX_{n,\mathrm{phys}}^i(T)
\Bigg].
\end{align*}
Hence a trader of mass $\varpi_n$ faces coefficients
$\Lambda/\varpi_n$, $Q/\varpi_n$, and $Q_T/\varpi_n$
on its physical trading rate and inventory, so the fixed parameters $\Lambda$,
$Q$, and $Q_T$ are best interpreted per unit of trader mass. For the trading
term, $\Lambda/\varpi_n$ means that a fixed absolute rate becomes more expensive as the
trader becomes smaller, as would be the case if execution capacity were
proportional to trader mass. For the inventory terms, $Q/\varpi_n$ and $Q_T/\varpi_n$ likewise
make a fixed absolute position more expensive. This amounts to greater
effective risk aversion in physical units and is reasonable under the assumption
that absolute risk-bearing capacity is proportional to trader mass.

If instead one requires the coefficients on absolute physical rates and
inventories to remain constant, one must take
$\Lambda_n=\varpi_n\Lambda_{\mathrm{phys}}$,
$Q_n=\varpi_nQ_{\mathrm{phys}}$, and
$Q_{T,n}=\varpi_nQ_{T,\mathrm{phys}}$. Under $\varpi_n=\varpi/n$,
these coefficients vanish as $n\to\infty$. In particular, the representative-agent Hamiltonian loses its
strictly convex quadratic term in the trading rate, while the running and
terminal inventory-restoring terms disappear. The limiting mean-field control
problem is therefore generally singular.

As discussed in Remark~\ref{rem:twap_approximation}, modeling the indexer's execution as a TWAP
is a simplification. In practice, index funds concentrate much of their
rebalance trading in the closing auction, where opposing buy and sell orders
offset in the auction imbalance. We therefore typically set $\chi_u=\chi_g=1$
to approximate the offsetting of buy and sell orders in the closing auction
within our continuous-trading game. With this choice, the trader's own rate
$u_n^i$, other opportunists' aggregate physical flow $\varpi_nU_n^{-i}$, and
the indexer's flow $g$ enter the execution price through the same matrix
$\Lambda$, allowing opposing contributions to offset in the execution-cost term.

In the atomic specification \(\varpi_n=1\), this choice $\chi_u=\chi_g=1$ gives
$\widehat S^i=S+\Lambda(U_n+g)$, so every trader faces the same execution price.
Opposite physical trading rates in the same asset cancel
in $U_n+g$ at each time.
For fixed $\varpi_n=\varpi\ne1$, the trader's own physical rate retains the
coefficient $\Lambda/\varpi$, so execution prices generally differ across traders.

Under the mean-field scaling \(\varpi_n=\varpi/n\), setting
$\chi_u=\chi_g=1$ gives the limiting execution price
\[
\widehat S^i(t)
=
S(t)
+
\Lambda\bigl(
u^i(t)+\varpi\bar u(t)+g(t)
\bigr),
\qquad
\bar u(t)
=
\lim_{n\to\infty}\frac{1}{n}U_n(t).
\]
Unlike under \(\varpi_n=1\), this expression does not represent exact
cancellation of physical flows even when \(\varpi=1\): $u_n^i$ enters
unscaled, whereas the other opportunists' aggregate rate enters as
$\varpi U_n^{-i}/n$. Thus, $u^i$ is the trader's intensive own rate, while
$\varpi\bar u+g$ is aggregate physical flow. This asymmetry is necessary to avoid the
singular mean-field limit discussed earlier in this subsection. Letting
$u^i$, $\varpi\bar u$, and $g$ enter through the same matrix $\Lambda$ is our best
approximation of auction netting in the mean-field model: $\varpi\bar u$ and $g$ can
still offset each other, and opposing aggregate flow reduces the execution-price
displacement faced by an individual trader, even though the trader's own
physical flow does not cancel one-for-one as it does under \(\varpi_n=1\).
In the mean-field model, $\varpi$ serves as a knob for offsetting: increasing
it gives greater weight to aggregate opportunist flow.

\begin{remark}[Discrete Closing-Auction Extension]
\label{rem:discrete_closing_auction}
A more literal treatment of the implementation-date auction would let the indexer
trade only in a discrete closing auction. Trader \(i\) would submit an auction
quantity \(u_{\mathrm A}^i\in\mathbb R^d\), and
\(U_{\mathrm A}:=\sum_{j=1}^n u_{\mathrm A}^j\) would denote the
opportunists' aggregate auction quantity. For example, one could specify
\[
P_{\mathrm{auction}}
=
S(T_{\mathrm{impl}}-)
+
H_{\mathrm A}
\bigl(
U_{\mathrm A}+g_{\mathrm A}+\xi_{\mathrm A}
\bigr),
\qquad
\mathbb E[\xi_{\mathrm A}]=0,
\]
where \(H_{\mathrm A}\) is the auction-impact matrix,
\(g_{\mathrm A}\) is the indexer's auction quantity, and
\(\xi_{\mathrm A}\) is exogenous auction imbalance. Since
\(U_{\mathrm A}=u_{\mathrm A}^i+\sum_{j\ne i}u_{\mathrm A}^j\), trader
\(i\)'s auction payment \(P_{\mathrm{auction}}^\top u_{\mathrm A}^i\) is
quadratic in the submitted quantities. The linear--quadratic structure would
therefore largely be preserved, but the solution would have to stitch together
three continuous-time LQ games---before the announcement, between the
announcement and the implementation-date auction, and after the auction---with the
discrete LQ auction game at \(T_{\mathrm{impl}}\). Working backward, the
post-auction continuation value enters the auction game. The auction equilibrium
value then supplies the quadratic terminal condition for the continuous game
preceding the auction, which must in turn be linked to the pre-announcement
game. The main difficulty is proving an analogue of
Theorem~\ref{thm:global_riccati_solvability}: the auction-induced terminal
matrix must generate a Riccati solution that remains finite throughout the
earlier intervals. We expect such a solvability result to hold, but proving it
would be substantially more cumbersome. We therefore retain the continuous
trading approximation of the auction-dominated environment around the
implementation date.
\end{remark}

\paragraph{Relative Weights in the Execution Price.}
For each finite-player game, we impose $\chi_u\varpi_n\leq1$, so each other
opportunist's rate enters \eqref{eq:execution_price} with no greater weight
than the trader's own rate. Together with the remaining conditions collected
in the following assumption, this restriction is sufficient to prevent
blow-up of the equilibrium coefficients, as shown later in
Theorem~\ref{thm:global_riccati_solvability}. Usually one chooses \(\chi_u\in\{0, 1\}\)
and either \(\varpi_n=1\) or \(\varpi=\varpi/n\). The restriction \(\chi_u\varpi_n\leq1\)
thus only becomes relevant in the mean-field scaling setting with \(\chi_u=1\) where it is equivalent to
\(n\geq\varpi\).

\begin{assumption}[Standing Assumptions]\label{ass:standing_model}
Under every \(\mathbb P_n^i\), \(S^0\) is a square-integrable càdlàg martingale
independent of \(\Xi\). Moreover, \(\Lambda\in\mathbb S_{++}^d\),
\(Q,Q_T\in\mathbb S_+^d\), \(\mathcal C(v)\geq0\) for every deterministic
\(v\in L^2([0,T],\mathbb R^d)\), \(\chi_u,\chi_g\geq0\), and
\(\varpi_n>0\). In each finite-player game, assume
\(\chi_u\varpi_n\leq1\).
\end{assumption}
\subsection{Admissible Closed-Loop Policies}

We next define the admissible controls, focusing on \emph{closed-loop policies}. In
general, a closed-loop policy maps the history of the processes observed by a
player up to time \(t\) to that player's current control. In our model, trader
\(i\) observes their own inventory \(X_n^i\), the average inventory \(\bar X_n\),
the midprice \(S\), and the announcement process \(A\). Importantly, traders do
not observe the individual inventories \(X_n^j\) of the other traders
\(j\ne i\). The formal definition is as follows.

\begin{definition}[Admissible Closed-Loop Policies]
\label{def:admissible_feedback_policies}
A profile \(\Phi_n=(\Phi_n^1,\ldots,\Phi_n^n)\) determines the induced
trading rates
\begin{equation}\label{eq:finite_feedback_induced_rate}
u_n^i(t)
:=
\Phi_n^i\left(
t,X_n^i,\bar X_n,S,A;\bar\pi_n,\pi_n^i
\right),
\qquad i=1,\ldots,n,
\end{equation}
where each \(\Phi_n^i\) is a Borel measurable map
\begin{equation}\label{eq:finite_feedback_policy_maps}
\Phi_n^i:
[0,T]\times D([0,T];\mathbb R^d)^4\times\Delta_m^2
\longrightarrow\mathbb R^d,
\qquad i=1,\ldots,n,
\end{equation}
where \(D([0,T];\mathbb R^d)\) denotes the space of \(\mathbb R^d\)-valued
càdlàg paths, equipped with the Skorokhod topology. The four path arguments
are ordered as \(\mathsf h=(\mathbf{x}^i,\bar{\mathbf{x}},\mathbf{s},\mathbf{a})\), corresponding
to the histories of \((X_n^i,\bar X_n,S,A)\), and the final two arguments are
ordered as \((\bar p,p^i)\), corresponding to the average prior and trader
\(i\)'s own prior. Each map \(\Phi_n^i\) is non-anticipative in its path
arguments: for all \(t\in[0,T]\), all paths \(\mathsf h,\widetilde{\mathsf h}\), and all
\((\bar p,p^i)\in\Delta_m^2\),
\begin{equation}\label{eq:finite_feedback_nonanticipativity}
\mathsf h|_{[0,t]}=\widetilde{\mathsf h}|_{[0,t]}
\quad\Longrightarrow\quad
\Phi_n^i(t,\mathsf h;\bar p,p^i)
=\Phi_n^i(t,\widetilde{\mathsf h};\bar p,p^i).
\end{equation}
The profile \(\Phi_n\) is admissible if the system formed by
\eqref{eq:inventory_dynamics}, \eqref{eq:midprice},
\eqref{eq:transient_impact}, and
\eqref{eq:finite_feedback_induced_rate} admits a unique strong solution and the
induced rates are square-integrable under every trader's subjective
probability:
\[
\mathbb E_n^i\left[
\int_0^T \sum_{j=1}^n \|u_n^j(t)\|^2\,dt
\right]
<\infty,
\qquad i=1,\dots,n.
\]
The set of admissible policy profiles is denoted by \(\mathcal U_n\).
\end{definition}

Above we defined policies as functions of the histories of
\((X_n^i,\bar X_n,S,A)\), because this is economically natural
as \(S\) is the actual midprice observed by traders. However,
from \eqref{eq:midprice} and \eqref{eq:transient_impact}, it is easy to see that
observing the histories of \((X_n^i,\bar X_n,S,A)\) is equivalent to observing
those of \((X_n^i,\bar X_n,I,A,S^0)\). We could therefore equivalently work
with policies of the form
\begin{equation}\label{eq:equivalent_history_policy_parametrization}
\Phi_n^i(t,\mathbf{x}^i,\bar{\mathbf{x}},\boldsymbol{\iota},\mathbf{a},\mathbf{s}^0;\bar p,p^i).
\end{equation}
The Nash equilibrium derived
in the next section will in fact consist of Markov policies in this alternative
history parametrization that ignore \(S^0\):
For current values \(x^i\), \(\bar x\), and \(\iota\), define trader
\(i\)'s \emph{deviation state} by \(\delta^i:=x^i-\bar x\) and call
\((\bar x,\iota)\) the \emph{common state}. We define the corresponding \emph{reduced state} by
\begin{equation}\label{eq:reduced_state}
z^i(x^i,\bar x,\iota)
:=
\bigl((\delta^i)^\top,\bar x^\top,\iota^\top\bigr)^\top
\in\mathbb R^{3d}
\end{equation}
and consider scenario-indexed feedback policies
\begin{equation}\label{eq:markov_feedback}
\phi_n^{i,k}(t,z^i;\bar\pi_n,\pi_n^i),
\qquad k=0,\ldots,m.
\end{equation}
For \(x_n=(x_n^1,\ldots,x_n^n)\), write
\(\bar x_n:=n^{-1}\sum_{j=1}^n x_n^j\).
Trader \(i\) uses \(\phi_n^{i,0}\) before the announcement and
\(\phi_n^{i,k}\) after scenario \(k\) is announced. The policy map on
histories of \((X_n^i,\bar X_n,I,A,S^0)\) induced by this feedback family is
\begin{equation}\label{eq:markov_feedback_lift}
\begin{aligned}
&\Phi_n^i(t,\mathbf{x}^i,\bar{\mathbf{x}},\boldsymbol{\iota},\mathbf{a},\mathbf{s}^0;
\bar p,p^i)
\\
&\quad:=
\begin{cases}
\phi_n^{i,0}
\bigl(t,z^i(\mathbf{x}^i(t),\bar{\mathbf{x}}(t),\boldsymbol{\iota}(t));\bar\pi_n,\pi_n^i\bigr),
&t<T_{\mathrm{ann}},\\[2pt]
\phi_n^{i,k}
\bigl(t,z^i(\mathbf{x}^i(t),\bar{\mathbf{x}}(t),\boldsymbol{\iota}(t));\bar\pi_n,\pi_n^i\bigr),
&t\geq T_{\mathrm{ann}},\ \mathbf{a}(t)=\xi^k.
\end{cases}
\end{aligned}
\end{equation}
Write \(\phi_n^i:=(\phi_n^{i,k})_{k=0}^m\) and
\(\phi_n:=(\phi_n^1,\ldots,\phi_n^n)\). Henceforth, when the Markov feedback
family \(\phi_n\) appears where a history-dependent policy profile is
required, it refers to the profile induced by
\eqref{eq:markov_feedback_lift}.
Because these Markov
feedbacks ignore \(S^0\), the induced states and rates are deterministic
conditional on each membership scenario. Thus, their square-integrability condition for membership
in \(\mathcal U_n\) reduces to deterministic \(L^2\)-integrability along the
finitely many induced paths for the \(m\) different scenarios.

For an admissible profile \(\Phi_n\), we write
\(J_n^i(\Phi_n^i;\Phi_n^{-i})\) for the objective in
\eqref{eq:overview_player_objective} evaluated at its induced rates and states.
An admissible deviation by trader \(i\) is a policy \(\Psi_n^i\) such
that \((\Psi_n^i,\Phi_n^{-i})\in\mathcal U_n\).

\begin{definition}[Closed-Loop Nash Equilibrium]
\label{def:nash_equilibrium}
An admissible policy profile \(\Phi_n^*\in\mathcal U_n\) is a closed-loop
Nash equilibrium if, for every trader \(i\) and every admissible deviation
\(\Psi_n^i\),
\begin{equation}\label{eq:ne_definition}
J_n^i(\Phi_n^{i,*};\Phi_n^{-i,*})
\leq
J_n^i(\Psi_n^i;\Phi_n^{-i,*}).
\end{equation}
\end{definition}

\begin{remark}[Closed-Loop vs Open-Loop Controls]
A comment on our choice of closed-loop rather than \emph{open-loop controls} is in
order. Open-loop controls, commonly used in the stochastic differential game literature
on market impact games, are
trading-rate processes adapted to the filtration generated by \(S^0\) and
\(A\). Since neither process is affected by the traders' actions, once an
open-loop profile has been chosen, the realized trading rates of the
nondeviating traders remain unchanged when another trader deviates. In our
closed-loop formulation, by contrast, policies are functions of observable
histories. A unilateral deviation changes the histories fed into the other
traders' policy functions. Those functions remain fixed, but their different
inputs generally produce different realized trading rates. This may be more
realistic because traders observe the resulting changes in the average
inventory and the midprice and can adjust their subsequent trading accordingly.
\end{remark}

\begin{remark}[A Market-Data Proxy for Average Inventory]
\label{rem:average_inventory_proxy}
In the definition of the admissible policies we treated \(\bar X_n\)
as observed by traders. In practice, however, \(\bar X_n\) is not directly
reported. Nevertheless, a reasonable proxy can often still be
constructed from market data if one is willing to make certain assumptions.
We briefly describe a crude procedure to compute such a proxy:
Suppose that the opportunists trade exclusively through liquidity-taking
market orders. Their aggregate inventory change can then be approximated by cumulative signed
liquidity-taking order flow. This flow can be estimated from consolidated
transaction and quotation data using the heuristic of \citet{lee_ready_1991},
which classifies trades above the prevailing quote midpoint as buyer-initiated,
trades below it as seller-initiated, and midpoint trades using a tick test.
Summing the signed quantities and subtracting the known cumulative indexer flow
would give a proxy for cumulative opportunist flow. Adding \(1/n\) times this
quantity to the initial estimate \(\bar X_n(0)\) would then yield a proxy for
\(\bar X_n\). Related empirical work estimates the flows of particular investor
groups by combining signed
transactions with observable transaction characteristics or lower-frequency
holdings data
\citep{campbell_ramadorai_schwartz_2009,boehmer_jones_zhang_zhang_2021}.
Possible limitations of this proxy are that (i) opportunists may provide
liquidity through resting limit orders, so their trades are signed according to
their counterparties' aggressive orders; (ii) aggressive orders from other
market participants are included; and (iii) auctions, crosses, and some
off-exchange or midpoint transactions have no reliably identifiable
liquidity-taking side. The model abstracts from these measurement errors and
treats \(\bar X_n\) as observed.
\end{remark}

\section{Solution of the Trading Game}\label{sec:solution}

In this section, we construct a subgame-perfect closed-loop Nash equilibrium
consisting of Markovian policies of the form \eqref{eq:markov_feedback}.

\subsection{Equilibrium Verification}
\label{sec:fundamental_price_reduction}

We start by removing the direct dependence of the objective in
\eqref{eq:overview_player_objective} on the fundamental price.
Fix an admissible profile \(\Phi_n\in\mathcal U_n\). By
\eqref{eq:inventory_dynamics}, \(X_n^i\) is continuous, adapted, and of finite
variation, and its quadratic covariation with \(S^0\) vanishes. Integration by parts therefore gives
\begin{equation}\label{eq:fundamental_price_integration_by_parts}
\begin{aligned}
&\int_0^T (S^0(t))^\top u_n^i(t)\,dt
-\left((S^0(T))^\top X_n^i(T)-(S^0(0))^\top x_0^i\right)
\\
&\qquad
=\int_{(0,T]}(S^0(t-))^\top\,dX_n^i(t)
-\left((S^0(T))^\top X_n^i(T)-(S^0(0))^\top x_0^i\right)
\\
&\qquad
=-\int_{(0,T]}(X_n^i(t-))^\top\,dS^0(t).
\end{aligned}
\end{equation}
By Lemma~\ref{lem:fundamental_price_stochastic_integral}, the last integral in
\eqref{eq:fundamental_price_integration_by_parts} is the terminal value of a
true martingale starting from zero and therefore has zero expectation.
Substituting this into \eqref{eq:overview_player_objective} gives
\begin{equation}\label{eq:reformulated_objective}
\begin{aligned}
J_n^i(\Phi_n^i;\Phi_n^{-i})
={}
\mathbb E_n^i\Bigg[&
\int_0^T
\mathcal L\left(
X_n^i(t),u_n^i(t);
I(t),\varpi_nU_n^{-i}(t),g(t)
\right)\,dt
\\
&+\frac{1}{2}(X_n^i(T))^\top Q_TX_n^i(T)
\Bigg]
\end{aligned}
\end{equation}
where
\begin{equation}\label{eq:reduced_running_cost}
\mathcal L(x,u;\iota,\mu,g)
:=
\bigl(\iota+\Lambda(u+\chi_u\mu+\chi_g g)\bigr)^\top u
+\frac{1}{2}x^\top Q x.
\end{equation}
Neither \eqref{eq:reformulated_objective} nor the dynamics of \(X_n\) and
\(I\) depend explicitly on \(S^0\). However, an admissible policy may still
use the observed history of \(S^0\) as an external randomization device. The
verification theorem \ref{thm:stitched_hjb_verification} below shows that sufficiently
smooth solutions of a coupled HJB system generate equilibrium policies that
ignore \(S^0\) as an argument
and remain robust against unilateral deviations from the full class
\(\mathcal U_n\) which is allowed to depend on \(S^0\).

Since we aim to solve for subgame-perfect Nash equilibria using dynamic
programming, we introduce continuation games from intermediate times and
states. A continuation game is specified by a tuple
\((t,x_n,\iota,a,\mathbf{s}^0_{[0,t]})\), consisting of the current time \(t\), the
inventory profile \(X_n(t)=x_n\), the impact state \(I(t)=\iota\), the
announcement state \(A(t)=a\), and the \emph{full history} \(\mathbf{s}^0_{[0,t]}\) of
the fundamental price. The full history of \(S^0\) is retained because its
conditional future law may depend on its past as \(S^0\) is not assumed to be Markov. The running cost and controlled
state equations have no direct dependence on the pre-\(t\) histories of
\(X_n\) and \(I\). Once these histories are fixed, any policy dependence on
them can be absorbed into the continuation policies, so only
\((x_n,\iota)\) must be recorded. The announcement process is
Markov, so its history is summarized by \((t,a)\).

At each \(\tau\in[t,T]\), an admissible continuation policy of trader \(i\)
is evaluated nonanticipatively on the histories of \(X_n^i\), \(\bar X_n\),
\(I\), and \(A\) over \([t,\tau]\), together with the history of \(S^0\) over
the larger interval \([0,\tau]\). Trader \(i\)'s continuation objective is
\begin{equation}\label{eq:reduced_continuation_objective}
\begin{aligned}
&J_n^i\left(
\Phi_n^i;\Phi_n^{-i}
\mid t,x_n,\iota,a,\mathbf{s}^0_{[0,t]}
\right)
\\
:={}&\mathbb E_n^i\Bigg[
\int_t^T
\mathcal L\left(
X_n^i(\tau),u_n^i(\tau);
I(\tau),\varpi_nU_n^{-i}(\tau),g(\tau)
\right)\,d\tau
\\
&\hspace{72pt}
+\frac{1}{2}(X_n^i(T))^\top Q_TX_n^i(T)
\,\Bigm|\,
S^0|_{[0,t]}=\mathbf{s}^0_{[0,t]},\ A(t)=a
\Bigg]
\end{aligned}
\end{equation}
where on \([t,T]\) the induced state processes satisfy
\begin{equation}\label{eq:continuation_state_dynamics}
\begin{aligned}
dX_n^i(\tau)
&=u_n^i(\tau)\,d\tau,
&X_n^i(t)&=x_n^i,
&i&=1,\ldots,n,\\
dI(\tau)
&=\left[-RI(\tau)
+\Gamma\left(
\varpi_n\sum_{j=1}^n u_n^j(\tau)+g(\tau)
\right)\right]d\tau,
&I(t)&=\iota.
\end{aligned}
\end{equation}
The announcement process continues according to
\eqref{eq:announcement_process}.

\begin{definition}[Subgame-Perfect Nash Equilibrium]
\label{def:subgame_perfect_nash_equilibrium}
An admissible profile \(\Phi_n^*\in\mathcal U_n\) is a subgame-perfect Nash
equilibrium if, in every continuation game starting from
\((t,x_n,\iota,a,\mathbf{s}^0_{[0,t]})\), for every trader
\(i\) and every admissible unilateral continuation
deviation \(\Psi_n^i\),
\begin{equation}\label{eq:subgame_perfect_ne}
J_n^i\left(
\Phi_n^{i,*};\Phi_n^{-i,*}
\mid t,x_n,\iota,a,\mathbf{s}^0_{[0,t]}
\right)
\leq
J_n^i\left(
\Psi_n^i;\Phi_n^{-i,*}
\mid t,x_n,\iota,a,\mathbf{s}^0_{[0,t]}
\right).
\end{equation}
The left-hand side of \eqref{eq:subgame_perfect_ne} is called trader \(i\)'s
value of the continuation game and is denoted by
\(V_n^i(t,x_n,\iota,a,\mathbf{s}^0_{[0,t]})\).
\end{definition}


We next motivate a verification theorem tailored to the announcement structure
of the index-rebalancing game. The candidate Markov feedbacks respect the
traders' observation structure and induce deterministic trajectories
conditional on the membership scenario.

For the Markov feedbacks in \eqref{eq:markov_feedback}, write
\(\phi_n^{-i,k}:=(\phi_n^{j,k})_{j\ne i}\). Define the aggregate feedback rate of trader
\(i\)'s opponents in each post-announcement scenario \(k=1,\ldots,m\) by
\[
U_n^{-i,k}(t,x_n,\iota)
:=\sum_{j\ne i}\phi_n^{j,k}
\left(t,z^j(x_n^j,\bar x_n,\iota);\bar\pi_n,\pi_n^j\right),
\]
and define \(U_n^{-i,0}\) analogously using the pre-announcement feedbacks
\(\phi_n^{j,0}\).
Let \(Z_n^i(t):=z^i(X_n^i(t),\bar X_n(t),I(t))\) denote trader \(i\)'s
reduced-state process. A full inventory vector \(x_n\) is \emph{compatible}
with a generic reduced state
\(z^i=((\delta^i)^\top,\bar x^\top,\iota^\top)^\top\) if
\(x_n^i=\delta^i+\bar x\) and \(\bar x_n=\bar x\). In each post-announcement
scenario \(k=1,\ldots,m\), the dynamics of \(Z_n^i\) are
\[
\dot Z_n^i(t)
=\widehat b_n^{i,k}\left(t,X_n(t),I(t),u_n^i(t)\right),
\]
where, for a full inventory vector \(x_n\), an impact state \(\iota\), and a
candidate trading rate \(v\in\mathbb R^d\), the drift is
\[
\widehat b_n^{i,k}(t,x_n,\iota,v)
:=\begin{bmatrix}
\frac{n-1}{n}v-\frac{1}{n}U_n^{-i,k}(t,x_n,\iota)\\[2pt]
\frac{1}{n}\bigl(v+U_n^{-i,k}(t,x_n,\iota)\bigr)\\[2pt]
-R\iota+\Gamma\bigl(
\varpi_n(v+U_n^{-i,k}(t,x_n,\iota))+g^k(t)
\bigr)
\end{bmatrix}.
\]
The same dynamics hold before the announcement with \(k=0\), under the
convention \(g^0\equiv0\).
For \(p\in\mathbb R^{3d}\), the costate conjugate to the reduced state, the
corresponding Hamiltonian is
\begin{equation}\label{eq:stitched_hamiltonian}
\begin{aligned}
\widehat{\mathscr H}_n^{i,k}(t,x_n,\iota,v,p;\phi_n^{-i,k})
:={}&
\mathcal L\left(
x_n^i,v;\iota,
\varpi_nU_n^{-i,k}(t,x_n,\iota),
g^k(t)
\right)
\\
&+p^\top\widehat b_n^{i,k}(t,x_n,\iota,v).
\end{aligned}
\end{equation}
The dependence on \(\phi_n^{-i,k}\) enters through \(U_n^{-i,k}\) in both
the running cost and the reduced-state drift.

Fix trader \(i\) and the opponents' feedbacks
\((\phi_n^{-i,k})_{k=0}^m\). We seek smooth candidate value functions
\(V_n^{i,k}(t,z^i)\) satisfying the following HJB equation away from the
announcement and switching times, for each
\(k=0,\ldots,m\) and every \(x_n\) compatible with \(z^i\),
\begin{equation}\label{eq:single_player_best_response_hjb}
0
=
\partial_tV_n^{i,k}(t,z^i)
+\min_{v\in\mathbb R^d}
\widehat{\mathscr H}_n^{i,k}\left(
t,x_n,\iota,v,
\nabla V_n^{i,k}(t,z^i);\phi_n^{-i,k}
\right)
\end{equation}
with terminal and announcement-stitching conditions
\[
\begin{aligned}
V_n^{i,k}(T,z^i)
&=\frac{1}{2}(\delta^i+\bar x)^\top Q_T(\delta^i+\bar x),
&&k=1,\ldots,m,\\
V_n^{i,0}(T_{\mathrm{ann}},z^i)
&=\sum_{k=1}^m\pi_n^{i,k}V_n^{i,k}(T_{\mathrm{ann}},z^i).
\end{aligned}
\]
If the minimizer of \eqref{eq:single_player_best_response_hjb} depends on
\(x_n\) only through \(z^i\) and defines an admissible feedback policy
\(\phi_n^i\) of the form \eqref{eq:markov_feedback}, verification shows that
this policy is an optimal response to the fixed opponents' feedbacks.
Requiring the HJB equations and minimizing-feedback conditions to hold
simultaneously for all traders yields a coupled system. We first solve the
coupled post-announcement HJBs separately for each scenario. Trader \(i\)'s
\(\pi_n^i\)-weighted average of the post-announcement value functions at
\(T_{\mathrm{ann}}\) then supplies the terminal condition for its
pre-announcement HJB.

\begin{theorem}[HJB Verification]
\label{thm:stitched_hjb_verification}
Fix a Markov feedback profile
\(\phi_n\) in the parametrization of \eqref{eq:markov_feedback} that is
admissible in every continuation game. Suppose that, for every trader \(i\),
there exist candidate functions
\(V_n^{i,0}:[0,T_{\mathrm{ann}}]\times\mathbb R^{3d}\to\mathbb R\) and
\(V_n^{i,k}:[T_{\mathrm{ann}},T]\times\mathbb R^{3d}\to\mathbb R\),
\(k=1,\ldots,m\), such that every \(V_n^{i,k}\), \(k=0,\ldots,m\), is
continuous and \(C^1\) between the switching times of the corresponding
\(g^k\). For every trader \(i\), assume the HJB and minimizing-feedback
conditions below hold at
every time between switching times, for every
\(z^i=((\delta^i)^\top,\bar x^\top,\iota^\top)^\top\in\mathbb R^{3d}\) and
every \(x_n\) compatible with \(z^i\), while the boundary conditions hold for
every \(z^i\in\mathbb R^{3d}\):
\begin{subequations}\label{eq:stitched_hjb_system}
\begin{equation}\label{eq:stitched_hjb}
0
=
\partial_tV_n^{i,k}(t,z^i)
+\min_{v\in\mathbb R^d}
\widehat{\mathscr H}_n^{i,k}\left(
t,x_n,\iota,v,
\nabla V_n^{i,k}(t,z^i);\phi_n^{-i,k}
\right),
\end{equation}
\begin{align}
\phi_n^{i,k}
\left(t,z^i;\bar\pi_n,\pi_n^i\right)
&=
\operatorname*{argmin}_{v\in\mathbb R^d}
\widehat{\mathscr H}_n^{i,k}\left(
t,x_n,\iota,v,
\nabla V_n^{i,k}(t,z^i);\phi_n^{-i,k}
\right),
\label{eq:stitched_hjb_argmin}\\
V_n^{i,k}(T,z^i)
&=
\frac{1}{2}(\delta^i+\bar x)^\top Q_T(\delta^i+\bar x),
\qquad k=1,\ldots,m,
\label{eq:stitched_hjb_terminal}\\
V_n^{i,0}(T_{\mathrm{ann}},z^i)
&=
\sum_{k=1}^m\pi_n^{i,k}
V_n^{i,k}(T_{\mathrm{ann}},z^i).
\label{eq:stitched_hjb_value_matching}
\end{align}
\end{subequations}
Here \eqref{eq:stitched_hjb}--\eqref{eq:stitched_hjb_argmin} hold for
\(k=0\) before the announcement and for \(k=1,\ldots,m\) after it.
Then \(\phi_n\) is a Markov subgame-perfect Nash equilibrium. Its continuation
value before the announcement satisfies
\[
V_n^i(s,x_n,\iota,0,\mathbf{s}^0_{[0,s]})
=V_n^{i,0}\left(s,z^i(x_n^i,\bar x_n,\iota)\right).
\]
After scenario \(k\) is announced, it satisfies
\[
V_n^i(s,x_n,\iota,\xi^k,\mathbf{s}^0_{[0,s]})
=V_n^{i,k}\left(s,z^i(x_n^i,\bar x_n,\iota)\right).
\]
\end{theorem}

The traders' information restriction is encoded in both the candidate
functions and the HJB system. For each \(k=0,\ldots,m\), \(V_n^{i,k}\) and
\(\phi_n^{i,k}\) depend on the current state only through the reduced state
\(z^i\). The minimum and argmin in
\eqref{eq:stitched_hjb}--\eqref{eq:stitched_hjb_argmin} range over all
\(v\in\mathbb R^d\). Because the Hamiltonian depends on \((x_n,\iota)\), its
minimizer could a priori be a general function of these variables.
Equation~\eqref{eq:stitched_hjb_argmin} instead requires the minimizer to be
the same for every \(x_n\) compatible with a given \(z^i\), so it depends on
the current state only through \(z^i\). Equation~\eqref{eq:stitched_hjb}
imposes the analogous restriction on the minimum value. Since the Hessian with
respect to \(v\) is \(2\Lambda\succ0\), the minimizer is unique.
Note that the theorem also shows that \(\phi_n\) remains a
subgame-perfect Nash equilibrium under the larger information structure where
each trader observes every other
trader's inventory. The proof is given in
Appendix~\ref{subsec:proof_stitched_hjb_verification}.

\begin{remark}[Closed-Loop and Open-Loop Stitching]
The coupled HJB system \eqref{eq:stitched_hjb_system} yields a natural
backward-induction procedure. First solve the post-announcement HJBs for
\(V_n^{i,k}\), \(k=1,\ldots,m\). Their \(\pi_n^i\)-weighted values at
\(T_{\mathrm{ann}}\) then provide, through
\eqref{eq:stitched_hjb_value_matching}, the terminal condition for the
pre-announcement HJB for \(V_n^{i,0}\). The \(m\) post-announcement value
functions \(V_n^{i,k}\) can thus be computed in parallel before solving the
pre-announcement problem.
\end{remark}

\subsection{Quadratic Solution and Backward Construction}
\label{sec:fixed_interval_lq_game}
\label{sec:full_procedure}

We now explicitly solve the coupled HJB system in
Theorem~\ref{thm:stitched_hjb_verification} using the quadratic ansatz
\begin{equation}\label{eq:quadratic_value_ansatz}
\begin{aligned}
V_n^{i,0}(t,z^i)
&=
\frac{1}{2}(z^i)^\top H_n(t)z^i
+(r_n^{i,0}(t))^\top z^i+c_n^{i,0}(t),\\
V_n^{i,k}(t,z^i)
&=
\frac{1}{2}(z^i)^\top H_n(t)z^i
+(r_n^k(t))^\top z^i+c_n^k(t),
\qquad k=1,\ldots,m.
\end{aligned}
\end{equation}
In this ansatz, \(H_n\) is common across traders and post-announcement
scenarios and is also the pre-announcement quadratic coefficient. The pairs
\((r_n^k,c_n^k)\) are scenario-specific after the announcement, whereas
\((r_n^{i,0},c_n^{i,0})\) are trader-specific before it.

To simplify notation, we first write the dynamics of \(Z_n^i\) in the
vectorized form
\begin{equation}\label{eq:fixed_interval_reduced_dynamics}
\dot Z_n^i=A_zZ_n^i+B_nu_n^i+B_n^{-}U_n^{-i}+B_gg,
\end{equation}
where
\[
A_z:=\begin{bmatrix}0&0&0\\0&0&0\\0&0&-R\end{bmatrix},
\quad
B_n:=\begin{bmatrix}\frac{n-1}{n}\operatorname{Id}_d\\
\frac{1}{n}\operatorname{Id}_d\\ \varpi_n\Gamma\end{bmatrix},
\quad
B_n^{-}:=\begin{bmatrix}-\frac{1}{n}\operatorname{Id}_d\\
\frac{1}{n}\operatorname{Id}_d\\ \varpi_n\Gamma\end{bmatrix},
\quad
B_g:=\begin{bmatrix}0\\0\\\Gamma\end{bmatrix}.
\]
With
\(E_x:=[\operatorname{Id}_d\ \operatorname{Id}_d\ 0]\) and
\(E_\iota:=[0\ 0\ \operatorname{Id}_d]\), one has
\(X_n^i=E_xZ_n^i\) and \(I=E_\iota Z_n^i\). Consequently, the inventory
penalty in the running cost \(\mathcal L\) from
\eqref{eq:reduced_running_cost} and the terminal inventory penalty in
\eqref{eq:reformulated_objective} are represented in the reduced state by
\(Q_z\) and \(H_T\), respectively, where
\begin{equation}\label{eq:fixed_interval_final_boundary}
Q_z:=E_x^\top QE_x,
\qquad
H_T:=E_x^\top Q_TE_x.
\end{equation}

Additionally, we introduce some notation required to compactly state the
ODE system satisfied by the coefficients in \eqref{eq:quadratic_value_ansatz}.
Set
\begin{equation}\label{eq:fixed_interval_a_d_def}
\vartheta_n:=\chi_u\varpi_n,
\qquad
a_n:=2-\vartheta_n,
\qquad
d_n:=2+\vartheta_n(n-1).
\end{equation}
Assumption~\ref{ass:standing_model} gives \(0\leq\vartheta_n\leq1\), so
\(1\leq a_n\leq2\) and \(d_n\geq2\). Define
\begin{equation}\label{eq:fixed_interval_D_def}
\begin{aligned}
D_n&:=\operatorname{diag}\left(
\tfrac{1}{a_n}\operatorname{Id}_d,
\tfrac{1}{d_n}\operatorname{Id}_d,
\tfrac{1}{d_n}\operatorname{Id}_d\right),
&
F_n&:=\operatorname{diag}\left(
\tfrac{1}{a_n}\operatorname{Id}_d,
-\tfrac{n-1}{d_n}\operatorname{Id}_d,
-\tfrac{n-1}{d_n}\operatorname{Id}_d\right).
\end{aligned}
\end{equation}
For \(H\in\mathbb S^{3d}\), define
\begin{equation}\label{eq:fixed_interval_matrix_maps_def}
\Theta_n(H):=E_\iota+B_n^\top H,
\qquad
M_n(H):=A_z+B_n^{-}\Lambda^{-1}\Theta_n(H)F_n,
\end{equation}
\begin{subequations}\label{eq:pre_ann_matrix_maps_def}
\begin{align}
\mathsf A_n^{\mathrm c}(H)
&:=M_n(H)^\top-\tfrac{1}{d_n}
\bigl((n-1)HB_n^{-}+2D_n^\top\Theta_n(H)^\top\bigr)
\Lambda^{-1}B_n^\top,\\
\mathsf A_n^{\mathrm d}(H)
&:=M_n(H)^\top+\tfrac{1}{a_n}
\bigl(HB_n^{-}-2D_n^\top\Theta_n(H)^\top\bigr)
\Lambda^{-1}B_n^\top.
\end{align}
\end{subequations}
For \(r\in\mathbb R^{3d}\) and \(g\in\mathbb R^d\), set
\begin{equation}\label{eq:post_ann_q_def}
q_n(r,g):=B_n^\top r+\chi_g\Lambda g.
\end{equation}


We now state the ODE system for the coefficients in
\eqref{eq:quadratic_value_ansatz}.
Theorem~\ref{thm:full_horizon_closed_loop_equilibrium} below then verifies
that these indeed give the value functions and also gives the
corresponding equilibrium feedbacks. Time arguments are suppressed where no
ambiguity arises.

\begin{itemize}
\item \textbf{Quadratic coefficient.}
On \([0,T]\), \(H_n\) solves the non-standard Riccati ODE
\begin{equation}\label{eq:fixed_interval_riccati_P}
\begin{aligned}
\dot H_n={}&-Q_z-M_n(H_n)^\top H_n-H_nM_n(H_n)\\[-2pt]
&+2D_n^\top\Theta_n(H_n)^\top\Lambda^{-1}\Theta_n(H_n)D_n,
\\[-2pt]
H_n(T)={}&H_T.
\end{aligned}
\end{equation}

\item \textbf{Linear coefficients.}
\begin{subequations}\label{eq:full_horizon_linear_coefficients}
On \([T_{\mathrm{ann}},T]\), \(k=1,\ldots,m\), \(r_n^k\) solves the linear ODE
\begin{equation}\label{eq:post_ann_common_r}
\begin{aligned}
\dot r_n^k={}&-\mathsf A_n^{\mathrm c}(H_n)r_n^k-H_nB_gg^k
\\[-2pt]
&+\tfrac{\chi_g}{d_n}
\bigl((n-1)H_nB_n^{-}+2D_n^\top\Theta_n(H_n)^\top\bigr)g^k,
\\[-2pt]
r_n^k(T)={}&0,
\end{aligned}
\end{equation}
On \([0,T_{\mathrm{ann}}]\), the pre-announcement affine coefficients solve
the coupled system
\begin{equation}\label{eq:pre_ann_coupled_r_ode}
\begin{aligned}
\dot r_n^{i,0}
={}&-\mathsf A_n^{\mathrm d}(H_n)r_n^{i,0}
+\tfrac{1}{n}
\bigl(\mathsf A_n^{\mathrm d}(H_n)-\mathsf A_n^{\mathrm c}(H_n)\bigr)
\sum_{j=1}^n r_n^{j,0},
\\[-2pt]
r_n^{i,0}(T_{\mathrm{ann}})
={}&\sum_{k=1}^m\pi_n^{i,k}r_n^k(T_{\mathrm{ann}}),
\qquad i=1,\ldots,n.
\end{aligned}
\end{equation}
Because \eqref{eq:pre_ann_coupled_r_ode} is coupled through
\(\sum_jr_n^{j,0}\), solving it directly appears to require the full belief
profile \((\pi_n^1,\ldots,\pi_n^n)\), which is not available to any trader.
To eliminate this apparent full-profile dependence, set
\begin{equation}\label{eq:pre_ann_affine_components}
r_n^0=\bar r_n^0:=\tfrac{1}{n}\sum_{j=1}^n r_n^{j,0},
\qquad
\widetilde r_n^{i,0}:=r_n^{i,0}-r_n^0
\end{equation}
The bar emphasizes that \(r_n^0\) is the population-average affine
coefficient. We use the unbarred notation in the sequel to keep the notation
uniform across pre- and post-announcement indices.
Averaging \eqref{eq:pre_ann_coupled_r_ode} and using
\(n^{-1}\sum_i\pi_n^i=\bar\pi_n\) gives
\begin{equation}\label{eq:pre_ann_common_r_ode}
\dot r_n^0
=-\mathsf A_n^{\mathrm c}(H_n)r_n^0,
\qquad
r_n^0(T_{\mathrm{ann}})
=\sum_{k=1}^m\bar\pi_n^kr_n^k(T_{\mathrm{ann}}).
\end{equation}
Also, subtracting the average equation from trader \(i\)'s equation gives
\begin{equation}\label{eq:pre_ann_deviation_r_ode}
\dot{\widetilde r}_n^{i,0}
=-\mathsf A_n^{\mathrm d}(H_n)\widetilde r_n^{i,0},
\qquad
\widetilde r_n^{i,0}(T_{\mathrm{ann}})
=\sum_{k=1}^m(\pi_n^{i,k}-\bar\pi_n^k)r_n^k(T_{\mathrm{ann}}).
\end{equation}
The ODE for \(r_n^0\) depends only on \(\bar\pi_n\), whereas the ODE
for \(\widetilde r_n^{i,0}\) depends only on
\(\pi_n^i-\bar\pi_n\). Hence trader \(i\) can compute
\(r_n^{i,0}\) from \((\bar\pi_n,\pi_n^i)\).
\end{subequations}

\item \textbf{Constant coefficients.}
\begin{subequations}\label{eq:full_horizon_constant_coefficients}
On \([T_{\mathrm{ann}},T]\), \(k=1,\ldots,m\), \(c_n^k\) solves
\begin{equation}\label{eq:post_ann_common_alpha}
\begin{aligned}
\dot c_n^k={}&-\left(
B_gg^k-\tfrac{n-1}{d_n}B_n^{-}\Lambda^{-1}q_n(r_n^k,g^k)
\right)^\top r_n^k
\\[-2pt]
&+\tfrac{1}{d_n^2}q_n(r_n^k,g^k)^\top
\Lambda^{-1}q_n(r_n^k,g^k),
\\[-2pt]
c_n^k(T)={}&0,
\end{aligned}
\end{equation}
On \([0,T_{\mathrm{ann}}]\), \(i=1,\ldots,n\), \(c_n^{i,0}\) solves
\begin{equation}\label{eq:pre_ann_constant_ode}
\begin{aligned}
\dot c_n^{i,0}
={}&-\left[
B_n^{-}\Lambda^{-1}
\left(
\tfrac{1}{a_n}B_n^\top\widetilde r_n^{i,0}
-\tfrac{n-1}{d_n}B_n^\top r_n^0
\right)
\right]^\top r_n^{i,0}
\\[-2pt]
&+\left(
\tfrac{1}{a_n}B_n^\top\widetilde r_n^{i,0}
+\tfrac{1}{d_n}B_n^\top r_n^0
\right)^\top
\Lambda^{-1}
\left(
\tfrac{1}{a_n}B_n^\top\widetilde r_n^{i,0}
+\tfrac{1}{d_n}B_n^\top r_n^0
\right),
\\[-2pt]
c_n^{i,0}(T_{\mathrm{ann}})
={}&\sum_{k=1}^m\pi_n^{i,k}c_n^k(T_{\mathrm{ann}}).
\end{aligned}
\end{equation}
\end{subequations}
\end{itemize}

\begin{theorem}[Explicit Quadratic Equilibrium]
\label{thm:full_horizon_closed_loop_equilibrium}
Suppose that the Riccati equation
\eqref{eq:fixed_interval_riccati_P} admits a solution
\(H_n\in C^1([0,T];\mathbb S^{3d})\), and construct the remaining
coefficients from \eqref{eq:full_horizon_linear_coefficients} and
\eqref{eq:full_horizon_constant_coefficients}. Define the pre-announcement
feedback on \([0,T_{\mathrm{ann}}]\) and, for \(k=1,\ldots,m\), the
scenario-\(k\) feedback on \([T_{\mathrm{ann}},T]\) by
\begin{equation}\label{eq:full_horizon_markov_feedback}
\begin{aligned}
\phi_n^{i,0}(t,z^i;\bar\pi_n,\pi_n^i)
={}&-\Lambda^{-1}\Bigl[
\Theta_n(H_n(t))D_nz^i
+\tfrac{1}{d_n}B_n^\top r_n^0(t)
+\tfrac{1}{a_n}B_n^\top\widetilde r_n^{i,0}(t)
\Bigr],\\[4pt]
\phi_n^{i,k}(t,z^i;\bar\pi_n,\pi_n^i)
={}&-\Lambda^{-1}\left[
\Theta_n(H_n(t))D_nz^i
+\tfrac{1}{d_n}
B_n^\top r_n^k(t)+\tfrac{1}{d_n}\chi_g\Lambda g^k(t)
\right].
\end{aligned}
\end{equation}
The value functions in \eqref{eq:quadratic_value_ansatz} and the feedbacks in
\eqref{eq:full_horizon_markov_feedback} satisfy
\eqref{eq:stitched_hjb_system}. The Markov feedback family \(\phi_n\) defined
by \eqref{eq:full_horizon_markov_feedback} is admissible in every continuation
game. Consequently, \(\phi_n\) is a subgame-perfect Nash equilibrium.
\end{theorem}

The proof is given in
Appendix~\ref{subsec:proof_full_horizon_closed_loop_equilibrium}.
Note that the Riccati equation has dimension independent of the number of traders. This
scalable structure is related to the decomposition of discrete-time LQ
deep-structured games in
\citet{arabneydi2020explicitsequentialequilibrialq}.

\begin{remark}[State Decomposition]\label{rem:full_procedure_decomposition}
For the equilibrium trajectories, let \(X_n^{i,0}\) and \(\bar X_n^0\)
denote the pre-announcement inventories and let \(X_n^{i,k}\) and
\(\bar X_n^k\) denote their continuations in scenario \(k\). Define
\(\delta_n^{i,0}:=X_n^{i,0}-\bar X_n^0\) before the announcement and
\(\delta_n^{i,k}:=X_n^{i,k}-\bar X_n^k\) in scenario \(k\) after the
announcement. Continuity of
the individual and average inventories gives
\(\delta_n^{i,k}(T_{\mathrm{ann}})=\delta_n^{i,0}(T_{\mathrm{ann}})\).
For the generic coordinates \((\delta,\bar x,\iota)\), write the block columns
of the feedback matrix as
\begin{equation}\label{eq:fixed_interval_theta_partition}
\Theta_n(H_n(t))
=
\begin{bmatrix}
[\Theta_n(H_n(t))]_{\delta}&
[\Theta_n(H_n(t))]_{\bar x}&
[\Theta_n(H_n(t))]_{\iota}
\end{bmatrix}.
\end{equation}
In scenario \(k\), define the common feedback signal
\begin{equation}\label{eq:post_ann_common_feedback_signal}
\mathcal K_n^k(t)
:=
[\Theta_n(H_n(t))]_{\bar x}\bar X_n^k(t)
+[\Theta_n(H_n(t))]_{\iota}I_n^k(t)
+B_n^\top r_n^k(t)+\chi_g\Lambda g^k(t).
\end{equation}
Using \eqref{eq:full_horizon_markov_feedback}, the individual feedback and
aggregate opportunist flow in scenario \(k\) are
\begin{align}
u_n^{i,k}(t)
&=
-\Lambda^{-1}
\left(
\frac1{a_n}[\Theta_n(H_n(t))]_{\delta}\delta_n^{i,k}(t)
+\frac1{d_n}\mathcal K_n^k(t)
\right),
\label{eq:post_ann_feedback_decomposition}\\
U_n^{k}(t)
&=
-\frac{n}{d_n}\Lambda^{-1}\mathcal K_n^k(t).
\label{eq:post_ann_total_flow_decomposition}
\end{align}
Therefore the common state \((\bar X_n^k,I_n^k)\) satisfies
\begin{align}
\begin{bmatrix}
\dot{\bar X}_n^k(t)\\
\dot I_n^k(t)
\end{bmatrix}
={}&
\begin{bmatrix}
-\frac1{d_n}\Lambda^{-1}[\Theta_n(H_n(t))]_{\bar x}
&
-\frac1{d_n}\Lambda^{-1}[\Theta_n(H_n(t))]_{\iota}
\\[5pt]
-\frac{n\varpi_n}{d_n}\Gamma\Lambda^{-1}[\Theta_n(H_n(t))]_{\bar x}
&
-R-\frac{n\varpi_n}{d_n}\Gamma\Lambda^{-1}[\Theta_n(H_n(t))]_{\iota}
\end{bmatrix}
\begin{bmatrix}
\bar X_n^k(t)\\
I_n^k(t)
\end{bmatrix}
\notag\\
&-
\frac1{d_n}
\begin{bmatrix}
\Lambda^{-1}\\
n\varpi_n\Gamma\Lambda^{-1}
\end{bmatrix}
B_n^\top r_n^k(t)
+
\begin{bmatrix}
-\frac{\chi_g}{d_n}\operatorname{Id}_d\\[4pt]
\left(1-\frac{n\varpi_n\chi_g}{d_n}\right)\Gamma
\end{bmatrix}
g^k(t).
\label{eq:post_ann_common_subsystem}
\end{align}
Trader \(i\)'s post-announcement inventory deviation satisfies
\begin{equation}\label{eq:post_ann_delta_dynamics}
\dot\delta_n^{i,k}(t)
=
-\frac1{a_n}\Lambda^{-1}
[\Theta_n(H_n(t))]_{\delta}\delta_n^{i,k}(t),
\qquad
\delta_n^{i,k}(T_{\mathrm{ann}})
=\delta_n^{i,0}(T_{\mathrm{ann}}).
\end{equation}
Since neither the ODE nor its initial condition depends on \(k\), uniqueness
implies that \(\delta_n^{i,k}\) is identical across scenarios. The scenario
only affects the common state through \(r_n^k\) and \(g^k\).

Before the announcement the individual and aggregate
controls are
\begin{align*}
u_n^{i,0}(t)
&=
-\Lambda^{-1}
\Biggl(
\frac1{a_n}
\Bigl(
[\Theta_n(H_n(t))]_{\delta}\delta_n^{i,0}(t)
+B_n^\top\widetilde r_n^{i,0}(t)
\Bigr)
\\
&\qquad
+\frac1{d_n}
\Bigl(
[\Theta_n(H_n(t))]_{\bar x}\bar X_n^0(t)
+[\Theta_n(H_n(t))]_{\iota}I_n^0(t)
+B_n^\top r_n^0(t)
\Bigr)
\Biggr),\\
U_n^{0}(t)
&=
-\frac{n}{d_n}\Lambda^{-1}
\Bigl(
[\Theta_n(H_n(t))]_{\bar x}\bar X_n^0(t)
+[\Theta_n(H_n(t))]_{\iota}I_n^0(t)
+B_n^\top r_n^0(t)
\Bigr).
\end{align*}
The common state \((\bar X_n^0,I_n^0)\)
satisfies
\begin{align}
\frac{d}{dt}
\begin{bmatrix}
\bar X_n^0(t)\\
I_n^0(t)
\end{bmatrix}
={}&
\begin{bmatrix}
-\frac1{d_n}\Lambda^{-1}[\Theta_n(H_n(t))]_{\bar x}
&
-\frac1{d_n}\Lambda^{-1}[\Theta_n(H_n(t))]_{\iota}
\\[5pt]
-\frac{n\varpi_n}{d_n}\Gamma\Lambda^{-1}[\Theta_n(H_n(t))]_{\bar x}
&
-R-\frac{n\varpi_n}{d_n}\Gamma\Lambda^{-1}[\Theta_n(H_n(t))]_{\iota}
\end{bmatrix}
\begin{bmatrix}
\bar X_n^0(t)\\
I_n^0(t)
\end{bmatrix}
\notag\\
&-
\frac1{d_n}
\begin{bmatrix}
\Lambda^{-1}\\
n\varpi_n\Gamma\Lambda^{-1}
\end{bmatrix}
B_n^\top r_n^0(t).
\label{eq:pre_ann_common_subsystem}
\end{align}
The pre-announcement deviation satisfies
\begin{equation}\label{eq:pre_ann_delta_dynamics}
\dot\delta_n^{i,0}(t)
=
-\frac1{a_n}\Lambda^{-1}
\left(
[\Theta_n(H_n(t))]_{\delta}\delta_n^{i,0}(t)
+B_n^\top\widetilde r_n^{i,0}(t)
\right),
\qquad
\delta_n^{i,0}(0)
=
x_0^i-\frac1n\sum_{j=1}^n x_0^j.
\end{equation}
The common-state and individual-deviation systems decouple both before and
after the announcement, as shown in
\eqref{eq:post_ann_common_subsystem}--\eqref{eq:post_ann_delta_dynamics} and
\eqref{eq:pre_ann_common_subsystem}--\eqref{eq:pre_ann_delta_dynamics}.
Unlike \eqref{eq:post_ann_delta_dynamics},
\eqref{eq:pre_ann_delta_dynamics} is forced by the centered affine coefficient
\(\widetilde r_n^{i,0}\) through 
\(B_n^\top\widetilde r_n^{i,0}\). Its terminal value encodes trader \(i\)'s
belief deviation, and its dynamics are given by
\eqref{eq:pre_ann_deviation_r_ode}.
\end{remark}

\subsection{Global Solvability}\label{sec:solvability_structure}

Theorem~\ref{thm:full_horizon_closed_loop_equilibrium} constructs a
subgame-perfect closed-loop Nash equilibrium conditional on the Riccati
equation admitting a solution on \([0,T]\). We now show that
Assumption~\ref{ass:standing_model} guarantees this full-horizon solvability a
priori.

\begin{theorem}[Uniform Global Solvability]
\label{thm:global_riccati_solvability}
Under Assumption~\ref{ass:standing_model}, for every \(n\ge2\), let
\(H_n\) be the maximal backward solution of
\eqref{eq:fixed_interval_riccati_P}, with the corresponding
\(n\)-player coefficients, and with terminal matrix \(H_T\) from
\eqref{eq:fixed_interval_final_boundary}.
Then \(H_n\) exists uniquely on the whole interval \([0,T]\).
More precisely, for \(t\in[0,T]\), define
\[
G_Q(t):=Q_T+(T-t)Q,
\qquad
G_R(t)
:=
\int_t^T
e^{-R^\top(s-t)}\Lambda^{-1}e^{-R(s-t)}\,ds,
\]
and set
\begin{align}
\overline H(t)
&:=
\begin{bmatrix}
G_Q(t)&G_Q(t)&0\\
G_Q(t)&G_Q(t)&0\\
0&0&0
\end{bmatrix},
\label{eq:global_riccati_upper_bound_matrix}\\
\underline H(t)
&:=
\left(
2\operatorname{tr}(G_Q(t))
+\frac12\operatorname{tr}(G_R(t))
\right)\operatorname{Id}_{3d}.
\label{eq:global_riccati_lower_bound_matrix}
\end{align}
Then, uniformly over \(n\ge2\),
\begin{equation}
\label{eq:global_riccati_two_sided_bound}
-\underline H(t)
\preceq
H_n(t)
\preceq
\overline H(t),
\qquad
0\le t\le T.
\end{equation}
The post-announcement systems for \(r_n^k\) and \(c_n^k\), the
pre-announcement systems for \(r_n^0\) and
\(\widetilde r_n^{i,0}\), and the scalar system for \(c_n^{i,0}\) also admit
unique finite solutions. Hence the construction in
Theorem~\ref{thm:full_horizon_closed_loop_equilibrium} is well defined.
\end{theorem}

The proof is given in
Appendix~\ref{subsec:proof_global_riccati_solvability}.

\begin{remark}[Impacted Terminal Marks and Solvability]
\label{rem:impacted_terminal_mark_solvability}
The global-solvability proof relies on marking terminal inventory to the
fundamental price. Marking it instead to \(S(T)=S^0(T)+I(T)\) adds
\(\iota^\top x_n^i-I(T)^\top X_n^i(T)\) to the continuation objective and changes
the terminal matrix in \eqref{eq:fixed_interval_final_boundary} to
\[
H_T^{\mathrm{imp}}
=
\begin{bmatrix}
Q_T&Q_T&-\operatorname{Id}_d\\
Q_T&Q_T&-\operatorname{Id}_d\\
-\operatorname{Id}_d&-\operatorname{Id}_d&0
\end{bmatrix}.
\]
If the Riccati equation with terminal condition
\(H_n(T)=H_T^{\mathrm{imp}}\) has a solution on \([0,T]\), then the
construction in Theorem~\ref{thm:full_horizon_closed_loop_equilibrium} still yields
a subgame-perfect closed-loop Nash equilibrium for the game where wealth is marked
to terminal impacted price. What fails is the a priori
argument for global existence. Indeed, one can construct parameter values for which the modified
Riccati solution blows up in finite time. We omit the construction because it is
not needed for the subsequent analysis.

Both \citet{NeumanVoss2023} and
\citet{campbell2026optimalexecutionntraders} mark terminal inventory at the
fundamental price to retain convexity of each player's optimization problem,
making their first-order conditions sufficient. We use the same valuation
convention, but for a different reason: it provides the a priori bounds that
prevent Riccati blow-up and guarantee full-horizon existence of the value
functions and the closed-loop Nash equilibrium.
\end{remark}

\section{Mean-Field Limit}\label{sec:mean_field_limit}

The finite-player closed-loop strategies constructed in
Section~\ref{sec:solution} are attractive because they respond to the realized
state. They can, however, be difficult to implement when an opportunist knows
only that the number of relevant players is large, rather than the exact value
of \(n\), and cannot observe the current average inventory \(\bar X_n\); see
Remark~\ref{rem:average_inventory_proxy}. A mean-field strategy avoids both
requirements. We assume that every opportunist starts with zero inventory,
which is natural when they participate only to exploit the rebalance. The equilibrium mean paths can then be computed ex
ante for every announcement scenario and need not be observed during execution.
The mean-field feedback therefore depends on the trader's own inventory and
belief and on the announced scenario, but not on the current population-average
inventory.

We define a mean-field game motivated by the \(n\to\infty\) limit of the
finite-player model under the mean-field scaling \(\varpi_n=\varpi/n\), construct
its Markov equilibrium, and then prove that the finite-player equilibria
converge to this equilibrium. Such convergence is not automatic:
\citet{Lacker2020ClosedLoop} shows that
subsequential limits of closed-loop Nash equilibria need only be weak
mean-field equilibria and may contain endogenous common randomness, so
convergence to a strong mean-field equilibrium is not automatic. In our LQ
model, if
\(\bar\pi_n\to\bar\pi\), the finite-player equilibria constructed in
Section~\ref{sec:full_procedure} converge uniformly to the Markov mean-field
equilibrium of Theorem~\ref{thm:mf_solution}, and the lifted mean-field
feedbacks form approximate Nash equilibria by
Theorem~\ref{thm:mf_approximate_nash}.

\subsection{Definition of the Mean-Field Game}

Let \(\Pi\) be a probability distribution on \(\Delta_m\). A representative
player of type \(\pi\in\Delta_m\) evaluates the announcement under the
subjective probability
\[
\mathbb P^\pi(\Xi=\xi^k)=\pi^k,
\qquad k=1,\ldots,m.
\]
The \emph{population-average belief} is
\[
\bar\pi:=\int_{\Delta_m}\pi\,\Pi(d\pi).
\]
It is the population analogue of the finite-population average belief
\(\bar\pi_n\) in \eqref{eq:finite_population_average_prior}.

As in the finite-player game, the objective and state dynamics can be written
without explicit dependence on the fundamental price. We thus seek Markov feedbacks
that ignore \(S^0\), with one feedback before the announcement and one for each
post-announcement scenario. Except for the randomness coming from the
announcement, the game is fully deterministic.

\begin{definition}[Mean-Field Equilibrium]\label{def:mf_equilibrium}
A candidate mean flow is a collection
\[
\bar u^0\in L^2([0,T_{\mathrm{ann}}],\mathbb R^d),
\qquad
\bar u^k\in L^2([T_{\mathrm{ann}},T],\mathbb R^d),
\quad k=1,\ldots,m.
\]
Its impact components satisfy
\begin{equation}\label{eq:mf_component_impact_dynamics}
\begin{aligned}
\dot I^0(t)
&=-RI^0(t)+\varpi\Gamma\bar u^0(t),
& I^0(0)&=0,
\\
\dot I^k(t)
&=-RI^k(t)+\Gamma\bigl(\varpi\bar u^k(t)+g^k(t)\bigr),
& I^k(T_{\mathrm{ann}})&=I^0(T_{\mathrm{ann}}),
\quad k=1,\ldots,m.
\end{aligned}
\end{equation}
For a type \(\pi\), an admissible Markov policy \(\phi^\pi\) is a collection of
Borel maps
\[
\phi^{\pi,0}:[0,T_{\mathrm{ann}}]\times\mathbb R^d\longrightarrow\mathbb R^d,
\qquad
\phi^{\pi,k}:[T_{\mathrm{ann}},T]\times\mathbb R^d\longrightarrow\mathbb R^d,
\quad k=1,\ldots,m,
\]
such that the rates
\[
\begin{aligned}
u^{\pi,0}(t)&:=\phi^{\pi,0}(t,X^{\pi,0}(t)),\\
u^{\pi,k}(t)&:=\phi^{\pi,k}(t,X^{\pi,k}(t)),
\quad k=1,\ldots,m,
\end{aligned}
\]
generate unique solutions of
\(\dot X^{\pi,0}=u^{\pi,0}\) and
\(\dot X^{\pi,k}=u^{\pi,k}\), \(k=1,\ldots,m\), satisfying
\[
X^{\pi,0}(0)=0,
\qquad
X^{\pi,k}(T_{\mathrm{ann}})=X^{\pi,0}(T_{\mathrm{ann}}),
\quad k=1,\ldots,m,
\]
and are square-integrable on their respective intervals. The type-\(\pi\)
objective is
\begin{equation}\label{eq:mf_type_objective}
\begin{aligned}
J^\pi(\phi^\pi;\bar u)
={}&
\int_0^{T_{\mathrm{ann}}}
\mathcal L\bigl(
X^{\pi,0}(t),u^{\pi,0}(t);
I^0(t),\varpi\bar u^0(t),0
\bigr)\,dt
\\
&+
\sum_{k=1}^m\pi^k
\int_{T_{\mathrm{ann}}}^T
\mathcal L\bigl(
X^{\pi,k}(t),u^{\pi,k}(t);
I^k(t),\varpi\bar u^k(t),g^k(t)
\bigr)\,dt
\\
&+
\frac{1}{2}\sum_{k=1}^m\pi^k
(X^{\pi,k}(T))^\top Q_T X^{\pi,k}(T).
\end{aligned}
\end{equation}
A family \(\phi=(\phi^{\pi})_{\pi\in\Delta_m}\) of admissible Markov
policies is a mean-field equilibrium if the following two conditions hold.
First, the componentwise averages
\begin{equation}\label{eq:mf_average_consistency}
\begin{aligned}
\bar u^0(t)
&=
\int_{\Delta_m}
\phi^{\pi,0}(t,X^{\pi,0}(t))\,\Pi(d\pi),
\\
\bar u^k(t)
&=
\int_{\Delta_m}
\phi^{\pi,k}(t,X^{\pi,k}(t))\,\Pi(d\pi),
\quad k=1,\ldots,m,
\end{aligned}
\end{equation}
are well defined, hold for a.e. time on their respective intervals, and form a
candidate mean flow. Second, for \(\Pi\)-a.e. \(\pi\), the policy
\(\phi^{\pi}\) is a best response to this mean flow:
\[
\phi^{\pi}
\in
\operatorname*{argmin}_{\psi^\pi}
J^\pi(\psi^\pi;\bar u),
\]
where the minimum is taken over admissible Markov policies and the mean flow
\(\bar u\) and its impact components are held fixed. The second condition also
requires, for every \(k=1,\ldots,m\), that \(\phi^{\pi,k}\) solve the
post-announcement control problem from every announcement-time inventory
\(x\in\mathbb R^d\), with indexer flow \(g^k\), mean flow \(\bar u^k\), and
impact component \(I^k\), regardless of whether \(\pi^k=0\).
\end{definition}

We impose \(X^{\pi,0}(0)=0\) for every type and \(I^0(0)=0\). Thus the
population heterogeneity represented by \(\Pi\) concerns beliefs rather than
initial inventories. This differs from the liquidation game of
\citet{NeumanVoss2023}, where agents have heterogeneous deterministic initial
positions and mean-field consistency is imposed through the limiting empirical
average of their trading rates.

Our model combines two features: heterogeneous beliefs and mean-field
consistency obtained by averaging over a distribution of belief types.
\citet{Casgrain2020} also allow agents to disagree about the dynamics of the
asset price and its latent factors, but divide the population into finitely
many positive-mass subpopulations, each sharing the same beliefs and objective.
The papers \citep{Lacker2019,Lacker2020} instead formulate
mean-field games with a general distribution of types, such as risk
tolerances, relative-performance preferences, and market parameters. Both
\citet{Casgrain2020} and \citet{Lacker2019,Lacker2020} retain stochastic
asset-price dynamics in their mean-field games. By contrast, \(S^0\) drops out
of our mean-field objective and state dynamics, so the common announcement
\(\Xi\) is the only remaining source of uncertainty and the dynamics are
deterministic conditional on its realization.


\subsection{Construction of the Mean-Field Equilibrium}

As in the finite-player game, we seek value functions for a type-\(\pi\)
player of the form
\begin{equation}\label{eq:mf_quadratic_value_ansatz}
\begin{aligned}
V^{\pi,0}(t,x)
&=\frac12x^\top P(t)x+(\ell^{\pi,0}(t))^\top x+c^{\pi,0}(t),
\\
V^k(t,x)
&=\frac12x^\top P(t)x+(\ell^k(t))^\top x+c^k(t),
\qquad k=1,\ldots,m.
\end{aligned}
\end{equation}
The quadratic coefficient is common across types, intervals, and scenarios.
The affine and constant coefficients depend on the scenario after the
announcement and on the type before the announcement.

To state the coefficient equations and the corresponding optimal controls,
define the type-average pre-announcement affine coefficient by
\[
\ell^0(t)=\bar\ell^0(t)
:=
\int_{\Delta_m}\ell^{\pi,0}(t)\,\Pi(d\pi).
\]
The bar emphasizes that this coefficient is a type average. We use the
unbarred notation in the sequel to keep the notation uniform across
pre- and post-announcement indices.
Let \(\bar X^0\) and \(I^0\) denote the equilibrium mean-inventory and
impact-state paths before the announcement, and let \(\bar X^k\) and \(I^k\)
denote their counterparts in post-announcement scenario \(k\). Set
\begin{equation}\label{eq:mf_augmented_common_state_paths}
Y^k(t):=(\bar X^k(t),I^k(t),\ell^k(t))^\top,
\qquad k=0,\ldots,m.
\end{equation}

We call \(Y^0,\ldots,Y^m\) the augmented mean-field common-state paths:
their \((\bar X,I)\) components are the mean-field analogue of the
finite-player common state, augmented by the affine coefficient \(\ell^k\).
These paths and the remaining equilibrium coefficients are determined as
follows.
Theorem~\ref{thm:mf_solution} below makes this construction rigorous and also
states the corresponding equilibrium feedbacks.

\begin{itemize}
\item \textbf{Quadratic Coefficient.}

On \([0,T]\), \(P\) solves
\begin{equation}\label{eq:mf_riccati}
\dot P(t)
=
- Q+\frac12P(t)\Lambda^{-1}P(t),
\qquad
P(T)=Q_T.
\end{equation}
Since \(\Lambda\succ0\) and \(Q,Q_T\succeq0\), standard finite-horizon LQR
theory gives a unique continuous symmetric nonnegative solution on
\([0,T]\).

\item \textbf{Augmented Mean-Field Common-State System.}

The augmented mean-field common-state paths \(Y^0,\ldots,Y^m\) satisfy
\begin{subequations}\label{eq:mf_consistency_system}
\begin{equation}\label{eq:mf_y_system}
\begin{aligned}
\dot Y^0(t)&=\mathcal A(t)Y^0(t),
&&t\in[0,T_{\mathrm{ann}}],
\\
\dot Y^k(t)&=\mathcal A(t)Y^k(t)+b^k(t),
&&t\in[T_{\mathrm{ann}},T],
\quad k=1,\ldots,m.
\end{aligned}
\end{equation}
with
\begin{equation}\label{eq:mf_pre_consistency_boundary}
\bar X^0(0)=0,
\qquad
I^0(0)=0,
\qquad
\ell^0(T_{\mathrm{ann}})
=\sum_{k=1}^m\bar\pi^k\ell^k(T_{\mathrm{ann}}),
\end{equation}
and, for \(k=1,\ldots,m\),
\begin{equation}\label{eq:mf_post_consistency_boundary}
\bar X^k(T_{\mathrm{ann}})=\bar X^0(T_{\mathrm{ann}}),
\qquad
I^k(T_{\mathrm{ann}})=I^0(T_{\mathrm{ann}}),
\qquad
\ell^k(T)=0.
\end{equation}
\end{subequations}
Here
\begin{equation}\label{eq:mf_A_matrix}
\mathcal A(t)
=
\frac{1}{2+\varpi\chi_u}
\begin{pmatrix}
-\Lambda^{-1}P(t)&-\Lambda^{-1}&-\Lambda^{-1}\\
-\varpi\Gamma\Lambda^{-1}P(t)&
-(2+\varpi\chi_u)R-\varpi\Gamma\Lambda^{-1}&
-\varpi\Gamma\Lambda^{-1}\\
-\frac{\varpi\chi_u}{2}P(t)\Lambda^{-1}P(t)&
P(t)\Lambda^{-1}&P(t)\Lambda^{-1}
\end{pmatrix},
\end{equation}
and \(b^0(t):=0\), while, for \(k=1,\ldots,m\),
\begin{equation}\label{eq:mf_b_vector}
b^k(t)
=
\frac{1}{2+\varpi\chi_u}
\begin{pmatrix}
-\chi_g g^k(t)\\
(2+\varpi\chi_u-\varpi\chi_g)\Gamma g^k(t)\\
\chi_g P(t)g^k(t)
\end{pmatrix}.
\end{equation}
Note that the conditions for \(\bar X^0\) and \(I^0\) are imposed at \(t=0\), whereas
the conditions for \(\ell^k\), \(k=1,\ldots,m\), are imposed at \(t=T\). The conditions at
\(T_{\mathrm{ann}}\) match the pre- and post-announcement mean inventories
and impact states and link \(\ell^0\) to the post-announcement affine
coefficients. Thus, \eqref{eq:mf_consistency_system} is a coupled
boundary-value problem rather than an initial-value problem, and its
solvability is not immediate.

The equilibrium mean flows are
\begin{subequations}\label{eq:mf_mu_closed_form}
\begin{align}
\bar u^0(t)
&=
-\frac{1}{2+\varpi\chi_u}\Lambda^{-1}
\left(P(t)\bar X^0(t)+I^0(t)+\ell^0(t)\right),
\label{eq:mf_pre_mu_closed_form}\\
\bar u^k(t)
&=
-\frac{1}{2+\varpi\chi_u}\Lambda^{-1}
\left(P(t)\bar X^k(t)+I^k(t)+\ell^k(t)+\Lambda\chi_g g^k(t)\right).
\label{eq:mf_post_mu_closed_form}
\end{align}
\end{subequations}

\item \textbf{Type-Specific Affine Coefficient.}

For \(\pi\in\Delta_m\), \(\ell^{\pi,0}\) solves on
\([0,T_{\mathrm{ann}}]\)
\begin{equation}\label{eq:mf_pre_ell}
\begin{aligned}
\dot\ell^{\pi,0}(t)
={}&
\frac12P(t)\Lambda^{-1}
\left(I^0(t)+\Lambda\varpi\chi_u\bar u^0(t)+\ell^{\pi,0}(t)\right),
\\
\ell^{\pi,0}(T_{\mathrm{ann}})
={}&\sum_{k=1}^m\pi^k\ell^k(T_{\mathrm{ann}}).
\end{aligned}
\end{equation}

\item \textbf{Constant Coefficients.}

For \(k=1,\ldots,m\), \(c^k\) solves on
\([T_{\mathrm{ann}},T]\)
\begin{equation}\label{eq:mf_post_alpha}
\begin{aligned}
\dot c^k(t)
={}&
\frac14
\left(I^k(t)+\Lambda(\varpi\chi_u\bar u^k(t)+\chi_g g^k(t))+\ell^k(t)\right)^\top
\\
&\qquad
\Lambda^{-1}
\left(I^k(t)+\Lambda(\varpi\chi_u\bar u^k(t)+\chi_g g^k(t))+\ell^k(t)\right),
\\
c^k(T)={}&0.
\end{aligned}
\end{equation}
For \(\pi\in\Delta_m\), \(c^{\pi,0}\) solves on
\([0,T_{\mathrm{ann}}]\)
\begin{equation}\label{eq:mf_pre_alpha}
\begin{aligned}
\dot c^{\pi,0}(t)
={}&
\frac14
\left(I^0(t)+\Lambda\varpi\chi_u\bar u^0(t)+\ell^{\pi,0}(t)\right)^\top
\\
&\qquad
\Lambda^{-1}
\left(I^0(t)+\Lambda\varpi\chi_u\bar u^0(t)+\ell^{\pi,0}(t)\right),
\\
c^{\pi,0}(T_{\mathrm{ann}})
={}&\sum_{k=1}^m\pi^kc^k(T_{\mathrm{ann}}).
\end{aligned}
\end{equation}
\end{itemize}

\begin{theorem}[Explicit Mean-Field Equilibrium]\label{thm:mf_solution}
Under Assumption~\ref{ass:standing_model}, let \(P\) be the solution of
\eqref{eq:mf_riccati}. Suppose that the augmented mean-field common-state
boundary-value problem
\eqref{eq:mf_consistency_system} admits a solution, and fix one such solution
\((Y^0,Y^1,\ldots,Y^m)\). Let \(\ell^{\pi,0}\), \(c^k\), and
\(c^{\pi,0}\) be determined by \eqref{eq:mf_pre_ell},
\eqref{eq:mf_post_alpha}, and \eqref{eq:mf_pre_alpha}.
Define \(\phi^{\pi,0}\) on \([0,T_{\mathrm{ann}}]\) and
\(\phi^{k}\) on \([T_{\mathrm{ann}},T]\) by
\begin{subequations}\label{eq:mf_equilibrium_feedbacks}
\begin{align}
\phi^{\pi,0}(t,x)
&=
-\frac12\Lambda^{-1}
\left(P(t)x+\ell^{\pi,0}(t)+I^0(t)+\Lambda\varpi\chi_u\bar u^0(t)\right),
\label{eq:mf_pre_feedback}\\
\phi^{k}(t,x)
&=
-\frac12\Lambda^{-1}
\left(P(t)x+\ell^k(t)+I^k(t)
+\Lambda(\varpi\chi_u\bar u^k(t)+\chi_g g^k(t))\right).
\label{eq:mf_post_feedback}
\end{align}
\end{subequations}
For \(k=1,\ldots,m\), set \(\phi^{\pi,k}:=\phi^{k}\). The feedback
family is a mean-field equilibrium with mean flow
\eqref{eq:mf_mu_closed_form} and value functions
\eqref{eq:mf_quadratic_value_ansatz}. The mean inventories
\(\bar X^0,\ldots,\bar X^m\) are the type averages of the induced inventories,
and \(\ell^0\) is the type average of \(\ell^{\pi,0}\).
\end{theorem}

The proof is given in Appendix~\ref{subsec:proof_mf_solution}.

\begin{remark}[State Decomposition]
\label{rem:mf_deviation_policy}
For the equilibrium trajectories, let \(X^{\pi,0}\) and \(\bar X^0\)
denote the type-\(\pi\) and mean pre-announcement inventories, respectively,
and let \(X^{\pi,k}\) and \(\bar X^k\) denote their continuations in
scenario \(k\). Define
\(\delta^{\pi,0}:=X^{\pi,0}-\bar X^0\) before the announcement and
\(\delta^{\pi,k}:=X^{\pi,k}-\bar X^k\) in scenario \(k\) after the
announcement. Continuity of the individual and mean inventories gives
\[
\delta^{\pi,k}(T_{\mathrm{ann}})
=\delta^{\pi,0}(T_{\mathrm{ann}}).
\]
Using \eqref{eq:mf_post_feedback} and
\eqref{eq:mf_post_mu_closed_form}, the individual feedback in scenario \(k\)
is
\begin{equation}\label{eq:mf_post_feedback_decomposition}
u^{\pi,k}(t)
=
\bar u^k(t)-\frac12\Lambda^{-1}P(t)\delta^{\pi,k}(t).
\end{equation}
The post-announcement deviation satisfies
\begin{equation}\label{eq:mf_post_deviation_dynamics}
\dot\delta^{\pi,k}(t)
=
-\frac12\Lambda^{-1}P(t)\delta^{\pi,k}(t),
\qquad
\delta^{\pi,k}(T_{\mathrm{ann}})
=\delta^{\pi,0}(T_{\mathrm{ann}}).
\end{equation}
Since neither the ODE nor its initial condition depends on \(k\), uniqueness
implies that \(\delta^{\pi,k}\) is identical across scenarios. The equilibrium
mean-inventory path \(\bar X^k\) around which \(\delta^{\pi,k}\) is centered
may nevertheless differ across scenarios because the augmented mean-field
common-state system
\eqref{eq:mf_y_system} is driven by the scenario-specific indexer flow \(g^k\).

Before the announcement, analogously to
\eqref{eq:pre_ann_affine_components}, define the centered affine coefficient
\[
\widetilde\ell^{\pi,0}:=\ell^{\pi,0}-\ell^0.
\]
Subtracting the type average of \eqref{eq:mf_pre_ell} from that equation gives
\begin{equation}\label{eq:mf_centered_affine_dynamics}
\dot{\widetilde\ell}^{\pi,0}(t)
=
\frac12P(t)\Lambda^{-1}\widetilde\ell^{\pi,0}(t),
\qquad
\widetilde\ell^{\pi,0}(T_{\mathrm{ann}})
=
\sum_{k=1}^m
(\pi^k-\bar\pi^k)\ell^k(T_{\mathrm{ann}}).
\end{equation}
Using \eqref{eq:mf_pre_feedback} and
\eqref{eq:mf_pre_mu_closed_form}, the individual feedback before the
announcement is
\begin{equation}\label{eq:mf_pre_feedback_decomposition}
u^{\pi,0}(t)
=
\bar u^0(t)-\frac12\Lambda^{-1}
\left(
P(t)\delta^{\pi,0}(t)+\widetilde\ell^{\pi,0}(t)
\right).
\end{equation}
The pre-announcement deviation satisfies
\begin{equation}\label{eq:mf_pre_deviation_dynamics}
\dot\delta^{\pi,0}(t)
=
-\frac12\Lambda^{-1}
\left(
P(t)\delta^{\pi,0}(t)+\widetilde\ell^{\pi,0}(t)
\right),
\qquad
\delta^{\pi,0}(0)=0.
\end{equation}
The dynamics of \(Y^0\) and \(\delta^{\pi,0}\), and likewise those of \(Y^k\)
and \(\delta^{\pi,k}\), decouple completely by
\eqref{eq:mf_consistency_system}, \eqref{eq:mf_post_deviation_dynamics}, and
\eqref{eq:mf_pre_deviation_dynamics}. Unlike
\eqref{eq:mf_post_deviation_dynamics},
\eqref{eq:mf_pre_deviation_dynamics} is forced by the centered affine
coefficient \(\widetilde\ell^{\pi,0}\), whose terminal value encodes the
belief deviation \(\pi-\bar\pi\) and whose dynamics are given by
\eqref{eq:mf_centered_affine_dynamics}. Consequently, a type-\(\pi\) player
can compute the equilibrium feedback from \((\bar\pi,\pi)\), without knowing
the full distribution \(\Pi\). This parallels the finite-player construction,
where trader \(i\) needs only \((\bar\pi_n,\pi_n^i)\), rather than the other
traders' individual beliefs.
\end{remark}

\subsection{Solvability of the Mean-Field System}

The equilibrium construction in Theorem~\ref{thm:mf_solution} is conditional
on solvability of the augmented mean-field common-state boundary-value problem
\eqref{eq:mf_consistency_system}. We now show that this problem is solvable
under Assumption~\ref{ass:standing_model} and show how it can be solved by
propagating the linear ODEs and solving a finite-dimensional linear system for
\(\ell^0(0)\). For
\(0\le s,t\le T\), let
\(\mathcal E(t,s)\in\mathbb R^{3d\times3d}\) be the state-transition matrix
of the homogeneous system
\(\dot{Y}(t)=\mathcal A(t)Y(t)\). Thus, for each fixed \(s\),
\[
\frac{\partial}{\partial t}\mathcal E(t,s)
=
\mathcal A(t)\mathcal E(t,s),
\qquad
\mathcal E(s,s)=\operatorname{Id}_{3d}.
\]
Following the coordinate order of \(Y^k\), use the block indices
\(X\), \(I\), and \(\ell\) for average inventory, the impact state, and the
affine coordinate, respectively. The \(d\times d\) block decomposition is
\[
\mathcal E(t,s)
=
\begin{pmatrix}
\mathcal E_{XX}(t,s)&\mathcal E_{XI}(t,s)&\mathcal E_{X\ell}(t,s)\\
\mathcal E_{IX}(t,s)&\mathcal E_{II}(t,s)&\mathcal E_{I\ell}(t,s)\\
\mathcal E_{\ell X}(t,s)&\mathcal E_{\ell I}(t,s)&\mathcal E_{\ell\ell}(t,s)
\end{pmatrix}.
\]
Define
\[
\mathcal E^{\mathrm{pre}}:=\mathcal E(T_{\mathrm{ann}},0),
\qquad
\mathcal E^{\mathrm{post}}:=\mathcal E(T,T_{\mathrm{ann}})
\]
and, for \(k=1,\ldots,m\),
\begin{equation}\label{eq:mf_post_forcing_ell}
\mathcal D_\ell^k
:=
\int_{T_{\mathrm{ann}}}^T
\begin{bmatrix}
\mathcal E_{\ell X}(T,s)&\mathcal E_{\ell I}(T,s)&\mathcal E_{\ell\ell}(T,s)
\end{bmatrix}
b^k(s)\,ds.
\end{equation}
Finally, set
\begin{equation}\label{eq:mf_matching_matrix}
\mathcal M:=\mathcal E_{\ell\ell}(T,0).
\end{equation}

The following lemma expresses the solution in terms of these transition blocks
and identifies the invertibility conditions that ensure uniqueness.

\begin{lemma}
\label{lem:mf_state_transition_representation}
Under Assumption~\ref{ass:standing_model}, if \(\mathcal M\) is invertible,
the initial affine coefficient is
\begin{equation}\label{eq:mf_pre_matching}
\ell^0(0)
=
-\mathcal M^{-1}\sum_{k=1}^m\bar\pi^k\mathcal D_\ell^k.
\end{equation}
The pre-announcement solution is then
\begin{equation}\label{eq:mf_pre_forward_solution}
Y^0(t)
=
\mathcal E(t,0)
\begin{pmatrix}
0\\
0\\
\ell^0(0)
\end{pmatrix},
\qquad
t\in[0,T_{\mathrm{ann}}].
\end{equation}
In particular,
\[
\begin{aligned}
\bar X^0(T_{\mathrm{ann}})
=\mathcal E^{\mathrm{pre}}_{X\ell}\ell^0(0),
\qquad
I^0(T_{\mathrm{ann}})
=\mathcal E^{\mathrm{pre}}_{I\ell}\ell^0(0),
\qquad
\ell^0(T_{\mathrm{ann}})
=\mathcal E^{\mathrm{pre}}_{\ell\ell}\ell^0(0).
\end{aligned}
\]
If \(\mathcal E^{\mathrm{post}}_{\ell\ell}\) is invertible, then, using these
announcement values, the post-announcement initial affine coefficient for
scenario \(k\) is
\begin{equation}\label{eq:mf_post_transition_ell}
\ell^k(T_{\mathrm{ann}})
=
-(\mathcal E^{\mathrm{post}}_{\ell\ell})^{-1}
\left(\mathcal E^{\mathrm{post}}_{\ell X}\bar X^0(T_{\mathrm{ann}})+
\mathcal E^{\mathrm{post}}_{\ell I}I^0(T_{\mathrm{ann}})+\mathcal D_\ell^k\right),
\qquad k=1,\ldots,m.
\end{equation}
Finally, the post-announcement solution in scenario \(k\) is
\begin{equation}\label{eq:mf_post_forward_solution}
Y^k(t)
=
\mathcal E(t,T_{\mathrm{ann}})
\begin{pmatrix}
\bar X^0(T_{\mathrm{ann}})\\
I^0(T_{\mathrm{ann}})\\
\ell^k(T_{\mathrm{ann}})
\end{pmatrix}
+
\int_{T_{\mathrm{ann}}}^t
\mathcal E(t,s)b^k(s)\,ds,
\qquad
t\in[T_{\mathrm{ann}},T].
\end{equation}
If both \(\mathcal M\) and \(\mathcal E^{\mathrm{post}}_{\ell\ell}\) are
invertible, these formulas determine a unique solution of the mean-field
boundary value problem.
\end{lemma}

The proof is given in
Appendix~\ref{subsec:proof_mf_state_transition_representation}. We now show that the procedure in
Lemma~\ref{lem:mf_state_transition_representation} can be carried out under the
market-parameter assumptions of Subsection~\ref{sec:opportunists}.

\begin{theorem}[Implementability of the Mean-Field Construction]
\label{thm:mf_implementability}
Under Assumption~\ref{ass:standing_model},
\(\mathcal E^{\mathrm{post}}_{\ell\ell}\) and the matching matrix
\(\mathcal M\) in
Lemma~\ref{lem:mf_state_transition_representation} are
invertible. Hence the mean-field construction is implementable and yields a
unique equilibrium within the affine construction of
Theorem~\ref{thm:mf_solution}.
\end{theorem}

The proof is given in
Appendix~\ref{subsec:proof_mf_implementability}.

\subsection{Convergence of the Scaled Finite-Player Game}

We consider finite-player games with \(x_0^i=0\) for every player and
\(\iota_0=0\) and show that the equilibria constructed in
Section~\ref{sec:full_procedure} under the scaling \(\varpi_n=\varpi/n\) converge to
the mean-field equilibrium of Definition~\ref{def:mf_equilibrium},
characterized in Theorem~\ref{thm:mf_solution}.

Recall that an \(n\)-subscript identifies a quantity from the \(n\)-player
equilibrium, while the corresponding symbol without an \(n\)-subscript denotes
its mean-field counterpart. Thus \(\bar\pi_n\) and \(\bar\pi\) denote the
finite-player and mean-field average beliefs, respectively.
\(X_n^{i,k}\), \(I_n^k\), and \(u_n^{i,k}\) denote player \(i\)'s finite-game
inventory, the finite-game impact state, and player \(i\)'s equilibrium trading
rate in scenario \(k\). The population-average inventory is \(\bar X_n^k\). Define the
population-average equilibrium rate by
\(\bar u_n^k:=n^{-1}U_n^k\).
The corresponding scenario components of the mean-field solution
 are \(\bar X^k\), \(I^k\), and \(\bar u^k\),
while \(X^{\pi,k}\) and \(u^{\pi,k}\) denote the type-\(\pi\)
representative player's \(k\)-component inventory and rate.

\begin{theorem}[Convergence to the Mean-Field Equilibrium]
\label{thm:finite_to_mf_convergence}
Under Assumption~\ref{ass:standing_model}, set
\(n_0:=\max\{2,\lceil\varpi\chi_u\rceil\}\). For each \(n\ge n_0\), consider the
finite-player game with population weight
\(\varpi_n=\varpi/n\) and
fixed instantaneous-flow loadings \(\chi_u,\chi_g\),
with \(x_0^i=0\) for every player \(i\) and \(\iota_0=0\). Let the
finite-player equilibrium be the one constructed in
Section~\ref{sec:full_procedure}. Set
\[
\varepsilon_n:=\frac1n+\|\bar\pi_n-\bar\pi\|.
\]
There is a
constant \(C\), independent of \(n\) and of the finite belief profile, such
that, for every \(n\ge n_0\),
\begin{equation}\label{eq:finite_to_mf_common_estimate}
\max\left\{
\sup_{t\in[0,T_{\mathrm{ann}}]}
\|\bar u_n^{0}(t)-\bar u^0(t)\|,
\max_{k=1,\ldots,m}\sup_{t\in[T_{\mathrm{ann}},T]}
\|\bar u_n^{k}(t)-\bar u^k(t)\|
\right\}
\le
C\varepsilon_n.
\end{equation}
Moreover, the player-level rate estimate holds uniformly over the finite
population:
\begin{equation}\label{eq:finite_to_mf_control_estimate}
\max_{1\le i\le n}\;
\max\left\{
\begin{aligned}
&\sup_{t\in[0,T_{\mathrm{ann}}]}
\|u_n^{i,0}(t)-u^{\pi_n^i,0}(t)\|,\\
&\max_{k=1,\ldots,m}\sup_{t\in[T_{\mathrm{ann}},T]}
\|u_n^{i,k}(t)-u^{\pi_n^i,k}(t)\|
\end{aligned}
\right\}
\le
C\varepsilon_n.
\end{equation}
Estimates \eqref{eq:finite_to_mf_common_estimate}--\eqref{eq:finite_to_mf_control_estimate}
also hold for the corresponding inventories and impact states.
\end{theorem}

The proof is given in
Appendix~\ref{subsec:proof_finite_to_mf_convergence}.
We next show that lifting the mean-field policies to the finite-player game
yields an approximate Nash equilibrium.

\begin{theorem}[Mean-Field Strategies Form an Approximate Nash Equilibrium]
\label{thm:mf_approximate_nash}
Under Assumption~\ref{ass:standing_model}, adopt the additional hypotheses and
notation of Theorem~\ref{thm:finite_to_mf_convergence}.
For each \(n\ge n_0\), lift the mean-field feedbacks from
Theorem~\ref{thm:mf_solution} to the finite-player history space by defining
the Borel non-anticipative policies
\begin{equation}\label{eq:mf_lifted_finite_profile}
\begin{aligned}
&\widehat\Phi_n^i
(t,\mathbf{x}^i,\bar{\mathbf{x}},\boldsymbol{\iota},\mathbf{a},\mathbf{s}^0;\bar p,p^i)
\\
&\quad:=
\mathbf 1_{\{t<T_{\mathrm{ann}}\}}
\phi^{p^i,0}(t,\mathbf{x}^i(t))
+
\mathbf 1_{\{t\geq T_{\mathrm{ann}}\}}
\sum_{k=1}^m
\mathbf 1_{\{\mathbf{a}(t)=\xi^k\}}
\phi^{k}(t,\mathbf{x}^i(t)),
\end{aligned}
\end{equation}
for \((\mathbf{x}^i,\bar{\mathbf{x}},\boldsymbol{\iota},\mathbf{a},\mathbf{s}^0)
\in D([0,T];\mathbb R^d)^5\) and
\((\bar p,p^i)\in\Delta_m^2\), and write
\(\widehat\Phi_n
:=(\widehat\Phi_n^1,\ldots,\widehat\Phi_n^n)\).
Then \(\widehat\Phi_n\in\mathcal U_n\) and there is a constant \(C\),
independent of \(n\) and of the finite belief profile, such that
\begin{equation}\label{eq:mf_profile_approximate_nash}
J_n^i(
\widehat\Phi_n^i;\widehat\Phi_n^{-i}
)
\leq
J_n^i(
\Psi_n^i;\widehat\Phi_n^{-i}
)
+
C\varepsilon_n^2
\end{equation}
for every \(i=1,\ldots,n\) and every policy \(\Psi_n^i\) satisfying
\((\Psi_n^i,\widehat\Phi_n^{-i})\in\mathcal U_n\). Thus
\(\widehat\Phi_n\) is a \(C\varepsilon_n^2\)-Nash equilibrium of the finite-player
game.
\end{theorem}

This approximation concerns the game from the prescribed initial state and is
not subgame perfect. Although the lifted policies respond to each player's own
inventory, they ignore the current average inventory and impact state. Hence
they need not remain approximate best responses after arbitrary off-equilibrium
histories.
The proof is given in
Appendix~\ref{subsec:proof_mf_approximate_nash}.

\begin{remark}[Random and Deterministic Belief Populations]\label{rem:mf_iid_beliefs}
Theorems~\ref{thm:finite_to_mf_convergence} and
\ref{thm:mf_approximate_nash} apply to every fixed belief profile.
If the profile is generated by i.i.d. sampling from \(\Pi\), then
\(\|\bar\pi_n-\bar\pi\|=O_p(n^{-1/2})\). The uniform path approximation
error is therefore \(O_p(n^{-1/2})\), and the lifted profile is an
\(O_p(n^{-1})\)-Nash equilibrium.
Alternatively, deterministic belief profiles satisfying
\(\|\bar\pi_n-\bar\pi\|=O(n^{-1})\) yield a uniform path approximation
error of \(O(n^{-1})\) and an \(O(n^{-2})\)-Nash equilibrium.
\end{remark}

\section{Illustrations}\label{sec:experiments}

We illustrate the model without seeking a complete empirical fit to real-world
index rebalances as the goal here is to showcase the qualitative properties
of the model. 

\subsection{Illustrative Setup}\label{sec:experiments_setup}

\paragraph{Market Parameters.}

We consider an index reconstitution involving \(d=5\) assets, with event times
\(T_{\mathrm{ann}}=10\), \(T_{\mathrm{TWAP}}=19\), \(T_{\mathrm{impl}}=20\), and \(T=22.\)
Time is measured in trading days, with one trading day equal to \(6.5\) trading
hours. Share quantities are
measured in millions of shares, and trading rates in millions of shares per
trading day. We assume that all coefficient matrices are diagonal:
\[
\Lambda=\operatorname{diag}(\lambda_1,\ldots,\lambda_d),
\qquad
\Gamma=\operatorname{diag}(\gamma_1,\ldots,\gamma_d),
\qquad
R=\operatorname{diag}(\rho_1,\ldots,\rho_d),
\]
\[
Q=\operatorname{diag}(q_1,\ldots,q_d),
\qquad
Q_T=\operatorname{diag}(q_{T,1},\ldots,q_{T,d}).
\]

We parameterize \(\Lambda\) and \(\Gamma\) using the estimates reported by
\cite[Table~A.7]{cartea_jaimungal_2016}. The instantaneous-impact coefficient
\(\hat k_a\) in \cite{cartea_jaimungal_2016} corresponds to our \(\lambda_a\), but because they measure trading
rates in shares per hour whereas we measure them in millions of shares per
trading day, we set \(\lambda_a=10^6\hat k_a/6.5\). Their permanent-impact
coefficient \(\hat b_a\) corresponds to our \(\gamma_a\), but because they
measure share quantities in individual shares whereas we measure them in
millions of shares, we set \(\gamma_a=10^6\hat b_a\). Their permanent-impact
model corresponds to \(R=0\) in our parameterization. To obtain transient
impact, we instead set \(\rho_a=2\log 2\) for every asset, giving the impact
state a half-life of half a trading day.

To choose \(q_a\), we suppress the impact state and strategic interactions
and consider an auxiliary scalar infinite-horizon problem with
\(\dot x_a=u_a\) and running cost
\(\lambda_a u_a^2+\frac{1}{2}q_a x_a^2\). Its optimal
feedback is \(u_a=-\sqrt{q_a/(2\lambda_a)}\,x_a\). We set
\(q_a/\lambda_a=0.25\,\mathrm{day}^{-2}\), which gives an inventory half-life
of approximately \(1.96\) trading days in the auxiliary problem. The
finite-horizon illustrations separately include a terminal penalty. We set
\(q_{T,a}/\lambda_a=13{,}000\,\mathrm{day}^{-1}\) to make terminal inventories
negligible.

We set \(f_a\) to \(100\) days of average daily volume (ADV).

Table~\ref{tab:experiment_asset_parameters} summarizes the data and
coefficients for the five selected stocks: Intel (INTC), Cisco (CSCO), eBay
(EBAY), Dow (DOW), and Costco (COST), with \(s_a^0\) denoting asset \(a\)'s
initial price.

\begin{table}[H]
\centering
\caption{Asset-Level Model Inputs and Coefficients}
\label{tab:experiment_asset_parameters}
\scriptsize
\resizebox{\textwidth}{!}{%
\begin{tabular}{clrrrrrrr}
\toprule
Asset & Ticker & \(s_a^0\) & \(\mathrm{ADV}_a\)
& \(\lambda_a\) & \(\gamma_a\) & \(\rho_a\) & \(q_a\) & \(q_{T,a}\) \\
\midrule
1 & INTC &  30.26 & 5.489503 & 0.067692 &  1.09 & 1.386294 & 0.016923 &    880 \\
2 & CSCO &  24.21 & 5.598870 & 0.037538 &  0.63 & 1.386294 & 0.009385 &    488 \\
3 & EBAY &  53.41 & 2.506048 & 0.384615 &  4.50 & 1.386294 & 0.096154 &  5,000 \\
4 & DOW  &  49.48 & 1.453332 & 0.463077 &  6.90 & 1.386294 & 0.115769 &  6,020 \\
5 & COST & 121.37 & 0.540008 & 1.630769 & 25.30 & 1.386294 & 0.407692 & 21,200 \\
\bottomrule
\end{tabular}}
\par\vspace{0.5\baselineskip}
\begin{minipage}{0.98\textwidth}
\footnotesize
\raggedright
\emph{Units.} The initial price \(s_a^0\) is measured in USD per share,
\(\mathrm{ADV}_a\) in millions of shares per trading day. The units of
\(\lambda_a\), \(\gamma_a\), \(\rho_a\),
\(q_a\), and \(q_{T,a}\) are, respectively,
\((\mathrm{USD/share})/(\mathrm{million\ shares/day})\),
\((\mathrm{USD/share})/(\mathrm{million\ shares})\),
\(\mathrm{day}^{-1}\),
\((\mathrm{USD/share})/(\mathrm{million\ shares}\cdot\mathrm{day})\), and
\((\mathrm{USD/share})/(\mathrm{million\ shares})\).
\par
\end{minipage}
\end{table}


We set \(x_0^i=0\) for every opportunist, \(\iota_0=0\), and
\(\chi_u=\chi_g=1\).
Further, we take \(S^0(t)\equiv s^0\), thereby suppressing fluctuations that would
affect realized execution prices and wealth but not the equilibrium strategies
or impact state.
Section~\ref{sec:experiments_finite_player} studies the
atomic-player model under the fixed-market-size scaling \(\varpi_n=1\).
Section~\ref{sec:experiments_mean_field} compares the finite-player model under
\(\varpi_n=30/n\) with its mean-field limit.

\paragraph{Membership Scenarios and Beliefs.}
The index contains exactly two of the five assets. Assets \(1\) and \(2\) are
the pre-rebalance constituents, and every two-asset subset is a possible
post-announcement scenario so that \(m=\binom{5}{2}=10\).

We model the population distribution of beliefs by
\(\Pi=\operatorname{Dirichlet}(\zeta\bar\pi)\) with
concentration parameter \(\zeta=50\). To construct the mean
\(\bar\pi\), we assign asset \(a\) the score \(\varsigma_a=6-a\) and compute
the unnormalized weight of pair \(\{a,a'\}\) as
\(\varsigma_a\varsigma_{a'}\). Normalizing the ten pair weights then gives
\(\bar\pi\). For a finite game with \(n\) opportunists, we draw their beliefs
independently from \(\Pi\).

We set \(h_a^0=0.07f_a\) for \(a=1,2\), corresponding to seven days of
ADV in each pre-rebalance constituent and USD~\(2.112\) billion in assets
under management.

Table~\ref{tab:experiment_scenarios} reports the
corresponding values of \(\bar\pi\). Table~\ref{tab:experiment_expected_flow}
reports the expected indexer flow under \(\bar\pi\), defined by
\(\bar g(t):=\sum_{k=1}^m\bar\pi^k g^k(t)\).

\begin{table}[H]
\centering
\begin{minipage}[t]{0.58\textwidth}
\centering
\captionsetup{justification=centering}
\caption{Post-Announcement Scenarios and Population-Average Belief}
\label{tab:experiment_scenarios}
\small
\begin{tabular}{crr}
\toprule
Scenario & Constituents & \(\bar\pi^k\) \\
\midrule
1  & \(\{1,2\}\) & 23.53\% \\
2  & \(\{1,3\}\) & 17.65\% \\
3  & \(\{1,4\}\) & 11.76\% \\
4  & \(\{1,5\}\) &  5.88\% \\
5  & \(\{2,3\}\) & 14.12\% \\
6  & \(\{2,4\}\) &  9.41\% \\
7  & \(\{2,5\}\) &  4.71\% \\
8  & \(\{3,4\}\) &  7.06\% \\
9  & \(\{3,5\}\) &  3.53\% \\
10 & \(\{4,5\}\) &  2.35\% \\
\bottomrule
\end{tabular}
\end{minipage}\hfill
\begin{minipage}[t]{0.36\textwidth}
\centering
\captionsetup{justification=centering}
\caption{Population-Average Expected Indexer Flow}
\label{tab:experiment_expected_flow}
\small
\begin{tabular}{cr}
\toprule
Asset & \([\bar g(t)]_a\) \\
\midrule
1 & \(-13.892459\) \\
2 & \(-15.644426\) \\
3 & \( 8.638579\) \\
4 & \( 4.487279\) \\
5 & \( 0.953451\) \\
\bottomrule
\end{tabular}
\par\vspace{0.5\baselineskip}
\footnotesize
\raggedright
\end{minipage}
\end{table}

\subsection{Atomic-Player Model}\label{sec:experiments_finite_player}

\begingroup
% Keep shorter figure floats on pages with text.
\renewcommand{\floatpagefraction}{0.85}

We begin with the finite-player model under the fixed-market-size interpretation
\(\varpi_n=1\), with \(n=20\) opportunists whose beliefs
\(\pi_n^i\), \(i=1,\ldots,n\), are i.i.d. samples from \(\Pi\).

\paragraph{Inventory.}

For scenario \(k\), let
\(L_n^k(t):=\sum_{i=1}^n X_n^{i,k}(t)\) denote the aggregate inventory
vector.
Figure~\ref{fig:experiments_aggregate_inventory} shows the aggregate inventory
\(L_n^k(t)\) while
Figure~\ref{fig:experiments_aggregate_trading_rates} shows the corresponding
aggregate trading rates, with separately scaled insets highlighting the short
bursts at \(t=0\) and \(T_{\mathrm{ann}}\). The paths in both figures depend on the sampled beliefs only through
\(\bar\pi_n\). Before the announcement, opportunists in aggregate have short positions in assets \(1\)
and \(2\) and long positions in assets \(3\), \(4\),
and \(5\), reflecting the expected indexer flow under \(\bar\pi_n\). After the
announcement, the aggregate inventory adjusts scenario by scenario to the realized indexer order.
Notably, for every asset added to or removed from the index, the aggregate inventory in that asset
passes through zero during the TWAP window, temporarily taking the opposite
sign, before being flattened. That is, during the indexer trading window the opportunits
sell more than what they held at the start of that time window.
The trading bursts at \(t=0\) and \(T_{\mathrm{ann}}\) in
Figure~\ref{fig:experiments_aggregate_trading_rates} reflect competition to
build positions before other opportunists' trades worsen entry prices.
The scale of \(\Lambda\) modulates the speed of entry by penalizing high
trading rates. Our relatively small choice of \(\Lambda\) permits pronounced
bursts. A larger \(\Lambda\) would make the bursts less pronounced and
inventory accumulation more gradual.

\begin{figure}[!htbp]
\centering
\includegraphics[width=\textwidth]{figures/finite_player_game/inventory/total_inventory/signed_total_inventory_all_scenarios_n020_varpi_1_chi_u_1_chi_g_1_closed_loop.pdf}
\caption{Aggregate inventory paths for \(n=20\). The vertical dotted lines mark \(T_{\mathrm{ann}}\), \(T_{\mathrm{TWAP}}\), and \(T_{\mathrm{impl}}\).}
\label{fig:experiments_aggregate_inventory}
\end{figure}

\begin{figure}[!htbp]
\centering
\includegraphics[width=\textwidth]{figures/finite_player_game/trading_rate/aggregate_trading_rate_all_scenarios_n020_varpi_1_chi_u_1_chi_g_1_closed_loop.pdf}
\caption{Aggregate opportunist trading rates \(U_n^k(t)\) for \(n=20\). The vertical dotted lines mark \(T_{\mathrm{ann}}\), \(T_{\mathrm{TWAP}}\), and \(T_{\mathrm{impl}}\).}
\label{fig:experiments_aggregate_trading_rates}
\end{figure}

Figure~\ref{fig:experiments_inventory_timing} plots the following three
quantities against various model parameters.
The first quantity is absolute inventory at \(T_{\mathrm{ann}}\),
averaged under the population-average belief \(\bar\pi\),
\(\mathsf{L}_n^{\mathrm{ann}}:=\sum_{k=1}^m\bar\pi^k\lVert L_n^k(T_{\mathrm{ann}})\rVert_1\). The second is
the corresponding quantity at \(T_{\mathrm{TWAP}}\),
\(\mathsf{L}_n^{\mathrm{TWAP}}:=\sum_{k=1}^m\bar\pi^k\lVert L_n^k(T_{\mathrm{TWAP}})\rVert_1\). The third,
\(\mathsf{L}_n^{\mathrm{chg}}:=\sum_{k=1}^m\bar\pi^k\lVert
L_n^k(T_{\mathrm{impl}})-\allowbreak L_n^k(T_{\mathrm{TWAP}})\rVert_1\), is the
absolute inventory change during the indexer's TWAP, averaged under
\(\bar\pi\).
It measures net aggregate inventory change during the TWAP, not total trading
volume.

In panel (a), the three quantities are plotted against the number of
opportunists \(n\). Aggregate inventory generally increases with \(n\),
with the largest increases at small \(n\). With multiple
opportunists,
\(\mathsf{L}_n^{\mathrm{chg}}>\mathsf{L}_n^{\mathrm{TWAP}}\), reflecting the
inventory sign reversals during the TWAP in
Figure~\ref{fig:experiments_aggregate_inventory}. Notably, in the case
of a single opportunistic trader there is no crossing through zero: During the
indexer implementation window the opportunist sells less than what he held 
at the start of the window. 

Panels (b)--(d) fix \(n=20\) and the sampled beliefs. Each panel varies one
coefficient while keeping the others at
their baseline values. In panels (b) and (d), the horizontal axis gives the
factor multiplying the corresponding baseline matrix.
In panel (b), increasing the scale of \(\Lambda\) reduces all three inventory
measures only slightly: Higher execution costs make inventory accumulation
before the TWAP more expensive, but also improve execution prices for trades
opposing the indexer as long as aggregate opportunist trading is greater than
aggregate opportunist trading. The slight decline in panel (b) shows
that these effects almost perfectly offset with the first one being slightly stronger.
In panel (c), increasing resilience reduces all three measures,
particularly inventory at \(T_{\mathrm{ann}}\): Faster decay weakens the price
impact carried from early opportunist trades into the TWAP, favoring later
accumulation, while quadratic execution costs limit how quickly positions can
be built thus prohibiting very large positions being built very shortly before \(T_{\mathrm{TWAP}}\).
In panel (d), the inventory measures generally increase with the
scale of \(\Gamma\), with smaller changes at larger scales: Stronger anticipated
price movements generated by the indexer's trades encourage opportunists to
build larger positions ahead of the TWAP pushing up the price at which they
can unload inventory on the indexer.

\begin{figure}[!htbp]
\centering
\includegraphics[width=\textwidth]{figures/finite_player_game/inventory/total_inventory/inventory_sensitivity_grid_n020_varpi_1_chi_u_1_chi_g_1_closed_loop.pdf}
\caption{Expected absolute aggregate inventory at \(T_{\mathrm{ann}}\) and \(T_{\mathrm{TWAP}}\), denoted by \(\mathsf{L}_n^{\mathrm{ann}}\) and \(\mathsf{L}_n^{\mathrm{TWAP}}\), and the absolute net inventory change \(\mathsf{L}_n^{\mathrm{chg}}\) during the indexer's TWAP. Panels (b)--(d) fix \(n=20\). The vertical dotted lines mark baseline values, and the horizontal dotted lines show the expected gross indexer order.}
\label{fig:experiments_inventory_timing}
\end{figure}

Figure~\ref{fig:experiments_player_asset2} shows the asset \(2\) inventories of
the players whose expected indexer flows are lowest, closest to the expected
flow under \(\bar\pi_n\), and highest for that asset. Even the player with
the highest expected indexer flow expects negative flow for asset \(2\).
However, because this expectation is less negative than the expected flow under
\(\bar\pi_n\), the player carries a positive inventory into the
announcement. Once the scenario is announced, the players' pre-announcement beliefs no longer
enter the continuation problem directly, but their effect persists through the
different inventories carried into the announcement. The continuation paths
converge toward zero from those different positions.

\begin{figure}[!htbp]
\centering
\includegraphics[width=\textwidth]{figures/finite_player_game/inventory/player_inventory/player_inventory_asset_02_n020_varpi_1_chi_u_1_chi_g_1_closed_loop.pdf}
\caption{Individual inventory paths for asset \(2\).}
\label{fig:experiments_player_asset2}
\end{figure}

\paragraph{Trading Volume.}
The differences in individual inventories can generate trading that cancels
in the aggregate. To isolate the effect of belief dispersion, we replace the
sampled beliefs by
\[
\pi_n^i(\delta)=(1-\delta)\bar\pi_n+\delta\pi_n^i,
\qquad \delta\in[0,1].
\]
The original illustration corresponds to \(\delta=1\), and \(\delta=0\)
gives identical beliefs. Writing \(u_n^{i,k}(t;\delta)\) for the resulting
full-horizon control path in scenario \(k\), including the common
pre-announcement interval, define expected cumulative gross volume and
cumulative absolute aggregate flow by
\begin{align}
G_{n,\delta}(t)
&:=\sum_{k=1}^m\bar\pi_n^k\int_0^t
\sum_{i=1}^n\lVert u_n^{i,k}(s;\delta)\rVert_1\,ds,
\label{eq:experiments_gross_volume}\\
N_{n,\delta}(t)
&:=\sum_{k=1}^m\bar\pi_n^k\int_0^t
\left\lVert\sum_{i=1}^n u_n^{i,k}(s;\delta)\right\rVert_1\,ds.
\label{eq:experiments_net_volume}
\end{align}
The half-difference \(\mathcal H_{n,\delta}(t):=
\frac{1}{2}[G_{n,\delta}(t)-N_{n,\delta}(t)]\) measures contemporaneous
purchases and sales that can be matched within the opportunist population,
counting each matched share once. Figure~\ref{fig:experiments_volume_time}
plots both cumulative quantities.
The shaded gap counts both sides of the offsetting trades.
Belief dispersion generates substantial additional gross trading volume in
this illustration.


\begin{figure}[!htbp]
\centering
\includegraphics[width=\textwidth]{figures/finite_player_game/traded_volume/opportunist_volume_over_time_n020_varpi_1_chi_u_1_chi_g_1_closed_loop.pdf}
\caption{Expected cumulative gross opportunist volume and cumulative absolute aggregate flow for \(n=20\). The shaded gap is \(2\mathcal H_{20,1}(t)\).}
\label{fig:experiments_volume_time}
\end{figure}

Varying \(\delta\) preserves the average belief \(\bar\pi_n\) and leaves aggregate
signed trading rates and hence \(L_n^k\) and \(N_{n,\delta}\) unchanged, but
individual rates can move in opposite directions, as shown in
Figure~\ref{fig:experiments_player_asset2}. In the left panel of 
Figure~\ref{fig:experiments_volume_dispersion} we plot total gross volume
\(G_{n,\delta}(T)\) as a function of the dispersion \(
\mathcal D_{n,\delta}
:=\left(\frac{1}{n}\sum_{i=1}^n
\lVert\pi_n^i(\delta)-\bar\pi_n\rVert_2^2\right)^{1/2}
\) for various values of \(n\).
The right panel shows the increase in
gross volume relative to identical beliefs, normalized by the number of opportunists:
\([G_{n,\delta}(T)-G_{n,0}(T)]/n\). Both panels plot against the
root-mean-square belief dispersion \(\mathcal D_{n,\delta}\).
Gross volume increases with belief dispersion,
even though aggregate inventory is unchanged.
The right panel also shows that, at the same positive belief dispersion,
the additional gross volume per opportunist is larger for higher \(n\).
The reason for this is the following: At larger \(n\), each opportunist's
share of aggregate trading is smaller.
Thus, belief differences more readily turn buying less than the average rate into selling,
or vice versa for selling. The resulting opposing trades increase gross volume per 
opportunist while canceling in aggregate.


\begin{figure}[!htbp]
\centering
\includegraphics[width=\textwidth]{figures/finite_player_game/traded_volume/belief_dispersion_sensitivity/opportunist_trading_volume_by_belief_dispersion_n002_005_010_020_050_100_varpi_1_chi_u_1_chi_g_1_closed_loop.pdf}
\caption{Total gross volume and its dispersion-driven increase per opportunist under mean-preserving changes in beliefs. Scenario expectations use the fixed sample-average belief \(\bar\pi_n\) for each \(n\).}
\label{fig:experiments_volume_dispersion}
\end{figure}

\paragraph{Impact.}

Figure~\ref{fig:experiments_impact} shows that the impact state follows the same
qualitative pattern as the aggregate inventory: its pre-announcement path
reflects the expected indexer order under \(\bar\pi_n\), while each
post-announcement path adjusts to the realized indexer order.
Although the impact state is continuous, it adjusts rapidly when trading begins
at \(t=0\) and when the scenario is revealed at \(T_{\mathrm{ann}}\), and more
modestly when the indexer's TWAP starts at \(T_{\mathrm{TWAP}}\) and ends at
\(T_{\mathrm{impl}}\). The rapid initial and announcement-time impact
adjustments reflect the trading bursts in
Figure~\ref{fig:experiments_aggregate_trading_rates}.
If \(\Lambda\) were larger, the trading bursts would be less pronounced and
impact would adjust more gradually.
Once a burst subsides and before the TWAP begins, continuing opportunist
trading largely offsets impact decay in \eqref{eq:transient_impact}, allowing
impact to remain nearly constant while inventories continue to change.

\begin{figure}[!htbp]
\centering
\includegraphics[width=\textwidth]{figures/finite_player_game/impact_state/impact_state_all_scenarios_n020_varpi_1_chi_u_1_chi_g_1_closed_loop.pdf}
\caption{Impact-state paths for \(n=20\).}
\label{fig:experiments_impact}
\end{figure}

Figure~\ref{fig:experiments_impact_sensitivity} reports the expected
absolute impact magnitude \(\sum_{k=1}^m\bar\pi^k\lVert I_n^k(t)\rVert_1\)
at \(T_{\mathrm{ann}}\), \(T_{\mathrm{TWAP}}\), \(T_{\mathrm{impl}}\), and
\(T\) for the same four sweeps as the inventory grid.
Over the plotted range in panel (a), the four magnitudes level off.
At larger \(n\), the impact magnitudes at \(T_{\mathrm{TWAP}}\),
\(T_{\mathrm{impl}}\), and \(T\) are nearly equal and remain higher than at
\(T_{\mathrm{ann}}\).
As \(n\) increases, announcement impact generally rises, while the 
implementation-time and TWAP-start impacts shrink and converge to the same value.
In panel (b), impact at \(T_{\mathrm{impl}}\) increases with \(\Lambda\),
while the other three magnitudes change little. The indexer's schedule is
fixed, so the increase mainly reflects weaker offsetting by opportunists
during the TWAP. This is consistent with the smaller TWAP inventory change
in Figure~\ref{fig:experiments_inventory_timing}.
In panels (c) and (d), all four magnitudes decrease with \(\rho\) and
increase with the scale of \(\Gamma\), consistent with faster impact decay
and stronger price displacement from a given net trading rate, respectively.

\begin{figure}[!htbp]
\centering
\includegraphics[width=\textwidth]{figures/finite_player_game/impact_state/impact_state_sensitivity_grid_n020_varpi_1_chi_u_1_chi_g_1_closed_loop.pdf}
\caption{Expected absolute impact magnitude at the announcement, TWAP start, implementation time, and terminal time. Panels (b)--(d) fix \(n=20\). The vertical dotted lines mark baseline values.}
\label{fig:experiments_impact_sensitivity}
\end{figure}

\paragraph{Indexer Cost.}

For scenario \(k\), we define the indexer's implementation shortfall relative
to the fundamental price by
\begin{equation}\label{eq:experiments_indexer_cost}
C_{n,\mathrm{idx}}^k
:=
\int_{T_{\mathrm{TWAP}}}^{T_{\mathrm{impl}}}
(g^k(t))^\top
\left[
I_n^k(t)+\Lambda\bigl(U_n^k(t)+g^k(t)\bigr)
\right]dt.
\end{equation}

Write \(\bar C_{n,\mathrm{idx}}:=\sum_{k=1}^m\bar\pi^k
C_{n,\mathrm{idx}}^k\), and let \(\bar C_{0,\mathrm{idx}}\) denote the
corresponding cost without opportunists. For the same initial portfolio and
fixed rebalance orders, \(\bar C_{0,\mathrm{idx}}-\bar C_{n,\mathrm{idx}}\)
equals the increase in expected indexer wealth at \(T_{\mathrm{impl}}\)
relative to the no-opportunist case, including cash and valuing stocks at
the fundamental price \(s^0\).

For the common resilience \(R=\rho\operatorname{Id}_d\), define
\(w_\rho(r):=\int_0^r e^{-\rho s}\,ds\). Splitting the opportunists' impact
at \(T_{\mathrm{TWAP}}\) in \eqref{eq:transient_impact} gives
\begin{equation}\label{eq:experiments_savings_decomposition}
\bar C_{0,\mathrm{idx}}-\bar C_{n,\mathrm{idx}}
=-\mathsf P_n+\mathsf B_n+\mathsf E_n,
\end{equation}
where
\begin{align}
\mathsf P_n
&:=\frac{w_\rho(\Delta_{\mathrm{TWAP}})}{\Delta_{\mathrm{TWAP}}}
\sum_{k=1}^m\bar\pi^k(\Delta h^k)^\top I_n^k(T_{\mathrm{TWAP}}),
\label{eq:experiments_savings_pre_twap}\\
\mathsf B_n
&:=-\frac{1}{\Delta_{\mathrm{TWAP}}}
\sum_{k=1}^m\bar\pi^k
\int_{T_{\mathrm{TWAP}}}^{T_{\mathrm{impl}}}
w_\rho(T_{\mathrm{impl}}-t)(\Delta h^k)^\top\Gamma U_n^k(t)\,dt,
\label{eq:experiments_savings_during_twap}\\
\mathsf E_n
&:=-\frac{1}{\Delta_{\mathrm{TWAP}}}
\sum_{k=1}^m\bar\pi^k(\Delta h^k)^\top\Lambda
\bigl[L_n^k(T_{\mathrm{impl}})-L_n^k(T_{\mathrm{TWAP}})\bigr].
\label{eq:experiments_savings_execution}
\end{align}
All three quantities are positive in the plotted sweeps.
The cost \(\mathsf P_n\) comes from price displacement caused by opportunists
before the TWAP that persists during the indexer's execution.
The benefit \(\mathsf B_n\) is the reduction in the indexer's impact cost
from opposing trades during the TWAP, and \(\mathsf E_n\) is the saving in the
instantaneous execution-cost term. For a stock the indexer buys, earlier
opportunist purchases contribute to \(\mathsf P_n\), while opportunist sales
during the TWAP contribute to \(\mathsf B_n\) and \(\mathsf E_n\).
The trading directions reverse for an indexer sale.

Figure~\ref{fig:experiments_indexer_cost}
plots total expected cost savings and the signed contributions
\(-\mathsf P_n\), \(\mathsf B_n\), and \(\mathsf E_n\) in
\eqref{eq:experiments_savings_decomposition} against the number of
opportunists, the scale of \(\Lambda\), resilience \(\rho\), and the scale
of \(\Gamma\).
Each comparison uses the no-opportunist benchmark at the same parameter values.
In panel (a), savings rise mainly at small \(n\), as additional opportunists
offset more of the indexer's trading and leave less adverse impact at the
TWAP start. In panel (b), larger \(\Lambda\) increases the execution-cost
saving \(\mathsf E_n\) despite the slight reduction in opposing trading.
The increase in \(\mathsf E_n\) outweighs the small decline in net
transient-impact savings \(\mathsf B_n-\mathsf P_n\).
In panel (c), savings peak at intermediate resilience, where much of the
adverse impact from earlier trading has decayed but opposing trades during
the TWAP still provide substantial benefits. In contrast, at small
resilience, impact is close to permanent and opportunists can be harmful
when \(\Lambda\) is small relative to \(\Gamma\), as in our
parameterization. Their net effect on the indexer's cost through the
mid-price is then adverse, \(\mathsf B_n-\mathsf P_n<0\), and the
execution-cost saving \(\mathsf E_n\) from opposing the indexer's trades
is too small to compensate. When \(\Lambda\) is larger (not shown here), the execution-cost
saving can outweigh the adverse net price effect, making opportunists
beneficial even at low resilience.
This dependence on execution costs relative to permanent impact is
consistent with \citet[Proposition~2]{PegoraroSammonShim2026}, whose
discrete-time model establishes cost reductions for an index fund trading
entirely at implementation when instantaneous execution costs are
sufficiently large relative to permanent impact.
At the other extreme, large resilience reduces the TWAP impact benefit
\(\mathsf B_n\), as the indexer's own price impact also decays rapidly.
In panel (d), larger \(\Gamma\) amplifies the benefit \(\mathsf B_n\) of
opposing trades during the TWAP more than the adverse impact cost
\(\mathsf P_n\) of earlier trading.

\begin{figure}[!htbp]
\centering
\includegraphics[width=\textwidth]{figures/finite_player_game/indexer_cost/indexer_cost_sensitivity_grid_n020_varpi_1_chi_u_1_chi_g_1_closed_loop.pdf}
\caption{Expected indexer cost savings (black) and the signed components of \eqref{eq:experiments_savings_decomposition}, in USD million. Panels (b)--(d) fix \(n=20\). Horizontal dashed lines mark zero, and vertical dotted lines mark baseline parameter values.}
\label{fig:experiments_indexer_cost}
\end{figure}

To examine timing sensitivity, we fix \(n=20\), retain the baseline
announcement time and TWAP duration, and vary
\(T_{\mathrm{impl}}-T_{\mathrm{ann}}\) and \(T-T_{\mathrm{impl}}\).
Figure~\ref{fig:experiments_indexer_cost_timing} reports
\(\bar C_{0,\mathrm{idx}}-\bar C_{20,\mathrm{idx}}\) in USD million.
Savings increase in both timing dimensions. The marginal benefit of extending either
interval is largest when the other interval is short. More time before
 \(T_{\mathrm{impl}}\) lets the opportunists adjust further ahead
of the indexer's TWAP. More time after \(T_{\mathrm{impl}}\) lets them take
larger positions against the indexer's flow during the TWAP, carry those
positions through \(T_{\mathrm{impl}}\), and unwind them only after the indexer
has finished. In both cases, aggregate opportunist flow offsets more of the
indexer order during the TWAP.

\begin{figure}[!htbp]
\centering
\includegraphics[width=0.97\textwidth]{figures/finite_player_game/indexer_cost/indexer_cost_timing_sensitivity_n020_varpi_1_chi_u_1_chi_g_1_closed_loop.pdf}
\caption{Timing sensitivity of expected indexer cost savings for \(n=20\), in USD million. The marker denotes baseline timing.}
\label{fig:experiments_indexer_cost_timing}
\end{figure}

\paragraph{Player Wealth.}

An opportunist's terminal trading wealth is terminal cash plus terminal inventory
marked at the fundamental price \(s^0\). We average this quantity across
players and post-announcement scenarios under \(\bar\pi\) to obtain mean
wealth \(\bar W_n\). The running and terminal
inventory penalties affect equilibrium strategies but are not deducted from
this financial-wealth measure.
\begin{samepage}
Under \(\varpi_n=\chi_u=\chi_g=1\) and \(x_0^i=0\),
\eqref{eq:execution_price} gives
\begin{equation}\label{eq:experiments_player_wealth_decomposition}
n\bar W_n=\mathsf G_n+\mathsf E_n-\mathsf K_n,
\end{equation}
where
\begin{align}
\mathsf G_n
&:=-\sum_{k=1}^m\bar\pi^k\int_0^T
(U_n^k(t))^\top I_n^k(t)\,dt,
\label{eq:experiments_player_wealth_midprice}\\
\mathsf K_n
&:=\sum_{k=1}^m\bar\pi^k\int_0^T
(U_n^k(t))^\top\Lambda U_n^k(t)\,dt.
\label{eq:experiments_player_wealth_execution}
\end{align}
\end{samepage}
The midprice contribution \(\mathsf G_n\) includes gains from trading
against the indexer's impact and costs from the opportunists' own impact.
The term \(\mathsf K_n\) is the instantaneous execution cost generated by
aggregate opportunist trading over the entire horizon. Subtracting it from
the benefit \(\mathsf E_n\) of opposing the indexer's flow gives the net
execution contribution.

Figure~\ref{fig:experiments_player_wealth} plots total wealth \(n\bar W_n\)
and its signed contributions \(\mathsf G_n\), \(\mathsf E_n\), and
\(-\mathsf K_n\) for the four parameter sweeps.
Panel (a) also shows mean wealth \(\bar W_n\). Both mean and total wealth
decline with \(n\): more opportunists offset a larger fraction of the fixed
indexer orders, reducing the price displacement from which the group profits.
In panel (b), total wealth increases with the scale of \(\Lambda\).
Larger \(\Lambda\) makes position building before \(T_{\mathrm{TWAP}}\) and
unwinding after \(T_{\mathrm{impl}}\) more costly, but improves execution
prices for sales during the TWAP when the aggregate opportunist selling
rate is smaller than the indexer's buying rate. The latter benefit
prevails, increasing the net execution contribution.
In panel (c), total wealth decreases with resilience. Opportunists carry
smaller positions into the TWAP, as shown in
Figure~\ref{fig:experiments_inventory_timing}(c), and faster impact decay
reduces the favorable price displacements from which they profit, lowering
\(\mathsf G_n\).
In panel (d), total wealth increases with the scale of \(\Gamma\).
For a stock bought by the indexer, opportunists' purchases before the TWAP
generate a larger price displacement, raising the prices at which they
subsequently sell. During the TWAP, larger \(\Gamma\) benefits opportunists
as long as the indexer's buying rate exceeds their aggregate selling rate.
Although, higher prices also make covering short positions after \(T_{\mathrm{impl}}\)
more costly, the gains from selling during the TWAP at a higher price outweigh the
additional acquisition and unwinding costs.

\begin{figure}[!htbp]
\centering
\includegraphics[width=\textwidth]{figures/finite_player_game/player_wealth/player_wealth_sensitivity_grid_n020_varpi_1_chi_u_1_chi_g_1_closed_loop.pdf}
\caption{Total terminal trading wealth and its signed contributions \(\mathsf G_n\), \(\mathsf E_n\), and \(-\mathsf K_n\), in USD million. Panel (a) also shows mean wealth per opportunist. Panels (b)--(d) fix \(n=20\). Vertical axes are linear except in panel (b), which uses a symmetric logarithmic scale with a linear region from \(-1\) to \(1\) USD million. The vertical dotted lines mark baseline values.}
\label{fig:experiments_player_wealth}
\end{figure}

\FloatBarrier
\endgroup

\subsection{Mean-Field Limit}\label{sec:experiments_mean_field}

We compare the finite-player model with \(\varpi_n=30/n\)
and its mean-field limit. The restriction \(\chi_u\varpi_n\leq1\) in
Assumption~\ref{ass:standing_model} requires \(n\geq30\) because \(\chi_u=1\).
The game with \(n\) players uses the first \(n\)
draws from a single sequence of independent beliefs with law \(\Pi\).
Figures~\ref{fig:experiments_mf_inventory_scenario2} and
\ref{fig:experiments_mf_impact_scenario2} compare the finite-player and
mean-field inventory and impact paths in Scenario \(2\) for
\(n=30,60,150,300\). The paths are close to their mean-field limits at
the larger population sizes. The midprice reacts
almost immediately when trading starts and when the membership scenario
is announced. The legend reports
\(\bar\pi_n^k\) for each finite-player curve and \(\bar\pi^k\) for the
mean-field curve.
\begin{figure}[H]
\centering
\includegraphics[width=\textwidth]{figures/mean_field_game/inventory/average_inventory_varpi_30_over_n_chi_u_1_chi_g_1_n030_060_150_300_scenario_02.pdf}
\caption{Mean-field convergence of sample-average inventory in
Scenario \(2\) under \(\varpi_n=30/n\).}
\label{fig:experiments_mf_inventory_scenario2}
\end{figure}

\begin{figure}[H]
\centering
\includegraphics[width=\textwidth]{figures/mean_field_game/impact_state/impact_state_varpi_30_over_n_chi_u_1_chi_g_1_n030_060_150_300_scenario_02.pdf}
\caption{Mean-field convergence of the impact state in Scenario \(2\)
under \(\varpi_n=30/n\).}
\label{fig:experiments_mf_impact_scenario2}
\end{figure}

\FloatBarrier

\section{Conclusion}\label{sec:conclusion}

We have constructed a subgame-perfect feedback equilibrium for strategic
trading before and after an index announcement, with global solvability under
the model's standing assumptions. The equilibrium separates the effects of
average membership beliefs on aggregate inventories and price impact from
the effects of disagreement on individual positions. The mean-field limit
provides tractable policies with quantitative finite-player error bounds.
Our numerical illustrations show how competition and impact decay determine
whether anticipatory trading raises or lowers the indexer's execution costs.

Future work could allow traders to update their beliefs about index membership
dynamically as new information arrives before the announcement. Another
extension is to let traders infer other players' beliefs from observed prices
and trading activity, including learning the population-average belief that
is taken as known here.

\FloatBarrier

\section*{Acknowledgements}
J. Blanchet gratefully acknowledges support from ONR under award N00014-24-1-2655, and the National Science Foundation (NSF) under grants 2312204 and 2403007.

\bibliographystyle{plainnat}
\bibliography{refs}

\clearpage
\appendix
\section{Proofs of Section~\ref{sec:solution}}\label{sec:proofs_solution}

\subsection{Fundamental-Price Stochastic Integral}
\label{subsec:proof_fundamental_price_stochastic_integral}

\begin{lemma}[Fundamental-Price Stochastic Integral]
\label{lem:fundamental_price_stochastic_integral}
Under Assumption~\ref{ass:standing_model}, for every admissible profile and
trader \(i\), the process
\[
\mathcal N_n^i(0):=0,
\qquad
\mathcal N_n^i(t)
:=
\int_{(0,t]}(X_n^i(s-))^\top\,dS^0(s),
\qquad 0<t\leq T,
\]
is a uniformly integrable martingale under \(\mathbb P_n^i\).
\end{lemma}

\begin{proof}
Fix an admissible profile and trader \(i\), and evaluate all processes under
that profile. Independence of \(S^0\) and \(\Xi\) ensures that \(S^0\) remains
a martingale in the filtration generated by \((S^0,A)\). The strong solution
makes \(X_n^i\) continuous and adapted to this filtration, so
\(\mathcal N_n^i\) is a local martingale starting from zero. Applying
\eqref{eq:fundamental_price_integration_by_parts} up to time \(t\) and using
\eqref{eq:inventory_dynamics} gives
\[
\sup_{t\leq T}|\mathcal N_n^i(t)|
\leq
2\sup_{t\leq T}\lVert S^0(t)\rVert
\left(\lVert x_0^i\rVert+\int_0^T\lVert u_n^i(s)\rVert\,ds\right).
\]
The price supremum is square-integrable under \(\mathbb P_n^i\) by Doob's
\(L^2\) maximal inequality. The other factor is square-integrable by
Cauchy--Schwarz and admissibility. Hence \(\mathcal N_n^i\) has an integrable
absolute supremum. Localization and dominated convergence therefore show that
\(\mathcal N_n^i\) is a true martingale. Domination by the integrable absolute
supremum also gives uniform integrability.
\end{proof}

\subsection{Proof of Theorem~\ref{thm:stitched_hjb_verification}}
\label{subsec:proof_stitched_hjb_verification}

\begin{proof}
We first verify the best-response inequality for a continuation game starting
at or after the announcement, so \(s\geq T_{\mathrm{ann}}\) and \(a=\xi^k\).
Let trader \(i\) choose an admissible unilateral continuation deviation
\(\Psi_n^i\), which induces a potentially stochastic rate \(u_n^i(t)\) through
its dependence on the observed history of \(S^0\). To make this rate
deterministic, fix a full path \(\mathbf{s}^0\) extending \(\mathbf{s}^0_{[0,s]}\) and write
\(u_n^i(t;\mathbf{s}^0)\) for the resulting rate. We suppress the same path argument in
the other induced processes.
Equation \eqref{eq:stitched_hjb} and the chain rule
give outside of the switching times of \(g^k\)
\[
\begin{aligned}
&\mathcal L\left(
X_n^i(t),u_n^i(t;\mathbf{s}^0);
I(t),\varpi_nU_n^{-i}(t),g^k(t)
\right)
+\frac{d}{dt}V_n^{i,k}\bigl(t,Z_n^i(t)\bigr)
\\
&\quad=
\partial_tV_n^{i,k}\bigl(t,Z_n^i(t)\bigr)
+
\widehat{\mathscr H}_n^{i,k}\left(
t,X_n(t),I(t),u_n^i(t;\mathbf{s}^0),
\nabla V_n^{i,k}\bigl(t,Z_n^i(t)\bigr);
\phi_n^{-i,k}
\right)
\\
&\quad=
\widehat{\mathscr H}_n^{i,k}\left(
t,X_n(t),I(t),u_n^i(t;\mathbf{s}^0),
\nabla V_n^{i,k}\bigl(t,Z_n^i(t)\bigr);
\phi_n^{-i,k}
\right)
\\
&\qquad-
\min_{v\in\mathbb R^d}
\widehat{\mathscr H}_n^{i,k}\left(
t,X_n(t),I(t),v,
\nabla V_n^{i,k}\bigl(t,Z_n^i(t)\bigr);
\phi_n^{-i,k}
\right)
\geq 0.
\end{aligned}
\]
Integrating interval by interval over the intervals where \(g^k\) is constant
and using \eqref{eq:stitched_hjb_terminal} gives
\begin{equation}\label{eq:post_announcement_pathwise_bound}
\begin{aligned}
V_n^{i,k}\left(s,z^i(x_n^i,\bar x_n,\iota)\right)
\leq{}&
\int_s^T
\mathcal L\left(
X_n^i(t),u_n^i(t;\mathbf{s}^0);
I(t),\varpi_nU_n^{-i}(t),g^k(t)
\right)\,dt
\\
&+
\frac{1}{2}X_n^i(T)^\top Q_TX_n^i(T).
\end{aligned}
\end{equation}
Taking the conditional expectation under \(\mathbb P_n^i\) given
\(S^0|_{[0,s]}=\mathbf{s}^0_{[0,s]}\) and \(A(s)=\xi^k\) gives
\[
J_n^i\left(
\Psi_n^i;\phi_n^{-i}
\mid s,x_n,\iota,\xi^k,\mathbf{s}^0_{[0,s]}
\right)
\geq
V_n^{i,k}\left(s,z^i(x_n^i,\bar x_n,\iota)\right).
\]

We next verify the best-response inequality for a continuation game starting
before the announcement, so \(s<T_{\mathrm{ann}}\) and \(a=0\). Applying the
same HJB argument on \([s,T_{\mathrm{ann}}]\), followed by
the value-matching condition \eqref{eq:stitched_hjb_value_matching} at the
state \(Z_n^i(T_{\mathrm{ann}})\) and
the post-announcement bound \eqref{eq:post_announcement_pathwise_bound} in each
scenario, gives
\begin{equation}\label{eq:pre_announcement_pathwise_bound}
\begin{aligned}
V_n^{i,0}\left(s,z^i(x_n^i,\bar x_n,\iota)\right)
&\leq
\int_s^{T_{\mathrm{ann}}}
\mathcal L\left(
X_n^i(t),u_n^i(t;\mathbf{s}^0);
I(t),\varpi_nU_n^{-i}(t),0
\right)\,dt
\\
&+
\sum_{k=1}^m\pi_n^{i,k}
V_n^{i,k}\bigl(
T_{\mathrm{ann}},Z_n^i(T_{\mathrm{ann}})
\bigr)\\
&\leq
\int_s^{T_{\mathrm{ann}}}
\mathcal L\left(
X_n^i(t),u_n^i(t;\mathbf{s}^0);
I(t),\varpi_nU_n^{-i}(t),0
\right)\,dt
\\
&+
\sum_{k=1}^m\pi_n^{i,k}
\Bigg[
\int_{T_{\mathrm{ann}}}^T
\mathcal L\left(
X_n^i(t),u_n^i(t;\mathbf{s}^0);
I(t),\varpi_nU_n^{-i}(t),g^k(t)
\right)\,dt
\\
&\hspace{105pt}
+\frac{1}{2}X_n^i(T)^\top Q_TX_n^i(T)
\Bigg]
\\
&=
\mathbb E_n^i\Bigg[
\int_s^T
\mathcal L\left(
X_n^i(t),u_n^i(t;\mathbf{s}^0);
I(t),\varpi_nU_n^{-i}(t),g(t)
\right)\,dt
\\
&\hspace{48pt}
+\frac{1}{2}X_n^i(T)^\top Q_TX_n^i(T)
\,\Bigm|\,
S^0|_{[0,T]}=\mathbf{s}^0,\ A(s)=0
\Bigg].
\end{aligned}
\end{equation}
The final equality in \eqref{eq:pre_announcement_pathwise_bound} is the
conditional expectation over \(\Xi\) with the full path of \(S^0\) fixed at
\(\mathbf{s}^0\). Independence of \(S^0\) and \(\Xi\) leaves the conditional probability of
scenario \(k\) equal to \(\pi_n^{i,k}\). Integrating
\eqref{eq:pre_announcement_pathwise_bound} with respect to trader \(i\)'s
conditional future law of \(S^0\) given
\(S^0|_{[0,s]}=\mathbf{s}^0_{[0,s]}\) gives
\[
J_n^i\left(
\Psi_n^i;\phi_n^{-i}
\mid s,x_n,\iota,0,\mathbf{s}^0_{[0,s]}
\right)
\geq
V_n^{i,0}\left(s,z^i(x_n^i,\bar x_n,\iota)\right).
\]

For \(\Psi_n^i=\phi_n^i\), \eqref{eq:stitched_hjb_argmin} makes all inequalities
equalities. Thus the candidate attains the lower bounds above, whereas every
admissible deviation has weakly larger cost. Since the trader and
continuation game were arbitrary, \eqref{eq:subgame_perfect_ne} holds, and the
asserted continuation-value identities follow.
\end{proof}

\subsection{Proof of Theorem~\ref{thm:full_horizon_closed_loop_equilibrium}}
\label{subsec:proof_full_horizon_closed_loop_equilibrium}

\begin{proof}
We show that the feedbacks and value functions stated in
Theorem~\ref{thm:full_horizon_closed_loop_equilibrium} satisfy the conditions
of the verification Theorem~\ref{thm:stitched_hjb_verification}.
Fix either the pre-announcement interval, denoted by \(k=0\), or the
post-announcement interval together with a scenario
\(k\in\{1,\ldots,m\}\). Suppress time and the fixed index \(k\), and write
\[
V_n^i(z^i)
=\frac12(z^i)^\top H_nz^i+(r_n^i)^\top z^i+c_n^i.
\]
For \(k=0\), \(r_n^i=r_n^{i,0}\) and \(c_n^i=c_n^{i,0}\). For \(k\geq1\),
\(r_n^i=r_n^k\) and \(c_n^i=c_n^k\), so these coefficients are independent
of \(i\). Also write \(g:=g^k\), where \(g^0\equiv0\), and
\(U_n^{-i}:=U_n^{-i,k}\). We first perform the HJB calculation for generic
\(r_n^i,c_n^i,g\) and specialize afterward. For \(k\geq1\), the calculation
applies on each interval between jump times of \(g^k\). No derivative of
\(g^k\) appears, so the coefficient solutions can be pasted continuously
across these times.

Set \(q_n^i:=q_n(r_n^i,g)\). Substituting
\(\nabla V_n^i=H_nz^i+r_n^i\) in \eqref{eq:stitched_hamiltonian} yields:
\begin{equation}\label{eq:vectorized_quadratic_hamiltonian}
\begin{aligned}
&\widehat{\mathscr H}_n^{i,k}
\left(t,x_n,\iota,v,H_nz^i+r_n^i;\phi_n^{-i,k}\right)
\\
&\quad={}\frac12(z^i)^\top Q_zz^i
+\left(A_zz^i+B_n^{-}U_n^{-i}+B_gg\right)^\top(H_nz^i+r_n^i)
\\
&\quad+v^\top\Lambda v
+\left(\Theta_n(H_n)z^i+q_n^i
+\vartheta_n\Lambda U_n^{-i}\right)^\top v.
\end{aligned}
\end{equation}
Its Hessian in \(v\) is \(2\Lambda\succ0\). Hence the simultaneous
first-order conditions are
\begin{equation}\label{eq:vectorized_quadratic_foc}
2\Lambda\phi_n^i+\Theta_n(H_n)z^i+q_n^i
+\vartheta_n\Lambda U_n^{-i}=0,
\qquad i=1,\ldots,n.
\end{equation}
Let \(U_n:=\sum_{j=1}^n\phi_n^j\). Since
\(U_n^{-i}=U_n-\phi_n^i\), solving
\eqref{eq:vectorized_quadratic_foc} for \(\phi_n^i\) gives
\begin{equation}\label{eq:vectorized_quadratic_rate_given_aggregate}
\phi_n^i
=-\frac1{a_n}\Lambda^{-1}
\left(\Theta_n(H_n)z^i+q_n^i+\vartheta_n\Lambda U_n\right).
\end{equation}
Summing this equation over \(i\) gives
\[
U_n
=-\frac1{a_n}\Lambda^{-1}
\left(
\Theta_n(H_n)\sum_{j=1}^nz^j+\sum_{j=1}^nq_n^j
\right)
-\frac{n\vartheta_n}{a_n}U_n.
\]
Using \(d_n=a_n+n\vartheta_n\) and solving for \(U_n\) gives
\begin{equation}\label{eq:vectorized_quadratic_aggregate_rate}
U_n
=-\frac1{d_n}\Lambda^{-1}
\left(
\Theta_n(H_n)\sum_{j=1}^nz^j+
\sum_{j=1}^nq_n^j
\right).
\end{equation}
Since \(\sum_{j=1}^n\delta^j=0\), the definitions of \(D_n\) and \(F_n\)
give
\begin{equation}\label{eq:vectorized_quadratic_state_identities}
\frac1{a_n}z^i-\frac{\vartheta_n}{a_nd_n}\sum_{j=1}^nz^j=D_nz^i,
\qquad
\frac1{a_n}z^i-\frac2{a_nd_n}\sum_{j=1}^nz^j=F_nz^i.
\end{equation}
Substituting \eqref{eq:vectorized_quadratic_aggregate_rate} into
\eqref{eq:vectorized_quadratic_rate_given_aggregate}, applying the first identity in
\eqref{eq:vectorized_quadratic_state_identities}, and then using
\(U_n^{-i}=U_n-\phi_n^i\) with the second identity in
\eqref{eq:vectorized_quadratic_state_identities} gives
\begin{subequations}\label{eq:vectorized_quadratic_feedbacks}
\begin{align}
\phi_n^i
&=-\Lambda^{-1}\left[
\Theta_n(H_n)D_nz^i+\frac1{a_n}q_n^i
-\frac{\vartheta_n}{a_nd_n}\sum_{j=1}^nq_n^j
\right],
\label{eq:vectorized_quadratic_individual_feedback}
\\
U_n^{-i}
&=\Lambda^{-1}\left[
\Theta_n(H_n)F_nz^i+\frac1{a_n}q_n^i
-\frac2{a_nd_n}\sum_{j=1}^nq_n^j
\right].
\label{eq:vectorized_quadratic_opponent_feedback}
\end{align}
\end{subequations}
Substituting \eqref{eq:vectorized_quadratic_feedbacks} into the HJB and using
\eqref{eq:vectorized_quadratic_foc} gives
\begin{equation}\label{eq:vectorized_quadratic_hjb_identity}
\begin{aligned}
0={}&\frac12(z^i)^\top\dot H_nz^i+(z^i)^\top\dot r_n^i+\dot c_n^i
+\frac12(z^i)^\top Q_zz^i
\\
&+\left(A_zz^i+B_n^{-}U_n^{-i}+B_gg\right)^\top(H_nz^i+r_n^i)
-(\phi_n^i)^\top\Lambda\phi_n^i,
\end{aligned}
\end{equation}
Grouping the terms that are quadratic in \(z^i\) gives the Riccati equation
\eqref{eq:fixed_interval_riccati_P}.

After the announcement \(r_n^i=r_n^k\), \(c_n^i=c_n^k\), and \(g=g^k\)
for every trader. The individual formula in
\eqref{eq:vectorized_quadratic_individual_feedback} becomes \(\phi_n^{i,k}\) in
\eqref{eq:full_horizon_markov_feedback}, while
\[
U_n^{-i,k}
=\Lambda^{-1}\left[
\Theta_n(H_n)F_nz^i-\frac{n-1}{d_n}q_n(r_n^k,g^k)\right].
\]
Matching the linear and constant coefficients in
\eqref{eq:vectorized_quadratic_hjb_identity}, expanding
\(q_n(r_n^k,g^k)=B_n^\top r_n^k+\chi_g\Lambda g^k\), and using
\eqref{eq:pre_ann_matrix_maps_def} gives
\eqref{eq:post_ann_common_r} and \eqref{eq:post_ann_common_alpha}.

Before the announcement, \(g=0\). Grouping the terms linear in \(z^i\) in
\eqref{eq:vectorized_quadratic_hjb_identity} and using
\eqref{eq:vectorized_quadratic_feedbacks} along with the shorthands in
\eqref{eq:pre_ann_matrix_maps_def}, gives the differential equation in
\eqref{eq:pre_ann_coupled_r_ode}. Matching the terms linear in \(z^i\) in
\eqref{eq:stitched_hjb_value_matching} gives its boundary condition. As shown by
\eqref{eq:pre_ann_common_r_ode} and
\eqref{eq:pre_ann_deviation_r_ode}, trader \(i\) can compute \(r_n^{i,0}\)
using only \((\bar\pi_n,\pi_n^i)\) by solving the ODEs for \(r_n^0\)
and \(\widetilde r_n^{i,0}\).
Equations~\eqref{eq:post_ann_q_def} and
\eqref{eq:pre_ann_affine_components} give
\[
q_n^i=B_n^\top\bigl(r_n^0+\widetilde r_n^{i,0}\bigr),
\qquad
\sum_{j=1}^nq_n^j=nB_n^\top r_n^0.
\]
Substituting these identities into
\eqref{eq:vectorized_quadratic_individual_feedback} gives the
pre-announcement feedback in
\eqref{eq:full_horizon_markov_feedback}. Grouping the constant terms in
\eqref{eq:vectorized_quadratic_hjb_identity} gives
\eqref{eq:pre_ann_constant_ode}.

The coefficient boundary conditions establish
\eqref{eq:stitched_hjb_terminal} and
\eqref{eq:stitched_hjb_value_matching}. Since \(H_n\in C^1([0,T])\), the
remaining coefficient ODEs have finite solutions that are continuous and
\(C^1\) between switching times. In every continuation game, the affine
feedbacks induce a unique linear closed-loop system with square-integrable
rates and are therefore admissible. The claim follows from
Theorem~\ref{thm:stitched_hjb_verification}.
\end{proof}

\subsection{Proof of Theorem~\ref{thm:global_riccati_solvability}}
\label{subsec:proof_global_riccati_solvability}

\begin{proof}
Fix \(n\geq2\).  The right-hand side of
\eqref{eq:fixed_interval_riccati_P} is polynomial in the entries of \(H_n\),
so there is a unique local symmetric solution backward from \(H_T\).  Let
\((\underline t_n,T]\) be its maximal backward interval.  If
\(\underline t_n>0\), the continuation criterion gives
\(\|H_n(t)\|\to\infty\) as
\(t\downarrow\underline t_n\). Hence it suffices to prove the finite bound
\eqref{eq:global_riccati_two_sided_bound} on \((\underline t_n,T]\), which
rules out such blow-up.

The calculations in the proofs of
Theorems~\ref{thm:stitched_hjb_verification} and
\ref{thm:full_horizon_closed_loop_equilibrium} show that, whenever \(H_n\)
exists on \([t,T]\), \(\frac12z^\top H_nz\) is the equilibrium value
function of the \(n\)-player game without indexer trading, obtained by setting
\(g^k\equiv0\) for every \(k\). In this game the announcement scenarios are
identical, so no stitching is required and the affine and constant
coefficients vanish. Since \(H_n\) is bounded on \([t,T]\), the induced
feedbacks are admissible. Theorem~\ref{thm:stitched_hjb_verification} then
yields
\[
V_n^i(t,x_n,\iota)=\frac12z^\top H_n(t)z,
\qquad z=(\delta,\bar x,\iota)^\top.
\]

Fix \(t\in(\underline t_n,T]\). Trader \(i\) may
deviate by choosing \(u_n^i\equiv0\) on \([t,T]\). Its inventory stays at
\(\bar x+\delta\). Because \(u_n^i\equiv0\), the execution-price term
vanishes, irrespective of the opponents' controls and the impact state. Hence
\[
V_n^i(t,x_n,\iota)
\leq\frac12(\bar x+\delta)^\top G_Q(t)(\bar x+\delta)
\]
for every \(z\), which implies
\begin{equation}
\label{eq:uniform_solvability_upper_bound_proof}
H_n(t)\preceq\overline H(t).
\end{equation}

For the lower bound, fix \(\delta,\bar x,\iota\in\mathbb R^d\) and choose
\[
X_n^i(t)=\bar x+\delta,
\qquad
X_n^j(t)=\bar x-\frac{\delta}{n-1},
\quad j\ne i
\]
so that the average of the \(X^j_n(t)\) is \(\bar x\).  Let \(V_n^j(t)\) be trader \(j\)'s equilibrium
value from this profile and impact state \(\iota\).  Summing the objectives,
dropping the nonnegative inventory penalties, and using
\(\sum_{i\ne j}u_n^i=U_n-u_n^j\) gives
\[
\begin{aligned}
\sum_{j=1}^nV_n^j(t)
\geq{}&\int_t^T\Bigg[
I(s)^\top U_n(s)
+\sum_{j=1}^n(u_n^j(s))^\top\Lambda u_n^j(s)
\\
&\qquad
+\vartheta_n\sum_{j=1}^n
(u_n^j(s))^\top\Lambda(U_n(s)-u_n^j(s))
\Bigg]ds.
\end{aligned}
\]
The execution-cost terms satisfy
\[
\begin{aligned}
\sum_j(u_n^j)^\top\Lambda
\left(u_n^j+\vartheta_n(U_n-u_n^j)\right)
&=\vartheta_nU_n^\top\Lambda U_n
+(1-\vartheta_n)\sum_j(u_n^j)^\top\Lambda u_n^j
\\
&\geq
\left(\vartheta_n+\frac{1-\vartheta_n}{n}\right)
U_n^\top\Lambda U_n.
\end{aligned}
\]
The inequality uses \(0\leq\vartheta_n\leq1\) from
Assumption~\ref{ass:standing_model} and
\(\sum_j(u_n^j)^\top\Lambda u_n^j\geq n^{-1}U_n^\top\Lambda U_n\).  Set
\(\alpha_n:=\vartheta_n+(1-\vartheta_n)/n\).  Then
\begin{equation}
\label{eq:uniform_solvability_aggregate_execution_bound}
\sum_{j=1}^nV_n^j(t)
\geq
\int_t^T\left(
I(s)^\top U_n(s)+\alpha_nU_n(s)^\top\Lambda U_n(s)
\right)ds.
\end{equation}

Because \(g\equiv0\),
\[
I(s)=e^{-R(s-t)}\iota
+\varpi_n\int_t^se^{-R(s-t')}\Gamma U_n(t')dt'.
\]
Therefore
\begin{equation}\label{eq:uniform_solvability_impact_cost_decomposition}
\begin{aligned}
\int_t^TI(s)^\top U_n(s)ds
={}&\int_t^TU_n(s)^\top e^{-R(s-t)}\iota ds
\\
&+\varpi_n\int_t^TU_n(s)^\top
\left(\int_t^se^{-R(s-t')}\Gamma U_n(t')dt'\right)ds.
\end{aligned}
\end{equation}
The second term on the right-hand side of
\eqref{eq:uniform_solvability_impact_cost_decomposition} is nonnegative because
\(\varpi_n>0\) and \(\mathcal C(v)\geq0\) under
Assumption~\ref{ass:standing_model}. Substituting
\eqref{eq:uniform_solvability_impact_cost_decomposition} into
\eqref{eq:uniform_solvability_aggregate_execution_bound}, dropping the
nonnegative second term, and completing the square pointwise in \(U_n(s)\)
gives
\begin{equation}\label{eq:uniform_solvability_aggregate_value_bound}
\begin{aligned}
\sum_{j=1}^nV_n^j(t)
&\geq
\int_t^T\left[
U_n(s)^\top e^{-R(s-t)}\iota
+\alpha_nU_n(s)^\top\Lambda U_n(s)
\right]ds
\\
&\geq
-\frac1{4\alpha_n}\int_t^T
\iota^\top e^{-R^\top(s-t)}\Lambda^{-1}e^{-R(s-t)}\iota ds
\\
&=-\frac1{4\alpha_n}\iota^\top G_R(t)\iota.
\end{aligned}
\end{equation}

Let \(P_n(t):=[H_n(t)]_{\delta\delta}\), and partition \(H_n(t)\) according
to the deviation coordinate \(\delta\) and the common-state coordinates
\((\bar x,\iota)\):
\[
H_n(t)=
\begin{bmatrix}
P_n(t)&[H_n(t)]_{\delta,(\bar x,\iota)}\\
[H_n(t)]_{(\bar x,\iota),\delta}
&[H_n(t)]_{(\bar x,\iota),(\bar x,\iota)}
\end{bmatrix}.
\]
For the inventory profile above,
\(z^i=(\delta,\bar x,\iota)^\top\) and
\(z^j=(-\delta/(n-1),\bar x,\iota)^\top\) for \(j\ne i\).  Since the
deviations sum to zero, the cross terms cancel:
\begin{equation}
\label{eq:uniform_solvability_aggregate_block_identity}
\begin{aligned}
\sum_{j=1}^nV_n^j(t)
={}&\frac{n}{2(n-1)}\delta^\top P_n(t)\delta
+\frac n2
\begin{bmatrix}\bar x\\\iota\end{bmatrix}^{\!\top}
[H_n(t)]_{(\bar x,\iota),(\bar x,\iota)}
\begin{bmatrix}\bar x\\\iota\end{bmatrix}.
\end{aligned}
\end{equation}
Setting \(\bar x=\iota=0\) in
\eqref{eq:uniform_solvability_aggregate_value_bound} and using
\eqref{eq:uniform_solvability_aggregate_block_identity} gives, for every
\(\delta\in\mathbb R^d\),
\[
0\leq\sum_{j=1}^nV_n^j(t)
=\frac{n}{2(n-1)}\delta^\top P_n(t)\delta.
\]
Hence
\begin{equation}
\label{eq:uniform_solvability_deviation_block_psd}
P_n(t)\succeq0.
\end{equation}
Setting \(\delta=0\) in
\eqref{eq:uniform_solvability_aggregate_value_bound} and using
\eqref{eq:uniform_solvability_aggregate_block_identity} gives, for every
\((\bar x,\iota)\in\mathbb R^{2d}\),
\[
\frac n2
\begin{bmatrix}\bar x\\\iota\end{bmatrix}^{\!\top}
[H_n(t)]_{(\bar x,\iota),(\bar x,\iota)}
\begin{bmatrix}\bar x\\\iota\end{bmatrix}
\geq-\frac1{4\alpha_n}\iota^\top G_R(t)\iota.
\]
This gives the first matrix inequality below, while
\(n\alpha_n=1+(n-1)\vartheta_n\geq1\) and \(G_R(t)\succeq0\) give the second:
\begin{equation}\label{eq:uniform_solvability_common_block_lower_bound}
\bigl[H_n(t)\bigr]_{(\bar x,\iota),(\bar x,\iota)}
\succeq
-\frac1{2n\alpha_n}
\begin{bmatrix}0&0\\0&G_R(t)\end{bmatrix}
\succeq
-\frac12\begin{bmatrix}0&0\\0&G_R(t)\end{bmatrix}.
\end{equation}

Taking traces in \eqref{eq:uniform_solvability_deviation_block_psd} and
\eqref{eq:uniform_solvability_common_block_lower_bound} gives
\begin{equation}\label{eq:uniform_solvability_trace_lower_bound}
\operatorname{tr}(H_n(t))
=\operatorname{tr}(P_n(t))
+\operatorname{tr}\!\left(
\bigl[H_n(t)\bigr]_{(\bar x,\iota),(\bar x,\iota)}
\right)
\geq-\frac12\operatorname{tr}(G_R(t)).
\end{equation}
Set \(W_n(t):=\overline H(t)-H_n(t)\). By
\eqref{eq:uniform_solvability_upper_bound_proof}, \(W_n(t)\succeq0\). Since
\(\operatorname{tr}(\overline H(t))=2\operatorname{tr}(G_Q(t))\),
\eqref{eq:uniform_solvability_trace_lower_bound} gives
\begin{equation}\label{eq:uniform_solvability_w_trace_upper_bound}
\operatorname{tr}(W_n(t))
=\operatorname{tr}(\overline H(t))-\operatorname{tr}(H_n(t))
\leq
2\operatorname{tr}(G_Q(t))
+\frac12\operatorname{tr}(G_R(t)).
\end{equation}
Positive semidefiniteness of \(W_n(t)\) and
\eqref{eq:uniform_solvability_w_trace_upper_bound} imply
\[
W_n(t)
\preceq\operatorname{tr}(W_n(t))\operatorname{Id}_{3d}
\preceq
\left(
2\operatorname{tr}(G_Q(t))
+\frac12\operatorname{tr}(G_R(t))
\right)\operatorname{Id}_{3d}.
\]
Since \(\overline H(t)\succeq0\),
\[
\begin{aligned}
H_n(t)
&=\overline H(t)-W_n(t)
\succeq-W_n(t)
\\
&\succeq
-\left(
2\operatorname{tr}(G_Q(t))
+\frac12\operatorname{tr}(G_R(t))
\right)\operatorname{Id}_{3d}
=-\underline H(t).
\end{aligned}
\]
Together with \eqref{eq:uniform_solvability_upper_bound_proof}, this proves
\eqref{eq:global_riccati_two_sided_bound} on the maximal interval.

The functions \(G_Q\) and \(G_R\) are continuous on \([0,T]\), so the bound
is finite there and uniform over \(n\geq2\).  A finite \(\underline t_n\)
would contradict the continuation criterion.  Hence \(\underline t_n=0\),
and \(H_n\) exists uniquely on \([0,T]\).

Finally, boundedness of \(H_n\) makes the equations for \(r_n^k\),
\(r_n^0\), and \(\widetilde r_n^{i,0}\) linear terminal-value problems
with square-integrable forcing. They have unique finite solutions. The
equations for \(c_n^k\) and \(c_n^{i,0}\) then have integrable right-hand
sides and likewise have unique finite solutions. The construction in
Theorem~\ref{thm:full_horizon_closed_loop_equilibrium} is therefore well
defined.
\end{proof}

\section{Proofs of Section~\ref{sec:mean_field_limit}}
\label{sec:proofs_mean_field_limit}

\subsection{Proof of Theorem~\ref{thm:mf_solution}}
\label{subsec:proof_mf_solution}

\begin{proof}
Fix a candidate mean flow and its impact components from
\eqref{eq:mf_component_impact_dynamics}. We first solve the type-\(\pi\)
control problem \eqref{eq:mf_type_objective} and then verify the consistency
condition \eqref{eq:mf_average_consistency}.

In a post-announcement scenario \(k\ge1\), the player's HJB equation on
\([T_{\mathrm{ann}},T]\) is
\[
\begin{aligned}
0
={}&\frac{\partial}{\partial t}V^k(t,x)
+\frac12 x^\top Q x
\\
&+\inf_{u\in\mathbb R^d}
\Bigl\{
u^\top\Lambda u
+\bigl(I^k(t)+\Lambda(\varpi\chi_u\bar u^k(t)+\chi_g g^k(t))
+V_x^k(t,x)\bigr)^\top u
\Bigr\},
\end{aligned}
\]
with terminal condition \(V^k(T,x)=\frac12x^\top Q_T x\). Insert the
post-announcement ansatz in \eqref{eq:mf_quadratic_value_ansatz}.
The first-order condition is
\[
2\Lambda u+I^k(t)+\Lambda(\varpi\chi_u\bar u^k(t)+\chi_g g^k(t))
+P(t)x+\ell^k(t)=0.
\]
Solving this for \(u\) gives the optimal feedback
\begin{equation}\label{eq:mf_post_feedback_derived}
\phi^{k}(t,x)
=
-\frac12\Lambda^{-1}
\left(P(t)x+\ell^k(t)+I^k(t)+
\Lambda(\varpi\chi_u\bar u^k(t)+\chi_g g^k(t))\right).
\end{equation}
Substituting the minimizer back into the HJB and collecting the quadratic,
linear, and constant terms gives
\begin{subequations}\label{eq:mf_post_value_system_derived}
\begin{align}
\dot P(t)
&=
- Q+\frac12P(t)\Lambda^{-1}P(t),
\label{eq:mf_post_riccati_derived}\\
\dot \ell^k(t)
&=
\frac12P(t)\Lambda^{-1}
\left(I^k(t)+\Lambda(\varpi\chi_u\bar u^k(t)+\chi_g g^k(t))+\ell^k(t)\right),
\label{eq:mf_post_ell_derived}\\
\dot c^k(t)
&=
\frac14
\left(I^k(t)+\Lambda(\varpi\chi_u\bar u^k(t)+\chi_g g^k(t))+\ell^k(t)\right)^\top
\\
&\qquad
{}\Lambda^{-1}
\left(I^k(t)+\Lambda(\varpi\chi_u\bar u^k(t)+\chi_g g^k(t))+\ell^k(t)\right)
\label{eq:mf_post_alpha_derived}
\end{align}
\end{subequations}
with terminal conditions
\begin{subequations}\label{eq:mf_post_terminal_conditions_derived}
\begin{align}
P(T)&=Q_T,
\label{eq:mf_P_terminal_derived}\\
\ell^k(T)&=0,
\label{eq:mf_post_ell_terminal_derived}\\
c^k(T)&=0.
\label{eq:mf_post_alpha_terminal_derived}
\end{align}
\end{subequations}

Before the announcement, the type-\(\pi\) player has continuation value
\[
\begin{aligned}
V^{\pi,0}(T_{\mathrm{ann}},x)
&=
\sum_{k=1}^m\pi^k V^k(T_{\mathrm{ann}},x)\\
&=
\frac12x^\top P(T_{\mathrm{ann}})x
+\left(\sum_{k=1}^m\pi^k\ell^k(T_{\mathrm{ann}})\right)^\top x
+\sum_{k=1}^m\pi^kc^k(T_{\mathrm{ann}}),
\end{aligned}
\]
so the pre-announcement boundary conditions are
\begin{subequations}\label{eq:mf_pre_boundary_derived}
\begin{align}
\ell^{\pi,0}(T_{\mathrm{ann}})
&=
\sum_{k=1}^m\pi^k\ell^k(T_{\mathrm{ann}}),
\label{eq:mf_pre_ell_boundary_derived}\\
c^{\pi,0}(T_{\mathrm{ann}})
&=
\sum_{k=1}^m\pi^kc^k(T_{\mathrm{ann}}).
\label{eq:mf_pre_alpha_boundary_derived}
\end{align}
\end{subequations}
The same HJB calculation applies before the announcement. Its quadratic
equation is again \eqref{eq:mf_post_riccati_derived}, so the ODE for \(P\) is independent of
\(\pi\) and common to the pre- and post-announcement problems. Thus it can be
solved once on \([0,T]\). The type-dependent affine and scalar equations on the
pre-announcement interval are
\begin{subequations}\label{eq:mf_pre_value_system_derived}
\begin{align}
\dot\ell^{\pi,0}(t)
&=
\frac12P(t)\Lambda^{-1}
\left(I^0(t)+\Lambda\varpi\chi_u\bar u^0(t)+\ell^{\pi,0}(t)\right),
\label{eq:mf_pre_ell_derived}\\
\dot c^{\pi,0}(t)
&=
\frac14
\left(I^0(t)+\Lambda\varpi\chi_u\bar u^0(t)+\ell^{\pi,0}(t)\right)^\top
\Lambda^{-1}
\left(I^0(t)+\Lambda\varpi\chi_u\bar u^0(t)+\ell^{\pi,0}(t)\right)
\label{eq:mf_pre_alpha_derived}
\end{align}
\end{subequations}
with the boundary conditions in \eqref{eq:mf_pre_boundary_derived}. The
first-order condition gives the pre-announcement feedback
\begin{equation}\label{eq:mf_pre_feedback_derived}
\phi^{\pi,0}(t,x)
=
-\frac12\Lambda^{-1}
\left(P(t)x+\ell^{\pi,0}(t)+I^0(t)+\Lambda\varpi\chi_u\bar u^0(t)\right).
\end{equation}

We now impose the mean-field consistency condition
\eqref{eq:mf_average_consistency}. Averaging
\eqref{eq:mf_pre_feedback_derived} and
\eqref{eq:mf_post_feedback_derived} over types and solving for the mean flows
gives
\begin{subequations}\label{eq:mf_mu_closed_form_derived}
\begin{align}
\bar u^0
&=-\frac{1}{2+\varpi\chi_u}\Lambda^{-1}
\left(P\bar X^0+I^0+\ell^0\right),
\label{eq:mf_pre_mu_closed_form_derived}\\
\bar u^k
&=-\frac{1}{2+\varpi\chi_u}\Lambda^{-1}
\left(P\bar X^k+I^k+\ell^k+\Lambda\chi_g g^k\right),
\qquad k=1,\ldots,m.
\label{eq:mf_post_mu_closed_form_derived}
\end{align}
\end{subequations}
Averaging the inventory dynamics gives
\begin{subequations}\label{eq:mf_mean_inventory_dynamics_derived}
\begin{align}
\dot{\bar X}^0(t)&=\bar u^0(t),
&\bar X^0(0)&=0,
\label{eq:mf_pre_mean_inventory_dynamics_derived}\\
\dot{\bar X}^k(t)&=\bar u^k(t),
&\bar X^k(T_{\mathrm{ann}})&=\bar X^0(T_{\mathrm{ann}}),
\quad k=1,\ldots,m.
\label{eq:mf_post_mean_inventory_dynamics_derived}
\end{align}
\end{subequations}
Averaging \eqref{eq:mf_pre_ell_derived} and
\eqref{eq:mf_pre_ell_boundary_derived}, and using
\(\bar\pi=\int_{\Delta_m}\pi\,\Pi(d\pi)\), gives
\eqref{eq:mf_pre_common_affine_derived} and
\eqref{eq:mf_pre_common_affine_boundary_derived}:
\begin{subequations}\label{eq:mf_pre_common_affine_system_derived}
\begin{align}
\dot \ell^0(t)
&=
\frac12P(t)\Lambda^{-1}
\left(I^0(t)+\Lambda\varpi\chi_u\bar u^0(t)+\ell^0(t)\right),
\label{eq:mf_pre_common_affine_derived}\\
\ell^0(T_{\mathrm{ann}})
&=
\sum_{k=1}^m\bar\pi^k\ell^k(T_{\mathrm{ann}}),
\label{eq:mf_pre_common_affine_boundary_derived}
\end{align}
\end{subequations}
Substituting
\eqref{eq:mf_mu_closed_form_derived} into
\eqref{eq:mf_mean_inventory_dynamics_derived},
\eqref{eq:mf_component_impact_dynamics},
\eqref{eq:mf_pre_common_affine_system_derived}, and
\eqref{eq:mf_post_ell_derived}, and collecting terms,
gives the linear ODE in \eqref{eq:mf_y_system} with \(\mathcal A\) and \(b^k\) defined in
\eqref{eq:mf_A_matrix} and \eqref{eq:mf_b_vector}.
The boundary conditions in
\eqref{eq:mf_mean_inventory_dynamics_derived},
\eqref{eq:mf_component_impact_dynamics},
\eqref{eq:mf_pre_common_affine_boundary_derived}, and
\eqref{eq:mf_post_ell_terminal_derived} establish the boundary conditions of the ODE in
\eqref{eq:mf_consistency_system}.

The equilibrium feedbacks induce rates that are square-integrable uniformly
in \(\pi\), and their type averages \(\bar u^0,\ldots,\bar u^m\) are therefore
square-integrable. Together with the optimality established by the HJB
calculation and the mean-field consistency established by averaging the
feedbacks, this proves that the feedback family satisfies
Definition~\ref{def:mf_equilibrium}.
\end{proof}

\subsection{Proof of Lemma~\ref{lem:mf_state_transition_representation}}
\label{subsec:proof_mf_state_transition_representation}

\begin{proof}
For the linear ODE \(\dot Y(t)=\mathcal A(t)Y(t)+b(t)\), the solution with
initial value \(Y(s)\) is
\begin{equation}\label{eq:mf_linear_solution_formula}
Y(t)
=
\mathcal E(t,s)Y(s)
+
\int_s^t\mathcal E(t,s')b(s')\,ds'.
\end{equation}
On the pre-announcement interval, \(b^0=0\) and
\(\bar X^0(0)=I^0(0)=0\). Hence the whole pre-announcement path is determined by
the single unknown \(\ell^0(0)\). Applying the formula with \(s=0\)
gives the
pre-announcement path and, at \(T_{\mathrm{ann}}\),
\[
\bar X^0(T_{\mathrm{ann}})
=
\mathcal E^{\mathrm{pre}}_{X\ell}\ell^0(0),
\qquad
I^0(T_{\mathrm{ann}})
=
\mathcal E^{\mathrm{pre}}_{I\ell}\ell^0(0),
\qquad
\ell^0(T_{\mathrm{ann}})
=
\mathcal E^{\mathrm{pre}}_{\ell\ell}\ell^0(0).
\]
For a fixed post-announcement scenario \(k\), applying the same formula on
\([T_{\mathrm{ann}},T]\) with initial state \((\bar X^0(T_{\mathrm{ann}}),I^0(T_{\mathrm{ann}}),\ell^k(T_{\mathrm{ann}}))\) gives the
post-announcement path. Its terminal \(\ell\)-block is
\[
\ell^k(T)
=
\mathcal E^{\mathrm{post}}_{\ell X}\bar X^0(T_{\mathrm{ann}})+
\mathcal E^{\mathrm{post}}_{\ell I}I^0(T_{\mathrm{ann}})+
\mathcal E^{\mathrm{post}}_{\ell\ell}\ell^k(T_{\mathrm{ann}})+
\mathcal D_\ell^k.
\]
The boundary condition \(\ell^k(T)=0\) therefore implies
\begin{equation}\label{eq:simplification_proof_zero_cond}
0
=
\mathcal E^{\mathrm{post}}_{\ell X}\bar X^0(T_{\mathrm{ann}})
+\mathcal E^{\mathrm{post}}_{\ell I}I^0(T_{\mathrm{ann}})
+\mathcal E^{\mathrm{post}}_{\ell\ell}\ell^k(T_{\mathrm{ann}})
+\mathcal D_\ell^k.
\end{equation}
Averaging these identities with weights \(\bar\pi^k\) and using the announcement matching condition
\(\ell^0(T_{\mathrm{ann}})=\sum_k\bar\pi^k\ell^k(T_{\mathrm{ann}})\) gives
\[
0
=
\mathcal E^{\mathrm{post}}_{\ell X}\bar X^0(T_{\mathrm{ann}})
+\mathcal E^{\mathrm{post}}_{\ell I}I^0(T_{\mathrm{ann}})
+\mathcal E^{\mathrm{post}}_{\ell\ell}\ell^0(T_{\mathrm{ann}})
+\sum_{k=1}^m\bar\pi^k\mathcal D_\ell^k.
\]
Using
\[
\bar X^0(T_{\mathrm{ann}})
=
\mathcal E^{\mathrm{pre}}_{X\ell}\ell^0(0),
\qquad
I^0(T_{\mathrm{ann}})
=
\mathcal E^{\mathrm{pre}}_{I\ell}\ell^0(0),
\qquad
\ell^0(T_{\mathrm{ann}})
=
\mathcal E^{\mathrm{pre}}_{\ell\ell}\ell^0(0),
\]
and the identity
\(\mathcal E^{\mathrm{post}}\mathcal E^{\mathrm{pre}}
=\mathcal E(T,0)\), we obtain
\begin{equation}\label{eq:mf_matching_equation}
\mathcal M\ell^0(0)
=
-\sum_{k=1}^m\bar\pi^k\mathcal D_\ell^k.
\end{equation}
Since \(\mathcal M\) is invertible, this determines \(\ell^0(0)\) and hence
the pre-announcement path. If \(\mathcal E^{\mathrm{post}}_{\ell\ell}\) is invertible, each
scenario-specific terminal condition \eqref{eq:simplification_proof_zero_cond} determines
\[
\ell^k(T_{\mathrm{ann}})
=
-(\mathcal E^{\mathrm{post}}_{\ell\ell})^{-1}
\left(\mathcal E^{\mathrm{post}}_{\ell X}\bar X^0(T_{\mathrm{ann}})+
\mathcal E^{\mathrm{post}}_{\ell I}I^0(T_{\mathrm{ann}})+\mathcal D_\ell^k\right).
\]
The solution formula \eqref{eq:mf_linear_solution_formula} then gives all
post-announcement paths. These paths satisfy the ODEs and all boundary
conditions. Uniqueness follows because the invertibility of \(\mathcal M\)
fixes \(\ell^0(0)\), the invertibility of
\(\mathcal E^{\mathrm{post}}_{\ell\ell}\) fixes each \(\ell^k(T_{\mathrm{ann}})\), and the resulting linear
initial-value problems have unique solutions.
\end{proof}

\subsection{Proof of Theorem~\ref{thm:mf_implementability}}
\label{subsec:proof_mf_implementability}

\begin{proof}
We first prove that the post-announcement block
\(\mathcal E^{\mathrm{post}}_{\ell\ell}\) is invertible. Recall from
Lemma~\ref{lem:mf_state_transition_representation} that
\(\mathcal E^{\mathrm{post}}=\mathcal E(T,T_{\mathrm{ann}})\), and that
\(\mathcal E^{\mathrm{post}}_{\ell\ell}\) denotes its \(\ell\ell\)-block. The matrix
\(\mathcal E(t,s)\) is the fundamental solution of the homogeneous ODE
\begin{equation}\label{eq:mf_homogeneous_transition_ode}
\dot{Y}(t)=\mathcal A(t)Y(t),
\qquad
Y(t)=(\bar X(t),I(t),\ell(t))^\top
\end{equation}
with \(\mathcal{A}(t)\) defined in \eqref{eq:mf_A_matrix}.
Let \(\mathfrak q\in\ker(\mathcal E^{\mathrm{post}}_{\ell\ell})\). Solve
\eqref{eq:mf_homogeneous_transition_ode} on
\([T_{\mathrm{ann}},T]\) with initial condition
\[
\bar X(T_{\mathrm{ann}})=0,
\qquad
I(T_{\mathrm{ann}})=0,
\qquad
\ell(T_{\mathrm{ann}})=\mathfrak q.
\]
By definition of \(\mathcal E^{\mathrm{post}}\) and the assumption
\(\mathfrak q\in\ker(\mathcal E^{\mathrm{post}}_{\ell\ell})\), the terminal affine coefficient is
\[
\begin{aligned}
\ell(T)
={}&
\mathcal E^{\mathrm{post}}_{\ell X}\bar X(T_{\mathrm{ann}})
+
\mathcal E^{\mathrm{post}}_{\ell I}I(T_{\mathrm{ann}})
+
\mathcal E^{\mathrm{post}}_{\ell\ell}\ell(T_{\mathrm{ann}})
=\mathcal E^{\mathrm{post}}_{\ell\ell}\mathfrak q
=0.
\end{aligned}
\]
The first row of \eqref{eq:mf_homogeneous_transition_ode} is
\begin{equation}\label{eq:mf_homogeneous_mean_inventory_ode}
\dot{\bar X}(t)
=
-\frac{1}{2+\varpi\chi_u}\Lambda^{-1}
\left(P(t)\bar X(t)+I(t)+\ell(t)\right).
\end{equation}
Denote the right-hand side of
\eqref{eq:mf_homogeneous_mean_inventory_ode} by \(\bar u(t)\) and define the
costate \(p(t):=P(t)\bar X(t)+\ell(t)\). Then
\eqref{eq:mf_homogeneous_mean_inventory_ode} can be reformulated as
\begin{equation}\label{eq:mf_homogeneous_first_row_identity}
p(t)+I(t)+(2+\varpi\chi_u)\Lambda\bar u(t)=0.
\end{equation}
The costate process \(p(t)\) satisfies
\[
\dot p(t)=- Q\bar X(t),
\qquad
p(T)=Q_T\bar X(T).
\]
Indeed, the terminal condition follows from \(P(T)=Q_T\) and \(\ell(T)=0\). The ODE
follows by differentiating \(p(t)=P(t)\bar X(t)+\ell(t)\) and using
\eqref{eq:mf_riccati} together with the first and third rows of
\eqref{eq:mf_homogeneous_transition_ode}.
The impact row of \eqref{eq:mf_homogeneous_transition_ode} gives
\begin{equation}\label{eq:mf_homogeneous_impact_row}
\dot I(t)=-RI(t)+\varpi\Gamma\bar u(t).
\end{equation}
Multiplying \eqref{eq:mf_homogeneous_first_row_identity} by
\(\bar u(t)^\top\) and
integrating over \([T_{\mathrm{ann}},T]\) gives
\begin{equation}\label{eq:mf_homogeneous_energy_identity}
0
=
\int_{T_{\mathrm{ann}}}^T\bar u(t)^\top p(t)\,dt
+
\int_{T_{\mathrm{ann}}}^T\bar u(t)^\top I(t)\,dt
+
(2+\varpi\chi_u)\int_{T_{\mathrm{ann}}}^T
\bar u(t)^\top\Lambda\bar u(t)\,dt.
\end{equation}
Integration by parts gives
\[
\begin{aligned}
\int_{T_{\mathrm{ann}}}^T\bar u(t)^\top p(t)\,dt
={}&\bar X(T)^\top p(T)
-\bar X(T_{\mathrm{ann}})^\top p(T_{\mathrm{ann}})-\int_{T_{\mathrm{ann}}}^T\bar X(t)^\top\dot p(t)\,dt\\
={}&\bar X(T)^\top Q_T\bar X(T)
+\int_{T_{\mathrm{ann}}}^T\bar X(t)^\top Q\bar X(t)\,dt.
\end{aligned}
\]
All three summands on the right-hand side of
\eqref{eq:mf_homogeneous_energy_identity} are nonnegative. The first is
nonnegative by the integration-by-parts identity since
\(Q,Q_T\succeq0\). The second is nonnegative because
\eqref{eq:mf_homogeneous_impact_row}, together with
\(I(T_{\mathrm{ann}})=0\), identifies it with \(\varpi\) times the transient-impact quadratic
form of \(\bar u\) on \([T_{\mathrm{ann}},T]\), which is nonnegative by
\(\varpi>0\) and the
kernel-positivity requirement in Assumption~\ref{ass:standing_model}. The third
is nonnegative by \(2+\varpi\chi_u>0\) and \(\Lambda\succ0\). Hence
\eqref{eq:mf_homogeneous_energy_identity} implies \(\bar u\equiv0\). Since
\(\bar X(T_{\mathrm{ann}})=0\) and
\(\dot{\bar X}(t)=\bar u(t)\), this implies \(\bar X\equiv0\). The impact row
\eqref{eq:mf_homogeneous_impact_row} with \(I(T_{\mathrm{ann}})=0\) then
gives \(I\equiv0\). The identity \(p+I+(2+\varpi\chi_u)\Lambda\bar u=0\) from
\eqref{eq:mf_homogeneous_first_row_identity} gives
\(p\equiv0\). Evaluating at \(T_{\mathrm{ann}}\) therefore yields
\[
0=p(T_{\mathrm{ann}})
=P(T_{\mathrm{ann}})\bar X(T_{\mathrm{ann}})+\ell(T_{\mathrm{ann}})
=\ell(T_{\mathrm{ann}})=\mathfrak q
\]
and so \(\mathfrak q=0\). Hence
\(\ker(\mathcal E^{\mathrm{post}}_{\ell\ell})=\{0\}\) and
\(\mathcal E^{\mathrm{post}}_{\ell\ell}\) is invertible.

It remains to prove that the matching matrix \(\mathcal M\)
from \eqref{eq:mf_matching_matrix} is invertible.
Let \(\mathfrak q\in\ker(\mathcal M)\). Run
\eqref{eq:mf_homogeneous_transition_ode} on \([0,T_{\mathrm{ann}}]\) from
\[
\bar X(0)=0,
\qquad
I(0)=0,
\qquad
\ell(0)=\mathfrak q.
\]
By definition of \(\mathcal E^{\mathrm{pre}}\), the state at \(T_{\mathrm{ann}}\) is
\[
\bar X(T_{\mathrm{ann}})=\mathcal E^{\mathrm{pre}}_{X\ell}\mathfrak q,
\qquad
I(T_{\mathrm{ann}})=\mathcal E^{\mathrm{pre}}_{I\ell}\mathfrak q,
\qquad
\ell(T_{\mathrm{ann}})=\mathcal E^{\mathrm{pre}}_{\ell\ell}\mathfrak q.
\]
Using
\(\mathcal E(T,0)=\mathcal E(T,T_{\mathrm{ann}})
\mathcal E(T_{\mathrm{ann}},0)\) and
\eqref{eq:mf_matching_matrix}, the assumption
\(\mathfrak q\in\ker(\mathcal M)\) gives
\[
\mathcal E^{\mathrm{post}}_{\ell X}\bar X(T_{\mathrm{ann}})
+
\mathcal E^{\mathrm{post}}_{\ell I}I(T_{\mathrm{ann}})
+
\mathcal E^{\mathrm{post}}_{\ell\ell}\ell(T_{\mathrm{ann}})
=\mathcal M\mathfrak q
=0.
\]
By definition of \(\mathcal E^{\mathrm{post}}\), this implies that the solution to
\eqref{eq:mf_homogeneous_transition_ode} on
\([T_{\mathrm{ann}},T]\), starting from
\((\bar X(T_{\mathrm{ann}}),I(T_{\mathrm{ann}}),\ell(T_{\mathrm{ann}}))\), has terminal
affine coefficient \(\ell(T)=0\). Equivalently, running the homogeneous ODE over
the combined interval \([0,T]\) with \(\bar X(0)=0\), \(I(0)=0\), and
\(\ell(0)=\mathfrak q\)
produces a solution satisfying \(\ell(T)=0\). The same energy argument used for
\(\mathcal E^{\mathrm{post}}_{\ell\ell}\), now on \([0,T]\), gives
\(\bar u\equiv0\), hence \(\bar X\equiv0\), \(I\equiv0\), and
\(p\equiv0\).
Since \(p\equiv0\), \(\bar X(0)=0\), and \(\ell(0)=\mathfrak q\), evaluating
\(p=P\bar X+\ell\) at \(t=0\) gives
\(0=p(0)=\mathfrak q\). Thus
\(\ker(\mathcal M)=\{0\}\), and \(\mathcal M\) is invertible.
\end{proof}

\subsection{Proof of Theorem~\ref{thm:finite_to_mf_convergence}}
\label{subsec:proof_finite_to_mf_convergence}

\begin{proof}
The proof has four steps.

\emph{Step 1: Finite Riccati Coefficients.}
Recall that \(H_n:[0,T]\to\mathbb S^{3d}\) is the \(3d\times3d\) finite-player
Riccati coefficient from \eqref{eq:fixed_interval_riccati_P}. In contrast, the
mean-field Riccati coefficient \(P:[0,T]\to\mathbb S^d\) from
\eqref{eq:mf_riccati} is \(d\times d\). Thus \(H_n\) and \(P\) cannot be
compared as whole matrices. The goal of this step is to show that \(H_n\)
converges to the solution \(H:[0,T]\to\mathbb S^{3d}\) of the limiting
three-block equation, and then to identify
\([H]_{\delta\delta}\) with \(P\).

The Riccati equation for \(H_n\) can be written as
\begin{equation}\label{eq:finite_riccati_vector_field_ode}
0=\dot H_n(t)+\mathscr F_n(H_n(t)),
\qquad
H_n(T)=H_T
\end{equation}
where the vector field \(\mathscr F_n(H_0)\) is given by
\[
\mathscr F_n(H_0)
:=
Q_z
+M_n(H_0)^\top H_0
+H_0 M_n(H_0)
-2D_n^\top
\Theta_n(H_0)^\top\Lambda^{-1}\Theta_n(H_0)D_n.
\]

We now take the limit of \(\mathscr F_n\) term by term. All estimates below are uniform
for \(H_0\) in bounded subsets of \(\mathbb S^{3d}\). First, from
\eqref{eq:fixed_interval_D_def},
\[
D_n
=
\operatorname{diag}
\left(
\frac1{2-\varpi\chi_u/n}\operatorname{Id}_d,
\frac1{2+\varpi\chi_u(n-1)/n}\operatorname{Id}_d,
\frac1{2+\varpi\chi_u(n-1)/n}\operatorname{Id}_d
\right).
\]
Hence
\[
D_n
=
D+O\!\left(\frac1n\right),
\qquad
D:=
\operatorname{diag}
\left(\frac12\operatorname{Id}_d,\frac{1}{2+\varpi\chi_u}\operatorname{Id}_d,
\frac{1}{2+\varpi\chi_u}\operatorname{Id}_d\right).
\]
Next, from \eqref{eq:fixed_interval_matrix_maps_def},
\[
\Theta_n(H_0)
=
\Theta(H_0)+O\!\left(\frac1n\right),
\qquad
\Theta(H_0):=E_\iota+B_\delta^\top H_0,
\qquad
B_\delta:=
\begin{bmatrix}
\operatorname{Id}_d\\0\\0
\end{bmatrix}.
\]
Finally, from \eqref{eq:fixed_interval_matrix_maps_def},
\[
M_n(H_0)
=
M(H_0)+O\!\left(\frac1n\right),
\]
where
\[
\begin{aligned}
M(H_0)&:=
A_z+\frac{1}{2+\varpi\chi_u} B_{\mathrm{com}}\Lambda^{-1}\Theta(H_0)E_{\mathrm{com}},
\\
B_{\mathrm{com}}&:=
\begin{bmatrix}
\operatorname{Id}_d\\-\operatorname{Id}_d\\-\varpi\Gamma
\end{bmatrix},
\\
E_{\mathrm{com}}&:=
\operatorname{diag}(0,\operatorname{Id}_d,\operatorname{Id}_d).
\end{aligned}
\]
Indeed,
\[
B_n^{-}=-\frac1n B_{\mathrm{com}},
\qquad
\frac1n F_n
=
-\frac{1}{2+\varpi\chi_u} E_{\mathrm{com}}+O\!\left(\frac1n\right),
\]
so the order-\(n\) \(\bar x\)- and \(\iota\)-blocks of \(F_n\) offset the
\(1/n\)-scale of \(B_n^{-}\).
Combining these three estimates gives, for every compact
\(\mathfrak C\subset\mathbb S^{3d}\), a constant
\(C_{\mathfrak C}<\infty\), independent of \(n\), such that
\begin{equation}\label{eq:finite_riccati_vector_field_convergence}
\sup_{H_0\in\mathfrak C}
\|\mathscr F_n(H_0)-\mathscr F(H_0)\|
\le
\frac{C_{\mathfrak C}}{n},
\end{equation}
where \(\|\cdot\|\) is any fixed matrix norm.
Here,
\[
\mathscr F(H_0)
:=
Q_z
+M(H_0)^\top H_0
+H_0 M(H_0)
-2D^\top
\Theta(H_0)^\top\Lambda^{-1}
\Theta(H_0)D.
\]

Since \(\mathscr F\) is locally Lipschitz, let \(H\) be the unique
maximal backward solution of
\begin{equation}\label{eq:limiting_three_block_riccati}
0=\dot{H}(t)+\mathscr F(H(t)),
\qquad
H(T)=H_T.
\end{equation}
Write its maximal interval as \((\underline t,T]\).
Theorem~\ref{thm:global_riccati_solvability} gives a constant \(C_0\), independent
of \(n\), such that
\begin{equation}\label{eq:finite_riccati_uniform_bound}
\sup_{n\ge n_0}\sup_{t\in[0,T]}\lVert H_n(t)\rVert\le C_0.
\end{equation}
For \(s\in(\underline t,T)\), define
\[
\varrho_s
:=
1+\max\left\{
C_0,\,
\sup_{t\in[s,T]}\lVert H(t)\rVert
\right\},
\qquad
\mathbb B_s
:=
\left\{H_0\in\mathbb S^{3d}:\lVert H_0\rVert\le \varrho_s\right\}.
\]
The compact set \(\mathbb B_s\) contains \(H(t)\) and \(H_n(t)\)
for every \(n\ge n_0\) and \(t\in[s,T]\). Choose \(L_s<\infty\) such that
\[
\lVert \mathscr F(H_0)-\mathscr F(H_1)\rVert
\le
L_s\lVert H_0-H_1\rVert,
\qquad
H_0,H_1\in\mathbb B_s,
\]
and, by \eqref{eq:finite_riccati_vector_field_convergence}, choose
\(c_s<\infty\) such that
\[
\sup_{H_0\in\mathbb B_s}
\lVert \mathscr F_n(H_0)-\mathscr F(H_0)\rVert
\le
\frac{c_s}{n},
\qquad n\ge n_0.
\]
Since
\(H_n(T)=H(T)\), for \(t\in[s,T]\),
\[
\begin{aligned}
\lVert H_n(t)-H(t)\rVert
&\le
\int_t^T
\left\lVert
\mathscr F_n(H_n(t'))-\mathscr F(H_n(t'))
\right\rVert\,dt'
\\
&\quad+
\int_t^T
\left\lVert
\mathscr F(H_n(t'))-\mathscr F(H(t'))
\right\rVert\,dt'
\\
&\le
\frac{c_s}{n}(T-t)
+L_s\int_t^T\lVert H_n(t')-H(t')\rVert\,dt'.
\end{aligned}
\]
Gronwall therefore gives a constant \(C_s<\infty\), independent of
\(n\), such that
\begin{equation}\label{eq:finite_riccati_local_convergence}
\sup_{t\in[s,T]}\lVert H_n(t)-H(t)\rVert
\le\frac{C_s}{n}.
\end{equation}
Combining \eqref{eq:finite_riccati_uniform_bound} and
\eqref{eq:finite_riccati_local_convergence}, and taking \(n\) sufficiently
large, shows that
\begin{equation}\label{eq:limiting_riccati_uniform_bound}
\sup_{t\in[s,T]}\lVert H(t)\rVert\le C_0+1.
\end{equation}
The bound in \eqref{eq:limiting_riccati_uniform_bound} does not depend on
\(s\), so \(H\) cannot blow up as \(t\downarrow\underline t\).
The continuation criterion therefore shows that \(H\)
can be extended to the full interval \([0,T]\) and, together with all \(H_n\), remains in one compact
set. On this set, \(\mathscr F\) is Lipschitz and
\eqref{eq:finite_riccati_vector_field_convergence} is uniform. The preceding
Gronwall argument with \(s=0\) yields
\begin{equation}\label{eq:finite_riccati_limit_to_three_block}
\sup_{t\in[0,T]}\|H_n(t)-H(t)\|
\le
\frac{C}{n}.
\end{equation}

Reading off the \((\delta,\delta)\)-block of
\eqref{eq:limiting_three_block_riccati} gives
\begin{equation}\label{eq:limiting_riccati_deviation_block}
\frac{d}{dt}[H]_{\delta\delta}
=- Q
+\frac12[H]_{\delta\delta}\Lambda^{-1}
[H]_{\delta\delta},
\qquad
[H(T)]_{\delta\delta}=Q_T.
\end{equation}
This is the mean-field Riccati equation \eqref{eq:mf_riccati}, so uniqueness
gives \([H]_{\delta\delta}=P\).

Set \(P_n(t):=[H_n(t)]_{\delta\delta}\). Taking the
\((\delta,\delta)\)-block in
\eqref{eq:finite_riccati_limit_to_three_block} gives
\begin{equation}\label{eq:finite_riccati_deviation_block_convergence}
\sup_{t\in[0,T]}
\lVert P_n(t)-P(t)\rVert
\le\frac{C}{n}.
\end{equation}
Step 4 also requires convergence of the deviation-feedback block, which
determines the coefficient in the individual-deviation equation. Indeed,
\[
\bigl[\Theta_n(H_n)\bigr]_\delta
=
P_n
+\frac1n
\left(
[H_n]_{\bar x\delta}-P_n
\right)
+\frac{\varpi}{n}\Gamma^\top[H_n]_{\iota\delta}.
\]
The uniform bound \eqref{eq:finite_riccati_uniform_bound} makes every block
in the last two terms uniformly bounded, so those terms are \(O(n^{-1})\).
Together with \eqref{eq:finite_riccati_deviation_block_convergence}, this gives
\begin{equation}\label{eq:finite_theta_delta_convergence}
\sup_{t\in[0,T]}
\left\lVert
\bigl[\Theta_n(H_n(t))\bigr]_\delta-P(t)
\right\rVert
\le
\frac{C}{n}.
\end{equation}

\emph{Step 2: Augmented Mean-Field Common-State System Approximation.}
The goal of this step is to augment the finite-player common state
\((\bar X_n^k,I_n^k)\) with an auxiliary affine coordinate \(\ell_n^k\).
We show that the vector field of the resulting three-block system
\(Y_n^k:=(\bar X_n^k,I_n^k,\ell_n^k)^\top\) differs by \(O(n^{-1})\)
from the mean-field vector field in \eqref{eq:mf_y_system} on bounded sets.

Under the scaling \(\varpi_n=\varpi/n\), the finite common-state
systems \eqref{eq:pre_ann_common_subsystem} and
\eqref{eq:post_ann_common_subsystem} can be written in the same form on
\([0,T_{\mathrm{ann}}]\) for \(k=0\) and on
\([T_{\mathrm{ann}},T]\) for \(k=1,\ldots,m\), as
\begin{subequations}\label{eq:finite_common_original_system}
\begin{align}
\dot{\bar X}_n^k
&=
-\frac1{d_n}\Lambda^{-1}
\left(
[\Theta_n(H_n)]_{\bar x}\bar X_n^k
+[\Theta_n(H_n)]_{\iota}I_n^k
+B_n^\top r_n^k
+\chi_g\Lambda g^k
\right),
\label{eq:finite_common_original_X}\\
\dot I_n^k
&=
-RI_n^k
-\frac{\varpi}{d_n}\Gamma\Lambda^{-1}
\left(
[\Theta_n(H_n)]_{\bar x}\bar X_n^k
+[\Theta_n(H_n)]_{\iota}I_n^k
+B_n^\top r_n^k
+\chi_g\Lambda g^k
\right)
+\Gamma g^k.
\label{eq:finite_common_original_I}
\end{align}
\end{subequations}
To put the finite equations in the form as their mean-field counterparts in
\eqref{eq:mf_y_system}, introduce the
\(d\)-dimensional auxiliary affine coefficient
\begin{equation}\label{eq:finite_common_auxiliary_ell_definition}
\ell_n^k
:=
\left(
[\Theta_n(H_n)]_{\bar x}-P
\right)\bar X_n^k
+
\left(
[\Theta_n(H_n)]_{\iota}-\operatorname{Id}_d
\right)I_n^k
+B_n^\top r_n^k,
\qquad k=0,\ldots,m.
\end{equation}
Substituting
\eqref{eq:finite_common_auxiliary_ell_definition} into
\eqref{eq:finite_common_original_system} gives
\begin{subequations}\label{eq:finite_common_auxiliary_ell_state_system}
\begin{align}
\dot{\bar X}_n^k
&=
-\frac1{d_n}\Lambda^{-1}
\left(
P\bar X_n^k+I_n^k+\ell_n^k+\chi_g\Lambda g^k
\right),
\label{eq:finite_common_auxiliary_ell_X}\\
\dot I_n^k
&=
-RI_n^k
-\frac{\varpi}{d_n}\Gamma\Lambda^{-1}
\left(
P\bar X_n^k+I_n^k+\ell_n^k+\chi_g\Lambda g^k
\right)
+\Gamma g^k.
\label{eq:finite_common_auxiliary_ell_I}
\end{align}
\end{subequations}
These equations coincide with the \((\bar X,I)\)-components of
\eqref{eq:mf_y_system}, except that \((2+\varpi\chi_u)\) is replaced by \(d_n\).
To close the finite system, it remains to derive the equation for
\(\ell_n^k\).

Differentiating \eqref{eq:finite_common_auxiliary_ell_definition} and using
\eqref{eq:finite_common_auxiliary_ell_state_system},
\eqref{eq:mf_riccati}, \eqref{eq:fixed_interval_riccati_P},
\eqref{eq:post_ann_common_r}, and
\eqref{eq:pre_ann_common_r_ode} gives
\begin{equation}\label{eq:finite_common_auxiliary_ell_equation}
\begin{aligned}
\dot\ell_n^k(t)
={}&
-\frac{\varpi\chi_u}{2(2+\varpi\chi_u)}
P(t)\Lambda^{-1}P(t)\bar X_n^k(t)
\\
&\quad
+\frac{1}{2+\varpi\chi_u} P(t)\Lambda^{-1}
\left(
I_n^k(t)+\ell_n^k(t)+\chi_g\Lambda g^k(t)
\right)
\\
&\quad
+\upsilon_n^k
\left(
t,\bar X_n^k(t),I_n^k(t),\ell_n^k(t)
\right).
\end{aligned}
\end{equation}
For \(t\in[0,T_{\mathrm{ann}}]\) when \(k=0\) and
\(t\in[T_{\mathrm{ann}},T]\) when \(k=1,\ldots,m\), and for
\(\bar x,\iota,\lambda\in\mathbb R^d\), the residual \(\upsilon_n^k\) is given by
\begin{equation}\label{eq:finite_common_auxiliary_ell_residual_identity}
\begin{aligned}
\upsilon_n^k(t,\bar x,\iota,\lambda)
:={}&
-\left(
M_n(H_n(t))B_n
\right)^\top
\left[
H_n(t)
\begin{pmatrix}
0\\
\bar x\\
\iota
\end{pmatrix}
+r_n^k(t)
\right]
\\
&\quad
+\frac1{d_n}
\left[
2\left(
\Theta_n(H_n(t))D_nB_n
\right)^\top
-B_n^\top H_n(t)B_n
\right]
\\
&\qquad {}\times
\Lambda^{-1}
\left(
P(t)\bar x+\iota+\lambda+\chi_g\Lambda g^k(t)
\right)
\\
&\quad
+\left(
\frac1{d_n}-\frac{1}{2+\varpi\chi_u}
\right)
P(t)\Lambda^{-1}
\left(
P(t)\bar x+\iota+\lambda+\chi_g\Lambda g^k(t)
\right).
\end{aligned}
\end{equation}

With this definition,
\eqref{eq:finite_common_auxiliary_ell_equation} mirrors the equation for
\(\ell^k\) in the mean-field system \eqref{eq:mf_y_system}, apart from the
remainder \(\upsilon_n^k\), which we now show is \(O(n^{-1})\), uniformly on
bounded sets of \((\bar x,\iota,\lambda)\). First,
\begin{equation}\label{eq:finite_common_affine_uniform_bound}
\max\left\{
\sup_{t\in[0,T_{\mathrm{ann}}]}\lVert r_n^0(t)\rVert,
\max_{k=1,\ldots,m}\sup_{t\in[T_{\mathrm{ann}},T]}
\lVert r_n^k(t)\rVert
\right\}
\leq
C.
\end{equation}
Indeed, \eqref{eq:finite_riccati_uniform_bound} implies that, after the
announcement, these linear equations have uniformly bounded coefficients and
forcing, and have zero terminal data. Before the announcement, the terminal
condition is a convex combination of the resulting bounded values, and the
average affine equation has the same boundedness properties.
Further, the matrix coefficients appearing in
\eqref{eq:finite_common_auxiliary_ell_residual_identity} satisfy, uniformly in
\(t\in[0,T]\),
\begin{equation}\label{eq:finite_common_residual_matrix_bounds}
\begin{aligned}
M_n(H_n(t))B_n
&=O(n^{-1}),
\\
2\left(
\Theta_n(H_n(t))D_nB_n
\right)^\top
-B_n^\top H_n(t)B_n
&=O(n^{-1}).
\end{aligned}
\end{equation}
Indeed, the \(O(n^{-1})\) estimates from Step~1 give
\[
M_n(H_n(t))B_n
=
M(H(t))B_\delta+O(n^{-1}),
\]
where \(A_zB_\delta=E_{\mathrm{com}}B_\delta=0\) gives
\[
\begin{aligned}
M(H(t))B_\delta
&=
A_zB_\delta
+\frac{1}{2+\varpi\chi_u}
B_{\mathrm{com}}\Lambda^{-1}
\Theta(H(t))E_{\mathrm{com}}B_\delta
=0.
\end{aligned}
\]
Likewise, the estimates from Step~1 also give
\[
\begin{aligned}
2\left(
\Theta_n(H_n(t))D_nB_n
\right)^\top
-B_n^\top H_n(t)B_n
&=
2\left(
\Theta(H(t))DB_\delta
\right)^\top
-B_\delta^\top H(t)B_\delta
+O(n^{-1})\\
&=O(n^{-1}).
\end{aligned}
\]
Here \(DB_\delta=\frac12B_\delta\),
\(E_\iota B_\delta=0\), and
\(B_\delta^\top H(t)B_\delta=P(t)\) give
\[
\begin{aligned}
2\left(
\Theta(H(t))DB_\delta
\right)^\top
-B_\delta^\top H(t)B_\delta
&=
\left(
B_\delta^\top H(t)B_\delta
\right)^\top
-B_\delta^\top H(t)B_\delta
\\
&=P(t)-P(t)
=0.
\end{aligned}
\]
Together with \(d_n^{-1}-(2+\varpi\chi_u)^{-1}=O(n^{-1})\),
\eqref{eq:finite_common_affine_uniform_bound}, and
\eqref{eq:finite_riccati_uniform_bound}, this shows that
\(\upsilon_n^k=O(n^{-1})\), uniformly in time and on compact subsets of
\(\mathbb R^{3d}\). Since \(\upsilon_n^k(t,\cdot)\) is affine, its linear and
constant coefficients are \(O(n^{-1})\) as well. Hence
\eqref{eq:finite_common_auxiliary_ell_state_system} and
\eqref{eq:finite_common_auxiliary_ell_equation} can be written as
\begin{equation}\label{eq:finite_common_vector_system}
\dot{Y}_n^k
=
\mathcal A_n(t)Y_n^k
+b_n^k(t),
\end{equation}
for a matrix \(\mathcal A_n\) and vectors \(b_n^k\). Above we have shown that their coefficients satisfy
\begin{equation}\label{eq:finite_common_coefficient_convergence}
\sup_{t\in[0,T]}
\left\|
\mathcal A_n(t)-\mathcal A(t)
\right\|
+
\max\left\{
\begin{aligned}
&\sup_{t\in[0,T_{\mathrm{ann}}]}
\left\|b_n^0(t)-b^0(t)\right\|,\\
&\max_{k=1,\ldots,m}\sup_{t\in[T_{\mathrm{ann}},T]}
\left\|b_n^k(t)-b^k(t)\right\|
\end{aligned}
\right\}
\leq
\frac{C}{n}.
\end{equation}

We next recall the boundary and matching conditions. For the mean-field
system, for \(k=1,\ldots,m\), they are
\begin{equation}\label{eq:mf_common_boundary_matching_conditions}
\begin{aligned}
\bar X^0(0)&=I^0(0)=0,
&
\begin{pmatrix}
\bar X^k(T_{\mathrm{ann}})\\
I^k(T_{\mathrm{ann}})
\end{pmatrix}
&=
\begin{pmatrix}
\bar X^0(T_{\mathrm{ann}})\\
I^0(T_{\mathrm{ann}})
\end{pmatrix},
\\
\ell^k(T)&=0,
&
\ell^0(T_{\mathrm{ann}})
&=
\sum_{k'=1}^m
\bar\pi^{k'}\ell^{k'}(T_{\mathrm{ann}}).
\end{aligned}
\end{equation}
The finite boundary conditions and
\eqref{eq:finite_common_auxiliary_ell_definition} give, for
\(k=1,\ldots,m\),
\begin{equation}\label{eq:finite_common_auxiliary_ell_matching}
\begin{aligned}
\bar X_n^0(0)&=I_n^0(0)=0,
&
\begin{pmatrix}
\bar X_n^k(T_{\mathrm{ann}})\\
I_n^k(T_{\mathrm{ann}})
\end{pmatrix}
&=
\begin{pmatrix}
\bar X_n^0(T_{\mathrm{ann}})\\
I_n^0(T_{\mathrm{ann}})
\end{pmatrix},
\\
\ell_n^k(T)&=0,
&
\ell_n^0(T_{\mathrm{ann}})
&=
\sum_{k'=1}^m
\bar\pi_n^{k'}\ell_n^{k'}(T_{\mathrm{ann}}).
\end{aligned}
\end{equation}

In summary, \eqref{eq:finite_common_coefficient_convergence} implies that the
finite and mean-field ODE vector fields are \(O(n^{-1})\) apart on bounded
sets, while \eqref{eq:mf_common_boundary_matching_conditions} and
\eqref{eq:finite_common_auxiliary_ell_matching} imply that their boundary and
matching conditions are \(O(\|\bar\pi_n-\bar\pi\|)\) apart.

\emph{Step 3: Stability of the Augmented Mean-Field Common-State
Boundary-Value Problem.}
The goal of this step is to show that the \(O(n^{-1})\) vector-field error and
the \(O(\|\bar\pi_n-\bar\pi\|)\) matching error imply that the corresponding
solutions are \(O(n^{-1}+\|\bar\pi_n-\bar\pi\|)\) apart.
Unlike for an initial-value problem, this does not follow directly from
Gronwall's inequality: \((\bar X,I)\) have initial conditions, the
post-announcement affine coefficients have terminal conditions, and the
pre- and post-announcement paths are coupled through the matching conditions
at \(T_{\mathrm{ann}}\).

The uniform bounds \eqref{eq:finite_riccati_uniform_bound} and
\eqref{eq:finite_common_affine_uniform_bound} imply that
\eqref{eq:finite_common_original_system} has uniformly bounded coefficients
and forcing. Gronwall's inequality and
\eqref{eq:finite_common_auxiliary_ell_definition} therefore give
\[
\max\left\{
\sup_{t\in[0,T_{\mathrm{ann}}]}\|Y_n^0(t)\|,
\max_{k=1,\ldots,m}\sup_{t\in[T_{\mathrm{ann}},T]}\|Y_n^k(t)\|
\right\}
\leq
C.
\]

Set
\[
e_n^k
:=
Y_n^k-Y^k
=:
\begin{pmatrix}
e_{n,X}^k\\
e_{n,I}^k\\
e_{n,\ell}^k
\end{pmatrix}.
\]
Subtracting the finite and mean-field ODEs gives
\[
\dot e_n^k
=
\mathcal Ae_n^k+\Delta b_n^k,
\qquad
\Delta b_n^k
:=
\left(
\mathcal A_n-\mathcal A
\right)Y_n^k
+
b_n^k-b^k.
\]
By \eqref{eq:finite_common_coefficient_convergence} and the uniform bound on
\(Y_n^k\), \(\Delta b_n^k=O(n^{-1})\) uniformly on the corresponding time
intervals.

For \(k=1,\ldots,m\), the boundary and matching conditions for the error
system are
\[
\begin{aligned}
e_{n,X}^0(0)&=0,
&
e_{n,I}^0(0)&=0,
\\
e_{n,X}^k(T_{\mathrm{ann}})
&=e_{n,X}^0(T_{\mathrm{ann}}),
&
e_{n,I}^k(T_{\mathrm{ann}})
&=e_{n,I}^0(T_{\mathrm{ann}}),
\\
e_{n,\ell}^k(T)&=0,
&
e_{n,\ell}^0(T_{\mathrm{ann}})
&=
\sum_{k'=1}^m
\bar\pi^{k'} e_{n,\ell}^{k'}(T_{\mathrm{ann}})
+
e_{n,\pi},
\end{aligned}
\]
where
\[
e_{n,\pi}
:=
\sum_{k'=1}^m
\left(
\bar\pi_n^{k'}-\bar\pi^{k'}
\right)
\ell_n^{k'}(T_{\mathrm{ann}}).
\]
The uniform bound on \(Y_n^k\) gives
\[
\|e_{n,\pi}\|
\leq
C\|\bar\pi_n-\bar\pi\|.
\]

Set \(q_n:=e_{n,\ell}^0(0)\). Variation of constants on the
pre-announcement interval gives
\[
e_n^0(T_{\mathrm{ann}})
=
\mathcal E^{\mathrm{pre}}
\begin{pmatrix}
0\\0\\q_n
\end{pmatrix}
+\mathcal R_n^0,
\qquad
\mathcal R_n^0
:=
\int_0^{T_{\mathrm{ann}}}
\mathcal E(T_{\mathrm{ann}},s)\Delta b_n^0(s)\,ds.
\]
For each \(k=1,\ldots,m\), variation of constants on the
post-announcement interval and \(e_{n,\ell}^k(T)=0\) give
\begin{equation}\label{eq:finite_common_post_error_terminal_identity}
\begin{aligned}
0
={}&
\mathcal E^{\mathrm{post}}_{\ell X}e_{n,X}^0(T_{\mathrm{ann}})
+\mathcal E^{\mathrm{post}}_{\ell I}e_{n,I}^0(T_{\mathrm{ann}})
+\mathcal E^{\mathrm{post}}_{\ell\ell}e_{n,\ell}^k(T_{\mathrm{ann}})
+\mathcal R_{n,\ell}^k,
\\
\mathcal R_{n,\ell}^k
:={}&
\int_{T_{\mathrm{ann}}}^T
\begin{bmatrix}
\mathcal E_{\ell X}(T,s)&
\mathcal E_{\ell I}(T,s)&
\mathcal E_{\ell\ell}(T,s)
\end{bmatrix}
\Delta b_n^k(s)\,ds.
\end{aligned}
\end{equation}
Averaging the post-announcement terminal identities with weights
\(\bar\pi^k\) and using the announcement-time matching condition yields
\[
\begin{aligned}
0
={}&
\mathcal E^{\mathrm{post}}_{\ell X}e_{n,X}^0(T_{\mathrm{ann}})
+\mathcal E^{\mathrm{post}}_{\ell I}e_{n,I}^0(T_{\mathrm{ann}})
+
\mathcal E^{\mathrm{post}}_{\ell\ell}
\left(e_{n,\ell}^0(T_{\mathrm{ann}})-e_{n,\pi}\right)
+\sum_{k=1}^m\bar\pi^k\mathcal R_{n,\ell}^k.
\end{aligned}
\]
Substituting the pre-announcement representation and using
\(\mathcal E^{\mathrm{post}}\mathcal E^{\mathrm{pre}}=\mathcal E(T,0)\)
therefore gives
\[
\mathcal Mq_n
=
\mathcal E^{\mathrm{post}}_{\ell\ell}e_{n,\pi}
+\zeta_n,
\]
where
\[
\zeta_n
:=
-
\begin{bmatrix}
\mathcal E^{\mathrm{post}}_{\ell X}&
\mathcal E^{\mathrm{post}}_{\ell I}&
\mathcal E^{\mathrm{post}}_{\ell\ell}
\end{bmatrix}
\mathcal R_n^0
-
\sum_{k=1}^m\bar\pi^k\mathcal R_{n,\ell}^k.
\]
The bound on \(\Delta b_n^k\) and boundedness of \(\mathcal E\) imply
\[
\|\mathcal R_n^0\|
+\max_{1\leq k\leq m}\|\mathcal R_{n,\ell}^k\|
\leq\frac{C}{n},
\qquad
\|\zeta_n\|\leq\frac{C}{n}.
\]
Since \(\mathcal M\) is invertible by
Theorem~\ref{thm:mf_implementability},
\(\|q_n\|\leq C\varepsilon_n\). The pre-announcement error starts from
\((0,0,q_n)^\top\), so variation of constants and boundedness of
\(\mathcal E\) give
\[
\sup_{t\in[0,T_{\mathrm{ann}}]}\|e_n^0(t)\|
\leq C\varepsilon_n.
\]

Solving \eqref{eq:finite_common_post_error_terminal_identity} for
\(e_{n,\ell}^k(T_{\mathrm{ann}})\) gives
\[
e_{n,\ell}^k(T_{\mathrm{ann}})
=
-\left(\mathcal E^{\mathrm{post}}_{\ell\ell}\right)^{-1}
\left(
\mathcal E^{\mathrm{post}}_{\ell X}e_{n,X}^0(T_{\mathrm{ann}})
+\mathcal E^{\mathrm{post}}_{\ell I}e_{n,I}^0(T_{\mathrm{ann}})
+\mathcal R_{n,\ell}^k
\right)
\]
where we again used that by Theorem~\ref{thm:mf_implementability},
\(\mathcal E^{\mathrm{post}}_{\ell\ell}\) is invertible. The right-hand side
is \(O(\varepsilon_n)\), since \(e_n^0(T_{\mathrm{ann}})\) and
\(\mathcal R_{n,\ell}^k\) are both \(O(\varepsilon_n)\). Variation of constants
and boundedness of \(\mathcal E\) on each post-announcement interval yield
\begin{equation}\label{eq:finite_common_state_convergence}
\max\left\{
\begin{aligned}
&\sup_{t\in[0,T_{\mathrm{ann}}]}
\left\|Y_n^0(t)-Y^0(t)\right\|,\\
&\max_{k=1,\ldots,m}\sup_{t\in[T_{\mathrm{ann}},T]}
\left\|Y_n^k(t)-Y^k(t)\right\|
\end{aligned}
\right\}
\leq
C\varepsilon_n.
\end{equation}

Since \(\dot{\bar X}_n^k=\bar u_n^k\) and
\(\dot{\bar X}^k=\bar u^k\), taking the first components of the finite and
mean-field ODEs and using
\eqref{eq:finite_common_coefficient_convergence} and
\eqref{eq:finite_common_state_convergence} gives
\begin{equation}\label{eq:finite_common_flow_convergence}
\max\left\{
\begin{aligned}
&\sup_{t\in[0,T_{\mathrm{ann}}]}
\|\bar u_n^{0}(t)-\bar u^0(t)\|,\\
&\max_{k=1,\ldots,m}\sup_{t\in[T_{\mathrm{ann}},T]}
\|\bar u_n^{k}(t)-\bar u^k(t)\|
\end{aligned}
\right\}
\leq
C\varepsilon_n.
\end{equation}
Equation~\eqref{eq:finite_common_flow_convergence} proves
\eqref{eq:finite_to_mf_common_estimate}, while the \((\bar X,I)\)-components
of \eqref{eq:finite_common_state_convergence} give the corresponding mean
inventory and impact-state estimates.


\emph{Step 4: Player-Level Convergence.}
In this step, we prove that each finite-player inventory deviation \(\delta^{i,k}_n(t)\) is
\(O(\varepsilon_n)\)-close, uniformly in \(i\), to the mean-field deviation \(\delta^{\pi_i,k}(t)\)
of type \(\pi_n^i\), both before and after the announcement. Combining the
common-state and deviation estimates then gives the player-level inventory
and rate estimates.

Set
\[
\widetilde\ell_n^{i,0}(t)
:=B_n^\top\widetilde r_n^{i,0}(t).
\]
By \eqref{eq:pre_ann_delta_dynamics} and
\eqref{eq:mf_pre_deviation_dynamics}, the pre-announcement deviations satisfy
\[
\dot\delta_n^{i,0}(t)
=
-\frac1{a_n}\Lambda^{-1}
\left(
[\Theta_n(H_n(t))]_{\delta}\delta_n^{i,0}(t)
+\widetilde\ell_n^{i,0}(t)
\right),
\qquad
\delta_n^{i,0}(0)=0,
\]
and
\[
\dot\delta^{\pi_n^i,0}(t)
=
-\frac12\Lambda^{-1}
\left(
P(t)\delta^{\pi_n^i,0}(t)
+\widetilde\ell^{\pi_n^i,0}(t)
\right),
\qquad
\delta^{\pi_n^i,0}(0)=0.
\]
By \eqref{eq:post_ann_delta_dynamics} and
\eqref{eq:mf_post_deviation_dynamics}, the post-announcement deviations
satisfy, for \(k=1,\ldots,m\),
\[
\dot\delta_n^{i,k}(t)
=
-\frac1{a_n}\Lambda^{-1}
[\Theta_n(H_n(t))]_{\delta}\delta_n^{i,k}(t),
\qquad
\delta_n^{i,k}(T_{\mathrm{ann}})
=\delta_n^{i,0}(T_{\mathrm{ann}}),
\]
and
\[
\dot\delta^{\pi_n^i,k}(t)
=
-\frac12\Lambda^{-1}P(t)\delta^{\pi_n^i,k}(t),
\qquad
\delta^{\pi_n^i,k}(T_{\mathrm{ann}})
=\delta^{\pi_n^i,0}(T_{\mathrm{ann}}).
\]
The pre-announcement initial conditions agree, and the post-announcement
initial difference equals the pre-announcement terminal difference. Moreover,
\eqref{eq:finite_theta_delta_convergence} and
\(a_n=2-\varpi\chi_u/n\) give
\begin{equation}\label{eq:finite_deviation_coefficient_convergence}
\sup_{t\in[0,T]}
\left\|
\frac1{a_n}[\Theta_n(H_n(t))]_{\delta}
-\frac12P(t)
\right\|
\leq
\frac{C}{n}.
\end{equation}
Thus the linear state coefficients in both pairs of deviation equations are
\(O(n^{-1})\) apart. The post-announcement equations have no affine forcing.
For the pre-announcement equations, it remains to prove that
\(\widetilde\ell_n^{i,0}-\widetilde\ell^{\pi_n^i,0}\) is
\(O(\varepsilon_n)\), uniformly in \(i\) and time. Once established, this
estimate and \(a_n^{-1}-1/2=O(n^{-1})\) make the pre-announcement vector
fields \(O(\varepsilon_n)\) apart. Gronwall then gives the pre-announcement
deviation estimate. The inherited terminal difference and the
post-announcement coefficient estimate give the post-announcement deviation
estimate by a second application of Gronwall.

We now turn to showing that \(\widetilde\ell_n^{i,0}-\widetilde\ell^{\pi_n^i,0}\) is
\(O(\varepsilon_n)\). Applying \(B_n^\top\) to \eqref{eq:pre_ann_deviation_r_ode} gives
\[
\begin{aligned}
\dot{\widetilde\ell}_n^{i,0}
={}&
-\left(M_n(H_n)B_n\right)^\top\widetilde r_n^{i,0}
+\frac1{a_n}
\left[
2\left(\Theta_n(H_n)D_nB_n\right)^\top
-B_n^\top H_nB_n^-
\right]
\Lambda^{-1}\widetilde\ell_n^{i,0}.
\end{aligned}
\]
Its mean-field counterpart satisfies
\[
\dot{\widetilde\ell}^{\pi_n^i,0}(t)
=
\frac12P(t)\Lambda^{-1}\widetilde\ell^{\pi_n^i,0}(t)
\]
by \eqref{eq:mf_centered_affine_dynamics}. The mean-field equation has no
additive forcing. We compare the two terminal-value problems in three parts:
their linear coefficients, the finite-only additive forcing, and their
terminal conditions at \(T_{\mathrm{ann}}\) and show that each of these is either \(O(n^{-1})\) or \(O(\varepsilon_n)\).

We first bound the difference in the linear coefficients. Since \(B_n=B_\delta+O(n^{-1})\),
\eqref{eq:finite_riccati_limit_to_three_block} and
\(B_\delta^\top HB_\delta=P\) give, uniformly in time,
\[
B_n^\top H_nB_\delta
=B_\delta^\top HB_\delta+O(n^{-1})
=P+O(n^{-1}).
\]
Using \(B_n-B_n^-=B_\delta\) and
\eqref{eq:finite_common_residual_matrix_bounds} then gives
\[
\begin{aligned}
2\left(\Theta_n(H_n)D_nB_n\right)^\top-B_n^\top H_nB_n^-
&=
\left[
2\left(\Theta_n(H_n)D_nB_n\right)^\top-B_n^\top H_nB_n
\right]
+B_n^\top H_nB_\delta\\
&=P+O(n^{-1}).
\end{aligned}
\]
Since \(a_n^{-1}-1/2=O(n^{-1})\), the linear coefficients satisfy
\[
\frac1{a_n}
\left[
2\left(\Theta_n(H_n)D_nB_n\right)^\top-B_n^\top H_nB_n^-
\right]\Lambda^{-1}
=\frac12P\Lambda^{-1}+O(n^{-1}).
\]

We next bound the finite-only additive forcing. For this purpose, we
first obtain a uniform bound for \(\widetilde r_n^{i,0}\). Since
\(\pi_n^i,\bar\pi_n\in\Delta_m\),
\(
\sum_{k=1}^m
|\pi_n^{i,k}-\bar\pi_n^k|
\leq 2
\).
The terminal condition in \eqref{eq:pre_ann_deviation_r_ode} and
\eqref{eq:finite_common_affine_uniform_bound} therefore give, uniformly in
\(n\), \(i\), and the finite belief profile,
\[
\begin{aligned}
\left\|\widetilde r_n^{i,0}(T_{\mathrm{ann}})\right\|
&=
\left\|
\sum_{k=1}^m
(\pi_n^{i,k}-\bar\pi_n^k)r_n^k(T_{\mathrm{ann}})
\right\|
\leq
\left(
\sum_{k=1}^m
|\pi_n^{i,k}-\bar\pi_n^k|
\right)
\max_{1\leq k\leq m}
\left\|r_n^k(T_{\mathrm{ann}})\right\|
\\
&\leq
2\max_{1\leq k\leq m}
\left\|r_n^k(T_{\mathrm{ann}})\right\|
\leq C.
\end{aligned}
\]
The coefficients in \eqref{eq:pre_ann_deviation_r_ode} are uniformly bounded
by the matrix definitions and \eqref{eq:finite_riccati_uniform_bound}.
Gronwall consequently gives
\begin{equation}\label{eq:finite_centered_r_uniform_bound}
\max_{1\leq i\leq n}
\sup_{t\in[0,T_{\mathrm{ann}}]}
\|\widetilde r_n^{i,0}(t)\|
\leq C.
\end{equation}
The first estimate in \eqref{eq:finite_common_residual_matrix_bounds} and
\eqref{eq:finite_centered_r_uniform_bound} now yield, uniformly in \(i\) and
time,
\[
\left(M_n(H_n)B_n\right)^\top\widetilde r_n^{i,0}
=O(n^{-1}).
\]
Thus the finite-only additive forcing is \(O(n^{-1})\).

Let \(\rho_n^i\)
denote the residual formed by the forcing and the linear-coefficient error.
The finite ODE can then be written as
\begin{equation}\label{eq:finite_centered_projection_ode}
\dot{\widetilde\ell}_n^{i,0}(t)
=
\frac12P(t)\Lambda^{-1}\widetilde\ell_n^{i,0}(t)+\rho_n^i(t),
\qquad
\max_{1\le i\le n}
\sup_{t\in[0,T_{\mathrm{ann}}]}
\|\rho_n^i(t)\|
\leq
\frac{C}{n}.
\end{equation}
It remains to compare the terminal conditions at
\(T_{\mathrm{ann}}\). Their comparison uses the post-announcement common
auxiliary coordinates \(\ell_n^k\) and \(\ell^k\), whose convergence was
established in Step~3: For each \(k=1,\ldots,m\),
\eqref{eq:finite_common_auxiliary_ell_matching} gives
\[
\bar X_n^k(T_{\mathrm{ann}})=\bar X_n^0(T_{\mathrm{ann}}),
\qquad
I_n^k(T_{\mathrm{ann}})=I_n^0(T_{\mathrm{ann}}).
\]
Thus \eqref{eq:finite_common_auxiliary_ell_definition} reads
\[
\begin{aligned}
\ell_n^k(T_{\mathrm{ann}})
={}&
\left(
[\Theta_n(H_n(T_{\mathrm{ann}}))]_{\bar x}-P(T_{\mathrm{ann}})
\right)\bar X_n^0(T_{\mathrm{ann}})
\\
&+
\left(
[\Theta_n(H_n(T_{\mathrm{ann}}))]_{\iota}-\operatorname{Id}_d
\right)I_n^0(T_{\mathrm{ann}})
+B_n^\top r_n^k(T_{\mathrm{ann}}).
\end{aligned}
\]
The first two terms are independent of \(k\). The terminal condition in
\eqref{eq:pre_ann_deviation_r_ode} and
\(\sum_k(\pi_n^{i,k}-\bar\pi_n^k)=0\) therefore give
\[
\begin{aligned}
\widetilde\ell_n^{i,0}(T_{\mathrm{ann}})
&=
\sum_{k=1}^m
(\pi_n^{i,k}-\bar\pi_n^k)
B_n^\top r_n^k(T_{\mathrm{ann}})
=
\sum_{k=1}^m
(\pi_n^{i,k}-\bar\pi_n^k)\ell_n^k(T_{\mathrm{ann}}).
\end{aligned}
\]
The mean-field terminal condition in
\eqref{eq:mf_centered_affine_dynamics} is
\[
\widetilde\ell^{\pi_n^i,0}(T_{\mathrm{ann}})
=
\sum_{k=1}^m
(\pi_n^{i,k}-\bar\pi^k)\ell^k(T_{\mathrm{ann}}).
\]
Subtracting the two terminal conditions gives
\begin{equation}\label{eq:finite_centered_terminal_error_identity}
\begin{aligned}
&\widetilde\ell_n^{i,0}(T_{\mathrm{ann}})
-\widetilde\ell^{\pi_n^i,0}(T_{\mathrm{ann}})
\\
&\quad=
\sum_{k=1}^m
(\pi_n^{i,k}-\bar\pi_n^k)
\left(
\ell_n^k(T_{\mathrm{ann}})
-\ell^k(T_{\mathrm{ann}})
\right)
+
\sum_{k=1}^m
(\bar\pi^k-\bar\pi_n^k)
\ell^k(T_{\mathrm{ann}}).
\end{aligned}
\end{equation}
The bound
\(\sum_{k=1}^m|\pi_n^{i,k}-\bar\pi_n^k|\leq2\) gives
\[
\left\|
\sum_{k=1}^m
(\pi_n^{i,k}-\bar\pi_n^k)
(\ell_n^k(T_{\mathrm{ann}})-\ell^k(T_{\mathrm{ann}}))
\right\|
\leq
2\max_{1\leq k\leq m}
\|\ell_n^k(T_{\mathrm{ann}})-\ell^k(T_{\mathrm{ann}})\|.
\]
The \(\ell_n^k-\ell^k\) component of
\eqref{eq:finite_common_state_convergence}, the boundedness of the
mean-field affine coefficients, and the fixed dimension \(m\) therefore
give
\begin{equation}\label{eq:finite_centered_terminal_error_bound}
\max_{1\leq i\leq n}
\left\|
\widetilde\ell_n^{i,0}(T_{\mathrm{ann}})
-\widetilde\ell^{\pi_n^i,0}(T_{\mathrm{ann}})
\right\|
\leq
C\varepsilon_n,
\end{equation}
where \(C\) is independent of \(n\) and the finite belief profile.
Combining the vector-field estimate in
\eqref{eq:finite_centered_projection_ode} with the terminal-condition
estimate \eqref{eq:finite_centered_terminal_error_bound}, backward Gronwall
gives
\begin{equation}\label{eq:finite_centered_projection_convergence}
\max_{1\le i\le n}
\sup_{t\in[0,T_{\mathrm{ann}}]}
\left\|
\widetilde\ell_n^{i,0}(t)-\widetilde\ell^{\pi_n^i,0}(t)
\right\|
\leq
C\varepsilon_n.
\end{equation}

Together with \(a_n^{-1}-1/2=O(n^{-1})\),
\eqref{eq:finite_deviation_coefficient_convergence} and
\eqref{eq:finite_centered_projection_convergence} allow Gronwall to be
applied before and after the announcement, giving
\begin{equation}\label{eq:finite_deviation_state_convergence}
\begin{aligned}
&\max_{1\le i\le n}
\sup_{t\in[0,T_{\mathrm{ann}}]}
\left\|
\delta_n^{i,0}(t)-\delta^{\pi_n^i,0}(t)
\right\|
\\
&\quad+
\max_{1\le i\le n}
\max_{1\le k\le m}
\sup_{t\in[T_{\mathrm{ann}},T]}
\left\|
\delta_n^{i,k}(t)-\delta^{\pi_n^i,k}(t)
\right\|
\leq
C\varepsilon_n.
\end{aligned}
\end{equation}
Since
\(X_n^{i,k}-X^{\pi_n^i,k}
=(\bar X_n^k-\bar X^k)
+(\delta_n^{i,k}-\delta^{\pi_n^i,k})\),
\eqref{eq:finite_common_state_convergence} and
\eqref{eq:finite_deviation_state_convergence} give
\(X_n^{i,k}-X^{\pi_n^i,k}=O(\varepsilon_n)\), uniformly over players,
scenarios, and the corresponding time intervals.

Before the announcement, the finite and mean-field feedbacks satisfy
\[
\begin{aligned}
u_n^{i,0}
&=\bar u_n^0
-\frac1{a_n}\Lambda^{-1}
\left(
[\Theta_n(H_n)]_\delta\delta_n^{i,0}
+\widetilde\ell_n^{i,0}
\right),
\\
u^{\pi_n^i,0}
&=\bar u^0
-\frac12\Lambda^{-1}
\left(
P\delta^{\pi_n^i,0}
+\widetilde\ell^{\pi_n^i,0}
\right).
\end{aligned}
\]
After the announcement, for \(k=1,\ldots,m\), they satisfy
\[
\begin{aligned}
u_n^{i,k}
&=\bar u_n^k
-\frac1{a_n}\Lambda^{-1}
[\Theta_n(H_n)]_\delta\delta_n^{i,k},
\\
u^{\pi_n^i,k}
&=\bar u^k
-\frac12\Lambda^{-1}P\delta^{\pi_n^i,k}.
\end{aligned}
\]
Subtracting the corresponding formulas, every term is
\(O(\varepsilon_n)\) by the common-flow, coefficient, centered-affine, and
deviation estimates established above. This gives
\eqref{eq:finite_to_mf_control_estimate}.
\end{proof}

\subsection{Proof of Theorem~\ref{thm:mf_approximate_nash}}
\label{subsec:proof_mf_approximate_nash}

\begin{proof}
The mean-field feedbacks \eqref{eq:mf_pre_feedback} and
\eqref{eq:mf_post_feedback} are affine in the player's inventory, with
uniformly bounded deterministic coefficients. The boundary condition for
\(\ell^{\pi,0}\) is affine in \(\pi\), and its ODE is linear. Consequently,
\(u^{\pi,q}(t)\) is affine in \(\pi\) for \(q=0,\ldots,m\), with coefficients
uniformly bounded in \(q\) and time. Hence the lifted
pre-announcement feedback is jointly Borel in \((\pi,x)\). The lifted policies
are non-anticipative and globally Lipschitz in own inventory. Conditional on
each scenario, the inventory equations therefore have unique solutions with
square-integrable rates, which in turn uniquely determine the impact state.
Thus \(\widehat\Phi_n\in\mathcal U_n\).

For the Nash estimate, fix player \(i\) and retain the opponents' policies
\(\widehat\Phi_n^{-i}\). For \(j\ne i\),
\eqref{eq:mf_lifted_finite_profile} shows that \(\widehat\Phi_n^j\) depends only
on time, the announced scenario, player \(j\)'s own inventory, and the fixed
mean-field coefficients determined by \((\bar\pi,\pi_n^j)\). Under any deviation
by player \(i\), player \(j\)'s inventory therefore solves the same autonomous
ODE with the same initial condition. Hence \(\widehat\Phi_n^j\) induces the
rates \(u^{\pi_n^j,0}\) before the announcement and \(u^{\pi_n^j,k}\) after
scenario \(k\) is announced, independently of player \(i\)'s deviation.

Now let \(\Psi_n^i\) be an arbitrary admissible deviation. By independence of
\(S^0\) and \(\Xi\), the conditioning argument used in the proof of
Theorem~\ref{thm:stitched_hjb_verification} allows us to condition on \(S^0\)
and integrate at the end. It therefore suffices to consider deviations that are deterministic
\(L^2\)-rate paths \(v=(v^0,v^1,\ldots,v^m)\), where \(v^0\) is common before the
announcement and \(v^k\) is used after scenario \(k\) is announced. Define the seminorm
\[
\|v\|_i^2
:=
\int_0^{T_{\mathrm{ann}}}\|v^0(t)\|^2\,dt
+
\sum_{k=1}^m\pi_n^{i,k}
\int_{T_{\mathrm{ann}}}^T\|v^k(t)\|^2\,dt.
\]

Define
\(e_{n,i}^k\) as the difference, along scenario \(k\), between the
\(\varpi/n\)-scaled aggregate rate of player \(i\)'s opponents and the
mean-field flow \(\varpi\bar u\):
\[
e_{n,i}^{k}(t)
:=
\begin{cases}
\displaystyle
\frac{\varpi}{n}\sum_{j\ne i}u^{\pi_n^j,0}(t)-\varpi\bar u^0(t),
&0\leq t<T_{\mathrm{ann}},\\[6pt]
\displaystyle
\frac{\varpi}{n}\sum_{j\ne i}u^{\pi_n^j,k}(t)-\varpi\bar u^k(t),
&T_{\mathrm{ann}}\leq t\leq T.
\end{cases}
\]
As noted at the start of the proof, the map
\(\pi\mapsto u^{\pi,q}(t)\) is affine for \(q=0,\ldots,m\), with uniformly
bounded coefficients. Let \(A_u^q(t)\in\mathbb R^{d\times m}\) denote its
linear coefficient which is uniformly bounded. The consistency condition
\eqref{eq:mf_average_consistency} gives
\[
\bar u^q(t)
=
\int_{\Delta_m}u^{\pi,q}(t)\,\Pi(d\pi).
\]
Consequently,
\[
\frac{\varpi}{n}\sum_{j\ne i}u^{\pi_n^j,q}(t)-\varpi\bar u^q(t)
=
\varpi A_u^q(t)(\bar\pi_n-\bar\pi)
-\frac{\varpi}{n} u^{\pi_n^i,q}(t).
\]
Uniform boundedness of \(A_u^q\) and the mean-field rates therefore gives
\begin{equation}\label{eq:mf_profile_opponent_flow_bound}
\max_{1\leq i\leq n}
\max_{1\leq k\leq m}
\sup_{t\in[0,T]}
\|e_{n,i}^{k}(t)\|
\leq
C\left(
\|\bar\pi_n-\bar\pi\|+\frac1n
\right)
=C\varepsilon_n.
\end{equation}
Set
\[
F_{n,i}^k(t)
:=
\int_0^t e^{-R(t-s)}\Gamma e_{n,i}^k(s)\,ds
+\chi_u\Lambda e_{n,i}^k(t).
\]
Equation~\eqref{eq:mf_profile_opponent_flow_bound} and boundedness of the
convolution kernel give
\[
\max_{1\leq i\leq n}
\max_{1\leq k\leq m}
\sup_{t\in[0,T]}
\|F_{n,i}^k(t)\|
\leq
C\varepsilon_n.
\]

Write \(J^{\pi_n^i}(v;\bar u)\) and
\(J_n^i(v;\widehat\Phi_n^{-i})\) for the objectives induced by the
deterministic plan. Variation of constants in the impact equation and
\eqref{eq:reduced_running_cost} give the exact decomposition
\[
\begin{aligned}
J_n^i(v;\widehat\Phi_n^{-i})
={}&
J^{\pi_n^i}(v;\bar u)
+
\sum_{k=1}^m\pi_n^{i,k}
\left[
\int_0^{T_{\mathrm{ann}}}
\left(F_{n,i}^k(t)\right)^\top v^0(t)\,dt
+
\int_{T_{\mathrm{ann}}}^T
\left(F_{n,i}^k(t)\right)^\top v^k(t)\,dt
\right]
\\
&+
\frac{\varpi}{n}
\sum_{k=1}^m\pi_n^{i,k}
\mathcal C\left(
\mathbf 1_{[0,T_{\mathrm{ann}})}v^0
+\mathbf 1_{[T_{\mathrm{ann}},T]}v^k
\right).
\end{aligned}
\]

Define the rate plan induced by \(\widehat\Phi_n^i\) in
\eqref{eq:mf_lifted_finite_profile} by
\[
v^{i,*}
:=
\left(
u^{\pi_n^i,0},
u^{\pi_n^i,1},\ldots,u^{\pi_n^i,m}
\right)
\]
and set \(\Delta v:=v-v^{i,*}\). Applying the
decomposition to \(v\) and \(v^{i,*}\) and subtracting gives
\[
\begin{aligned}
&J_n^i(v;\widehat\Phi_n^{-i})
-J_n^i(v^{i,*};\widehat\Phi_n^{-i})
\\
&\quad=
J^{\pi_n^i}(v;\bar u)
-J^{\pi_n^i}(v^{i,*};\bar u)
\\
&\qquad+
\sum_{k=1}^m\pi_n^{i,k}
\left[
\int_0^{T_{\mathrm{ann}}}
\left(F_{n,i}^k(t)\right)^\top\Delta v^0(t)\,dt
+
\int_{T_{\mathrm{ann}}}^T
\left(F_{n,i}^k(t)\right)^\top\Delta v^k(t)\,dt
\right]
\\
&\qquad+
\frac{\varpi}{n}
\sum_{k=1}^m\pi_n^{i,k}
\left[
\mathcal C\left(
\mathbf 1_{[0,T_{\mathrm{ann}})}v^0
+\mathbf 1_{[T_{\mathrm{ann}},T]}v^k
\right)
-\mathcal C\left(
\mathbf 1_{[0,T_{\mathrm{ann}})}u^{\pi_n^i,0}
+\mathbf 1_{[T_{\mathrm{ann}},T]}u^{\pi_n^i,k}
\right)
\right].
\end{aligned}
\]
We bound each of the three terms on the right-hand side from below.

First, since every deterministic scenario-dependent rate plan is admissible,
the type-\(\pi\) HJB verification in the proof of
Theorem~\ref{thm:mf_solution} implies that \(v^{i,*}\) minimizes the
quadratic functional \(J^{\pi_n^i}(\,\cdot\,;\bar u)\) over deterministic
scenario-dependent plans. Its first variation vanishes at \(v^{i,*}\). Since
\(\Lambda\succ0\) and \(Q,Q_T\succeq0\),
\[
J^{\pi_n^i}(v;\bar u)
-J^{\pi_n^i}(v^{i,*};\bar u)
\geq
\lambda_{\min}(\Lambda)\|\Delta v\|_i^2.
\]
Second, the bound on \(F_{n,i}^k\) and Cauchy--Schwarz imply
\[
\sum_{k=1}^m\pi_n^{i,k}
\left[
\int_0^{T_{\mathrm{ann}}}
\left(F_{n,i}^k(t)\right)^\top\Delta v^0(t)\,dt
+
\int_{T_{\mathrm{ann}}}^T
\left(F_{n,i}^k(t)\right)^\top\Delta v^k(t)\,dt
\right]
\geq
-C\varepsilon_n\|\Delta v\|_i.
\]
Third, for \(w,h\in L^2([0,T],\mathbb R^d)\), the definition of
\(\mathcal C\) gives
\[
\begin{aligned}
\mathcal C(w+h)-\mathcal C(w)
={}&
\mathcal C(h)
+
\int_0^T
w(t)^\top
\left(
\int_0^t e^{-R(t-s)}\Gamma h(s)\,ds
\right)dt
\\
&+
\int_0^T
h(t)^\top
\left(
\int_0^t e^{-R(t-s)}\Gamma w(s)\,ds
\right)dt.
\end{aligned}
\]
For scenario \(k\), apply this identity with
\[
\begin{aligned}
w
&=
\mathbf 1_{[0,T_{\mathrm{ann}})}u^{\pi_n^i,0}
+\mathbf 1_{[T_{\mathrm{ann}},T]}u^{\pi_n^i,k},
\\
h
&=
\mathbf 1_{[0,T_{\mathrm{ann}})}\Delta v^0
+\mathbf 1_{[T_{\mathrm{ann}},T]}\Delta v^k.
\end{aligned}
\]
Assumption~\ref{ass:standing_model} gives \(\mathcal C(h)\geq0\). The two
remaining integrals are bounded using the transient-impact kernel and the
uniform bound on the mean-field rates. Cauchy--Schwarz therefore gives
\[
\sum_{k=1}^m\pi_n^{i,k}
\left[
\mathcal C\left(
\mathbf 1_{[0,T_{\mathrm{ann}})}v^0
+\mathbf 1_{[T_{\mathrm{ann}},T]}v^k
\right)
-\mathcal C\left(
\mathbf 1_{[0,T_{\mathrm{ann}})}u^{\pi_n^i,0}
+\mathbf 1_{[T_{\mathrm{ann}},T]}u^{\pi_n^i,k}
\right)
\right]
\geq
-C\|\Delta v\|_i.
\]
Substituting these three bounds and using \(n^{-1}\leq\varepsilon_n\) yields
the first inequality below for some \(C_1>0\). Completing the square yields
the second with \(C_2:=C_1^2/(4\lambda_{\min}(\Lambda))\):
\[
\begin{aligned}
&J_n^i(v;\widehat\Phi_n^{-i})
-J_n^i(v^{i,*};\widehat\Phi_n^{-i})
\geq \lambda_{\min}(\Lambda)\|\Delta v\|_i^2-C_1\varepsilon_n\|\Delta v\|_i
\geq -C_2\varepsilon_n^2.
\end{aligned}
\]
The bound is uniform over deterministic scenario-dependent plans. Integrating it over
the law of \(S^0\) gives
\eqref{eq:mf_profile_approximate_nash} for every admissible deviation.
\end{proof}


\end{document}
